\section{Finite-order Hecke characters at a complex place}
\label{sec:complex-place}
The zero-distribution estimate extends directly to Hecke functions, but
the archimedean cutoff changes from a Fourier to a Hankel concentration
problem.

Let $K=\Q(\sqrt{-3})$, let $\mathbb A_K$ denote its adele ring, and let
$C_K=K^\times\backslash\mathbb A_K^\times$ have its idelic norm
$|\cdot|:C_K\to\R_{>0}$. Give the compact norm-one subgroup Haar
measure one and the norm variable measure $dr/r$. The finite-order
test algebra considered here consists of the finite sums
\[
 h(c)=\sum_{\chi\in F}h_\chi(\log|c|)\overline{\chi(c)},
 \qquad h_\chi\in C_c^\infty(\R),
\]
where $F$ is a finite set of continuous finite-order characters of
$C_K$. These characters are trivial on the connected norm subgroup
and at the complex infinite place. This algebra specifies the
finite-order sectors; it is not the full Bruhat algebra with all
infinity types. Write $\mathfrak f_\chi$ for the primitive conductor,
$Q_\chi=3N\mathfrak f_\chi$, and $c_\chi=2\pi/\sqrt{Q_\chi}$.
The completed Hecke function has factor
\[
 \Lambda_K(s,\chi)=c_\chi^{-s}\Gamma(s)L_K(s,\chi),
 \qquad
 \Lambda_K(s,\chi)=\epsilon_\chi
                      \Lambda_K(1-s,\overline\chi),
 \quad |\epsilon_\chi|=1,
\]
up to a fixed nonzero normalization constant. In particular its gamma
factor has argument $s$, rather than a real-place half argument.

\begin{proposition}[The complex-place concentration operator]
\label{prop:complex-concentration}
For $c>0$ define
\[
 (\mathcal H_c\varphi)(y)
      =c\int_0^\infty J_0(2c\sqrt{xy})\varphi(x)\,dx,
 \qquad
 (T_X\varphi)(x)=\mathbf1_{(0,X)}(x)\varphi(x).
\]
Here $J_0$ is the Bessel function of the first kind and order zero.
The operator $\mathcal H_c$ extends to a unitary involution of
$L^2((0,\infty),dx)$. For $\varphi\in C_c^\infty((0,\infty))$ its
Mellin transform satisfies, initially for $0<\Re s<1$ and then
by meromorphic continuation,
\begin{equation}\label{eq:hankel-mellin}
 \mathcal M(\mathcal H_c\varphi)(s)
   =c^{1-2s}\frac{\Gamma(s)}{\Gamma(1-s)}
                       \mathcal M\varphi(1-s),
 \qquad
 \mathcal M\varphi(s)=\int_0^\infty\varphi(x)x^{s-1}\,dx.
\end{equation}
The positive contraction
$A_{X,Y}=T_X\mathcal H_cT_Y\mathcal H_cT_X$ is trace class and, with
$Z=2c\sqrt{XY}$,
\begin{equation}\label{eq:hankel-trace}
 \Tr A_{X,Y}
   =\int_0^Z tJ_0(t)^2\log(Z/t)\,dt
   =\frac Z\pi+O(\log(2+Z))\qquad(Z\to\infty).
\end{equation}
\end{proposition}

\begin{proof}
For the Euclidean Fourier transform on $\R^2$ with kernel
$e^{-2\pi i x\cdot\xi}$, polar coordinates applied to
$f(z)=\varphi(|z|^2)$ give
\[
 \widehat f(\xi)
   =\pi\int_0^\infty J_0(2\pi\sqrt{x|\xi|^2})\varphi(x)\,dx.
\]
Moreover $\|f\|_2^2=\pi\|\varphi\|_2^2$. Fourier unitarity and
Fourier inversion on radial functions prove the assertion for $c=\pi$.
For general $c$, apply the unitary dilation
$g(y)\mapsto(c/\pi)g((c/\pi)^2y)$ after $\mathcal H_\pi$.
The dilation covariance of Fourier inversion shows that its square
is again the identity.

In the Mellin integral substitute $t=2c\sqrt{xy}$ and use
\[
 \int_0^\infty t^{2s-1}J_0(t)\,dt
       =2^{2s-1}\frac{\Gamma(s)}{\Gamma(1-s)}.
\]
The interchange is absolutely justified for $0<\Re s<1/4$; continuation
gives \eqref{eq:hankel-mellin}. The two projections and the unitary
involution show $0\leq A_{X,Y}\leq I$. Its trace is the squared
Hilbert--Schmidt norm of $T_Y\mathcal H_cT_X$, namely
\[
 c^2\int_0^X\int_0^Y J_0(2c\sqrt{xy})^2\,dy\,dx.
\]
Set $x=Xu^2$, $y=Yv^2$, and integrate first with respect to the
product $uv$ to obtain the first equality in \eqref{eq:hankel-trace}.
The identity
$\int_0^ZtJ_0(t)^2dt=\tfrac12Z^2(J_0(Z)^2+J_1(Z)^2)$ gives
\[
 \frac d{dZ}\Tr A_{X,Y}
       =\frac Z2\bigl(J_0(Z)^2+J_1(Z)^2\bigr)
       =\frac1\pi+O(Z^{-1}).
\]
The last equality follows from the Bessel asymptotic expansions.
Integration from a fixed positive value proves the asserted error.
\end{proof}

This is the finite Hankel, or circular prolate, concentration problem;
see \citet[Section~1]{boulsane2017} for its standard radius-variable
normalization. The norm-variable theta map must also be centered:
\[
\Theta_\chi\varphi(x)=\sqrt{x}
       \sum_{\mathfrak a}\chi(\mathfrak a)\varphi(N\mathfrak a\,x).
\]
Initially for $\Re s>1$ the uncorrected sum satisfies
$\int_0^\infty\Theta_\chi\varphi(x)x^{s-1/2}dx/x
=L_K(s,\chi)\mathcal M\varphi(s)$.
The logarithmic map used in the cutoff is
$E_\chi\varphi(t)=\Theta_\chi\varphi(e^t)
-\mathbf1_{\chi=1}r_KI\varphi\,e^{-t/2}$.
Its principal continuation requires the pole subtraction. Formula
\eqref{eq:hankel-mellin} matches the Hecke functional equation; a
one-dimensional sine or cosine transform does not. The concentration
trace \eqref{eq:hankel-trace} is therefore a normalization result,
not a comparison of the Hecke theta-image projection with a zero
distribution.

\section{A finite Hankel selection and its interior approximation}
\label{sec:hankel-selection}
This section constructs a complex-place concentration cutoff without
using a zero-location theorem and proves its archimedean approximation
estimates. It does not estimate the inverse Gram matrix of its Hecke
theta images.

Fix $c>0$ and use the unitary involution $\mathcal H_c$ of
Proposition~\ref{prop:complex-concentration}. On $L^2((0,\infty),dx)$
put
\[
 P_\Lambda=1_{(0,\Lambda)},\qquad
 B_\Lambda=\mathcal H_cP_\Lambda\mathcal H_c,\qquad
 A_\Lambda=P_\Lambda B_\Lambda P_\Lambda,
 \qquad \Pi_\Lambda=1_{[1/2,1]}(A_\Lambda).
\]
The last expression is the spectral projection. It defines a finite
Hankel, or circular prolate, selection by concentration, rather than
by an arbitrarily rounded index. It depends on $c$ and $\Lambda$ and
not on the zeros of an $L$-function.
Let $(Uf)(t)=e^{t/2}f(e^t)$, and write $H_D^s$ for the pullback
of the ordinary Sobolev space $H^s(\R,dt)$ by $U$.

\begin{proposition}[Size of the Hankel selection]
\label{prop:hankel-selection-size}
The selected concentration eigenvalues are at least $1/2$, while
every omitted positive eigenvalue is less than $1/2$. With $c$ fixed,
\[
 \operatorname{rank}\Pi_\Lambda
     =\frac{2c\Lambda}{\pi}+O_c(\log(2+\Lambda)).
\]
More precisely,
\[
 \Tr(A_\Lambda-A_\Lambda^2)=O_c(\log(2+\Lambda)),
 \qquad
 |\operatorname{rank}\Pi_\Lambda-\Tr A_\Lambda|
          \leq2\Tr(A_\Lambda-A_\Lambda^2).
\]
\end{proposition}
\begin{proof}
The concentration operator is positive, compact, and at most the
identity. Its trace is finite by
Proposition~\ref{prop:complex-concentration}, so the spectral
selection is finite dimensional. The assertions about selected and
omitted eigenvalues follow from its definition.

Here is a direct estimate of the transition mass. Change variables
$r=2c\sqrt{\Lambda x}$ by the unitary map
\[
 (J_\Lambda f)(r)
    =\left(\frac{r}{2c^2\Lambda}\right)^{1/2}
                        f\left(\frac{r^2}{4c^2\Lambda}\right).
\]
Under this map $B_\Lambda$ becomes the Bessel projection with kernel
\[
 K(r,s)=\sqrt{rs}\int_0^1vJ_0(rv)J_0(sv)\,dv,
\]
and $P_\Lambda$ becomes restriction to $(0,Z)$, where $Z=2c\Lambda$.
The Bessel differential equation, integrated once, gives for $r\ne s$
\[
 K(r,s)=\frac{\sqrt{rs}}{s^2-r^2}
       \bigl(sJ_1(s)J_0(r)-rJ_1(r)J_0(s)\bigr),
 \qquad
 K(r,r)=\frac r2\bigl(J_0(r)^2+J_1(r)^2\bigr).
\]
The standard bounds $|J_0(t)|+|J_1(t)|\leq C(1+t)^{-1/2}$
follow, for example, by splitting their trigonometric integral
representations at distance $t^{-1/2}$ from the stationary points
and integrating by parts on the complement. Their power series
give $J_0(t)=O(1)$ and $J_1(t)=O(t)$ near zero.
For $r,s\geq1$, the displayed off-diagonal formula gives
$|K(r,s)|\leq C|r-s|^{-1}$, while Cauchy--Schwarz in its integral
formula and the diagonal bound give $|K(r,s)|\leq C$.
For $0<r<1$ and $s\geq2$, the same formula gives
$|K(r,s)|\leq C\sqrt r/s$.

Since $B_\Lambda$ is a projection, its kernel therefore yields,
for $Z\geq2$,
\[
 \begin{split}
 \Tr(A_\Lambda-A_\Lambda^2)
 &=\int_0^Z\int_Z^\infty|K(r,s)|^2\,ds\,dr\\
 &\leq C\int_1^Z\int_Z^\infty
          \min(1,|s-r|^{-2})\,ds\,dr+C/Z\\
 &\leq C\log(2+Z).
 \end{split}
\]
The final inequality follows by putting $u=Z-r$ and $v=s-Z$:
the inner integral is at most $C/(1+u)$.
For each $0\leq\nu\leq1$,
$|1_{\nu\geq1/2}-\nu|\leq2\nu(1-\nu)$.
Sum over the eigenvalues and use
$\Tr A_\Lambda=Z/\pi+O(\log(2+Z))$ from
\eqref{eq:hankel-trace}. This proves all assertions.
\end{proof}

\begin{proposition}[Band extensions and dilation regularity]
\label{prop:hankel-raw-estimates}
There is $d>0$ such that, for all sufficiently large $\Lambda$ and
every $\psi\in\operatorname{Ran}\Pi_\Lambda$,
\[
 \int_{d/(c^2\Lambda)}^\Lambda|\psi(x)|^2dx
       \geq\tfrac12\|\psi\|_2^2,
 \qquad
 \|\psi\|_{H_D^s}\leq C_{c,s}(1+\Lambda)^s\|\psi\|_2,
                   \quad0<s<1/2.
\]
The band extension
$W_\Lambda\psi=B_\Lambda(A_\Lambda|_{\operatorname{Ran}\Pi_\Lambda})^{-1}\psi$
satisfies $P_\Lambda W_\Lambda\psi=\psi$ and
$\|W_\Lambda\|\leq\sqrt2$.
Its ordinary raw derivatives obey
\[
 \sup_{x\geq0}|(W_\Lambda\psi)^{(j)}(x)|
             \leq C_{c,j}\Lambda^{j+1/2}\|\psi\|_2
                 \qquad(j=0,1,\ldots).
\]
\end{proposition}
\begin{proof}
On the selected space the inverse of $A_\Lambda$ has norm at most
two, and
\[
 \|W_\Lambda\psi\|_2^2
     =\langle\psi,A_\Lambda^{-1}\psi\rangle\leq2\|\psi\|_2^2.
\]
The equality after restriction is the definition of $A_\Lambda$.
Write $w=W_\Lambda\psi$. Its Hankel transform is supported in
$(0,\Lambda)$. The bound $|J_0|\leq1$ and Cauchy--Schwarz give
$|w(x)|\leq c\sqrt\Lambda\|w\|_2$.
Thus the squared mass in $(0,d/(c^2\Lambda))$ is at most
$2d\|\psi\|_2^2$. Choose $d\leq1/4$.
For the derivatives, the identity
\[
 \partial_x^jJ_0(2c\sqrt{xy})
       =(-c)^j(y/x)^{j/2}J_j(2c\sqrt{xy})
\]
and $|J_j(t)|\leq(t/2)^j/j!$ imply
$|\partial_x^jJ_0(2c\sqrt{xy})|\leq c^{2j}y^j/j!$.
The latter Bessel bound follows directly from its Poisson integral
representation. Integrate the squared bound over $0<y<\Lambda$
and apply Cauchy--Schwarz. The formula extends continuously to
$x=0$ by the Bessel series.

To estimate the dilation derivative, put $F(z)=w(|z|^2)$ on
$\R^2$. The radial Fourier formula shows that $\widehat F$ is
supported in $|\xi|\leq(c/\pi)\sqrt\Lambda$, and
$\|F\|_2^2=\pi\|w\|_2^2$. Plancherel gives
$\|\nabla F\|_2\leq2c\sqrt\Lambda\|F\|_2$.
Since $x\partial_xw$ corresponds to $\tfrac12z\cdot\nabla F$,
\[
 \|xw'(x)\|_{L^2(0,\Lambda)}
             \leq c\Lambda\|w\|_2.
\]
Consequently $f=U\psi$, on $(-\infty,L)$ with $L=\log\Lambda$,
has interior derivative norm at most
$C_c(1+\Lambda)\|\psi\|_2$. Its value tends to zero at the lower
end, since $w$ is bounded near zero. The elementary trace estimate
$|f(L-)|^2\leq2\|f\|_2\|f'\|_2$ and the distributional Fourier
identity for its zero extension show that, for $M\geq1$,
\[
 \|f\|_{H^s}^2\leq C_s\bigl(
 M^{2s}\|f\|_2^2+M^{2s-2}\|f'\|_2^2
                         +M^{2s-1}|f(L-)|^2\bigr).
\]
The last frequency integral converges exactly when $s<1/2$.
Take $M=1+\Lambda$ to obtain the stated estimate.
\end{proof}

\begin{proposition}[Strict-margin Hankel approximation]
\label{prop:hankel-interior-approximation}
Let $p\in C_c^\infty([0,1))$ equal one near zero, with
$0\leq p\leq1$, and set
\[
 P^p_\Lambda=p(x/\Lambda),\qquad
 B^p_\Lambda=\mathcal H_cp(x/\Lambda)\mathcal H_c.
\]
For every $0\leq\sigma<1/2$ and $A>0$,
\[
 \|(1-\Pi_\Lambda)P^p_\Lambda B^p_\Lambda P^p_\Lambda
                         \|_{L^2\to H_D^\sigma}
                   \leq C_{c,p,\sigma,A}\Lambda^{-A}.
\]
The same conclusion holds for each fixed smooth profile supported
strictly inside $(0,1)$, dilated in the norm variable by $\Lambda$,
and for the Hankel transforms of such dilates. For a family already
satisfying $I\psi=\psi(0)=0$, with the stated superpolynomial
projection error and polynomial ordinary local derivative bounds,
the principal constraints can be restored with the same arbitrarily
strong rate, using two projected profiles.
\end{proposition}
\begin{proof}
On radial $L^2(\R^2)$ the two soft operators are multiplication by
$p(|z|^2/\Lambda)$ and the Fourier multiplier
$p((\pi/c)^2|\xi|^2/\Lambda)$. Choose $b<1$ containing the support
of $p$. The latter multiplier is supported in the disk of radius
$(c/\pi)\sqrt{b\Lambda}$, while the hard band has radius
$(c/\pi)\sqrt\Lambda$. The Fourier transform of the former
multiplier is
$\Lambda\widehat{p(|\cdot|^2)}(\sqrt\Lambda\,\xi)$.
Its Schwartz decay and the fixed relative gap between these disks
give, by Schur's test,
\[
 \|(1-B_\Lambda)P^p_\Lambda B^p_\Lambda\|_{2\to2}
                                  =O_{c,p,M}(\Lambda^{-M})
\]
for every fixed $M$. The final soft time multiplication gives
zero hard time tail.
For a vector $f$ supported in $(0,\Lambda)$, the omitted-eigenvalue
bound gives directly
\[
 \|(1-\Pi_\Lambda)f\|_2^2
       \leq2\langle f,(1-A_\Lambda)f\rangle
       =2\|(1-B_\Lambda)f\|_2^2.
\]
This proves the raw superpolynomial estimate.

For any fixed $s<1/2$, the preceding proposition bounds the selected
part in $H_D^s$ by $C\Lambda^s$ times the raw input norm. The soft
part has the same bound: its first dilation derivative costs
$O(1+\Lambda)$ by the radial band estimate, and it vanishes before
the hard endpoint. Interpolate this coarse $H_D^s$ bound with the
raw estimate, taking $\sigma<s<1/2$ and then taking $M$ large.
This proves the operator estimate for every $A$.

Let $(D_r f)(x)=r^{-1/2}f(x/r)$. Direct substitution gives
$\mathcal H_cD_r=D_{1/r}\mathcal H_c$, so the Hankel transform
reverses the dilation generator and preserves every $H_D^s$ norm.
For a fixed smooth profile $\beta$ in the statement,
$\mathcal H_c\beta$ is Schwartz in the radial variable, hence has
every polynomial decay rate in the norm variable. The identity
$\mathcal H_cD_\Lambda\beta=D_{1/\Lambda}\mathcal H_c\beta$
therefore gives every polynomially small hard band tail, while
the raw time tail is zero. For its Hankel dual the roles reverse.
For a profile with a small raw time tail, first restrict it by
$P_\Lambda$; the resulting band tail increases by at most that time
tail. Thus the preceding projection bound is at most a fixed
constant times the sum of both raw tails. The same interpolation
proves the claimed fixed-profile conclusions.

Choose $\beta$ supported away from zero with $I\beta=1$.
The two profiles
\[
 b_{\Lambda,+}=D_\Lambda\beta,\qquad
 b_{\Lambda,-}=c^{-1}\mathcal H_cD_\Lambda\beta
\]
have identity matrix for the normalized constraints
$I/\sqrt\Lambda$ and $\operatorname{ev}_0/\sqrt\Lambda$:
indeed $(\mathcal H_cf)(0)=cIf$ and
$I\mathcal H_cf=f(0)/c$ for these Schwartz profiles.
Their projected constraint matrix tends to the identity with
every polynomial rate. For the mean, use Cauchy--Schwarz on the
time interval and the Schwartz exterior tail of the dual profile.
For evaluation, the projection error has every raw $L^2$ rate,
whereas its ordinary derivative on $(0,1)$ has polynomial growth
by the preceding band-extension estimate. The one-dimensional
trace inequality
\[
 |F(0)|^2\leq C\|F\|_{L^2(0,1)}
          \bigl(\|F'\|_{L^2(0,1)}+\|F\|_{L^2(0,1)}\bigr)
\]
proves the evaluation assertion. Subtracting the two projected
profiles multiplied by the inverse constraint matrix restores
the constraints exactly. The correction has every prescribed
$H_D^\sigma$ rate whenever the original constraint errors do.
\end{proof}

If a Hecke theta multiplier satisfies a fixed bound
$|L_K(1/2+i\tau,\chi)|\leq C(1+|\tau|)^\sigma$ for some
$\sigma<1/2$, the last proposition immediately gives theta graph
approximation, by Plancherel. Such a subconvex estimate is a
separate arithmetic input. Neither that estimate nor the purely
archimedean propositions prove the inverse theta Gram inequality
or the smoothed theta-tail estimate needed for a Hecke trace limit.
Proposition~\ref{prop:hecke-crude-gram},
Theorem~\ref{thm:hecke-deep-leakage}, and
Theorem~\ref{thm:hi-local-limit} establish those arithmetic and
local-limit comparisons for the deeper selection below.

\section{A high-concentration Hankel cutoff}
\label{sec:hankel-high-concentration}

The following selection is different from the fixed threshold
$1/2$ used in Section~\ref{sec:hankel-selection}. This difference is
part of the definition: the present cutoff retains only modes whose
concentration error is at most a prescribed negative power of
$\Lambda$. All constants in this section may depend on the fixed
Hankel parameter $c>0$.

On $L^2((0,\infty),dx)$ let
\[
 P_\Lambda=1_{(0,\Lambda)},\qquad
 B_\Lambda=\mathcal H_cP_\Lambda\mathcal H_c,\qquad
 A_\Lambda=P_\Lambda B_\Lambda P_\Lambda,
 \qquad F_\Lambda=P_\Lambda\mathcal H_cP_\Lambda.
\]
Thus $A_\Lambda=F_\Lambda^2$. Fix $M>0$, put
$\delta_\Lambda=\Lambda^{-M}$, and, for sufficiently large $\Lambda$,
define
\begin{equation}
 \label{eq:hankel-high-cutoff}
 \Pi_{\Lambda,M}=1_{[1-\delta_\Lambda,1]}(A_\Lambda),
 \qquad
 W_{\Lambda,M}=B_\Lambda
       (A_\Lambda|_{\operatorname{Ran}\Pi_{\Lambda,M}})^{-1}.
\end{equation}
The inverse in the last expression acts only on the selected space.
As before, $H_D^s$ denotes the Sobolev norm after the unitary change
of variables $Uf(t)=e^{t/2}f(e^t)$.

\begin{proposition}[Size, extension, and graph tail]
\label{prop:hankel-high-cutoff}
The selection in \eqref{eq:hankel-high-cutoff} is finite and
\begin{equation}
 \label{eq:hankel-high-rank}
 \operatorname{rank}\Pi_{\Lambda,M}
                 \sim\frac{2c\Lambda}{\pi}.
\end{equation}
Its selected eigenvalues are at least $1-\Lambda^{-M}>0$;
every omitted positive eigenvalue is strictly less than
$1-\Lambda^{-M}<1$. For $\psi$ in its range,
\begin{align}
 P_\Lambda W_{\Lambda,M}\psi&=\psi,\label{eq:high-right-inverse}\\
 \|W_{\Lambda,M}\psi\|_2^2
        &\leq(1-\delta_\Lambda)^{-1}\|\psi\|_2^2,\label{eq:high-extension-norm}\\
 \|(1-P_\Lambda)W_{\Lambda,M}\psi\|_2^2
        &\leq\frac{\delta_\Lambda}{1-\delta_\Lambda}
                                      \|\psi\|_2^2.\label{eq:high-raw-tail}
\end{align}
For every fixed $0\leq\sigma<s<1/2$,
\begin{equation}
 \label{eq:high-graph-tail}
 \|(1-P_\Lambda)W_{\Lambda,M}\psi\|_{H_D^\sigma}
 \leq C_{c,M,s,\sigma}
       \Lambda^{\sigma-\frac M2(1-\sigma/s)}\|\psi\|_2.
\end{equation}
In particular, any prescribed negative power on the right is
obtained by choosing a sufficiently large fixed $M$.
\end{proposition}

\begin{proof}
Finiteness follows from $\Tr A_\Lambda<\infty$. Since
$\Pi_{\Lambda,M}\leq1_{[1/2,1]}(A_\Lambda)$ for large $\Lambda$,
all the raw estimates in
Proposition~\ref{prop:hankel-raw-estimates} apply uniformly to this
smaller space. Directly,
\[
 \|W_{\Lambda,M}\psi\|_2^2
     =\langle\psi,A_\Lambda^{-1}\psi\rangle,
 \qquad P_\Lambda W_{\Lambda,M}\psi=\psi.
\]
Subtract $\|\psi\|_2^2$ from the first identity to obtain
\eqref{eq:high-raw-tail}.

For the global Sobolev bound, the finite Hankel transform commutes
with its square and therefore preserves every spectral subspace
used here. It is invertible on the selected space and
\[
 \|F_\Lambda^{-1}\Pi_{\Lambda,M}\|
                       \leq(1-\delta_\Lambda)^{-1/2}.
\]
Moreover
\begin{equation}
 \label{eq:high-extension-hankel}
 W_{\Lambda,M}\psi
            =\mathcal H_cP_\Lambda F_\Lambda^{-1}\psi.
\end{equation}
Indeed, $F_\Lambda A_\Lambda^{-1}=F_\Lambda^{-1}$ on this
subspace. The identity
$\mathcal H_cD_r=D_{1/r}\mathcal H_c$ shows that $\mathcal H_c$
is an isometry on every $H_D^s$. Apply the earlier raw estimate
to $F_\Lambda^{-1}\psi$ to get
\[
 \|W_{\Lambda,M}\psi\|_{H_D^s}
                         \leq C_{c,s}\Lambda^s\|\psi\|_2.
\]
Multiplication by $1_{(\log\Lambda,\infty)}$ is bounded on
$H^s(\mathbb R)$, uniformly in the endpoint, for $0<s<1/2$.
One direct proof uses the fractional difference quotient formula:
the cross-boundary contribution is a constant multiple of
$\int_0^\infty |f(L+u)|^2u^{-2s}du$, which is bounded by the
$H^s$ norm by the one-dimensional fractional Hardy inequality.
For completeness, apply the difference quotient formula on
$I_j=(2^{-j-1},2^{-j})$, $j\geq0$, and let $m_j$ be the mean
of $f(L+\cdot)$ on $I_j$. The oscillation terms satisfy
\[
 \sum_{j\geq0}|I_j|^{-2s}
       \int_{I_j}|f(L+u)-m_j|^2du
                    \leq C_s[f]_{\dot H^s}^2.
\]
The difference of two adjacent means is bounded by their
cross-interval difference quotient integral times
$C_s|I_j|^{2s-1}$. Put $\rho=2^{-(1-2s)}<1$.
The elementary discrete inequality
\[
 \sum_{j\geq0}\rho^j|m_j|^2
 \leq C_\rho\left(|m_0|^2+
                \sum_{j\geq0}\rho^j|m_{j+1}-m_j|^2\right)
\]
follows by writing each $m_j$ as a partial sum and applying
the $\ell^1*\ell^2$ convolution bound to the sequence
$\rho^{j/2}$. This controls the mean terms by the same
fractional seminorm and an $L^2$ term. The region $u\geq1$
is controlled directly by $\|f\|_2^2$. These estimates prove
the asserted Hardy inequality, and explain the restriction
$2s<1$.
Consequently the exterior part also has $H_D^s$ norm
$O_{c,s}(\Lambda^s\|\psi\|_2)$. Interpolate this estimate with
\eqref{eq:high-raw-tail}, using the Fourier definition of the
Sobolev norms and H\"older's inequality. The resulting powers are
$\Lambda^{-M(1-\sigma/s)/2}\Lambda^\sigma$, as asserted.

It remains to prove the lower rank asymptotic; a transition-width
theorem is unnecessary. Fix $0<a<1$, and choose a smooth function
$h$, equal to one on $[0,a]$, supported in $[0,b]$ for some $b<1$,
with $0\leq h\leq1$. Put $p=h^2$ and
\[
 K^p_\Lambda=P^p_\Lambda B^p_\Lambda P^p_\Lambda,
 \quad P^p_\Lambda=p(x/\Lambda),
 \quad B^p_\Lambda=\mathcal H_cp(x/\Lambda)\mathcal H_c.
\]
Then $0\leq K^p_\Lambda\leq1$. The strict-margin Fourier argument
in Proposition~\ref{prop:hankel-interior-approximation} gives
\[
 \|(1-B_\Lambda)P^p_\Lambda(B^p_\Lambda)^{1/2}\|_{2\to2}
                                         =O_N(\Lambda^{-N})
\]
for every $N$, since $(B^p_\Lambda)^{1/2}
=\mathcal H_ch(x/\Lambda)\mathcal H_c$ has the same strict
frequency margin. In addition,
\[
 \|(B^p_\Lambda)^{1/2}P^p_\Lambda\|_{\rm HS}^2
    =\Tr(P^p_\Lambda B^p_\Lambda P^p_\Lambda)
    \leq\Tr A_\Lambda=O_c(\Lambda).
\]
The inequality follows from the diagonal kernel, since both
soft factors lie between zero and their corresponding hard
cutoffs. Hence
$\|(1-B_\Lambda)K^p_\Lambda\|_{\rm HS}=O_N(\Lambda^{-N})$
after relabelling $N$.

For every vector $f$ supported in $(0,\Lambda)$ the omitted
spectral gap gives
\begin{equation}
 \label{eq:high-projection-gap}
 \|(1-\Pi_{\Lambda,M})f\|_2^2
 \leq\delta_\Lambda^{-1}\langle f,(1-A_\Lambda)f\rangle
 =\delta_\Lambda^{-1}\|(1-B_\Lambda)f\|_2^2.
\end{equation}
Apply this to the columns of $K^p_\Lambda$ and sum. Thus
\begin{equation}
 \label{eq:high-soft-HS}
 \|(1-\Pi_{\Lambda,M})K^p_\Lambda\|_{\rm HS}
                                      =O_N(\Lambda^{-N})
\end{equation}
for every fixed $N$. Since $K^p_\Lambda$ is a contraction,
\[
 \operatorname{rank}\Pi_{\Lambda,M}
 \geq\|\Pi_{\Lambda,M}K^p_\Lambda\|_{\rm HS}^2
 =\|K^p_\Lambda\|_{\rm HS}^2-o(1).
\]
Compress $K^p_\Lambda$ to $(0,a\Lambda)$. The time multiplier is
then the identity and its band multiplier is at least
$B_{a\Lambda}$, so this compression is at least $A_{a\Lambda}$.
For positive Hilbert--Schmidt operators, order implies order of
their squared traces: if $T=S+D$ with $S,D\geq0$, then
$\Tr T^2-\Tr S^2=2\Tr(SD)+\Tr D^2\geq0$.
It follows that
\[
 \operatorname{rank}\Pi_{\Lambda,M}
 \geq\Tr(A_{a\Lambda}^2)-o(1)
 =\frac{2ca\Lambda}{\pi}+O_{c,a}(\log\Lambda).
\]
Here the last equality uses the trace and variance estimates in
Proposition~\ref{prop:hankel-selection-size}.
On the other hand,
\[
 \operatorname{rank}\Pi_{\Lambda,M}
       \leq\frac{\Tr A_\Lambda}{1-\delta_\Lambda}
       =\frac{2c\Lambda}{\pi}+o(\Lambda).
\]
Let $a$ increase to one after taking the lower limit. This proves
\eqref{eq:hankel-high-rank}.
\end{proof}

\begin{proposition}[Interior approximation for the new cutoff]
\label{prop:hankel-high-interior}
For fixed $M>0$, the strict-margin and fixed-profile conclusions
of Proposition~\ref{prop:hankel-interior-approximation} remain
valid with $\Pi_{\Lambda,M}$ in place of $\Pi_\Lambda$:
for every $\sigma<1/2$ and every $N>0$,
\[
 \|(1-\Pi_{\Lambda,M})P^p_\Lambda B^p_\Lambda P^p_\Lambda
                          \|_{L^2\to H_D^\sigma}
                                  =O_{c,p,M,\sigma,N}(\Lambda^{-N}).
\]
The projected two-profile repair of the principal mean and value
constraints is also valid. Imposing these constraints changes
the rank by at most two, so it preserves
\eqref{eq:hankel-high-rank}.
\end{proposition}
\begin{proof}
The only change in the raw approximation proof is that the
omitted spectral gap $1/2$ is replaced by $\delta_\Lambda$.
Equation~\eqref{eq:high-projection-gap} therefore costs
$\Lambda^{M/2}$ in norm. Every strict-margin and fixed-profile
raw error has an arbitrarily large polynomial decay rate, so
this fixed loss is absorbed by choosing that rate larger.
The coarse $H_D^s$ bound remains $O(\Lambda^s)$, since the new
selected space is contained in the old one. The same interpolation
argument proves the stated graph estimate. The mean and value
of the projected repair profiles converge to their original
normalized values with every polynomial rate, by the same
Cauchy--Schwarz and ordinary trace estimates as before. Their
constraint matrix is consequently invertible for large $\Lambda$.
\end{proof}

\begin{corollary}[An abstract theta-tail criterion]
\label{cor:hankel-high-theta-tail}
Suppose a fixed theta multiplier $E$ obeys
\[
 \|Ef\|_2\leq C\|f\|_{H_D^\sigma},\qquad 0\leq\sigma<1/2.
\]
Let $\Pi^0_{\Lambda,M}$ be either the full high-concentration
selection or its subspace satisfying the two principal
constraints. Suppose that a lower-tail operator $T_{\Lambda,g}$
has the exact identity
\[
 T_{\Lambda,g}\psi
 =-V(g)1_{(-\infty,-\log\Lambda)}
                   E(1-P_\Lambda)W_{\Lambda,M}\psi
       \qquad(\psi\in\operatorname{Ran}\Pi^0_{\Lambda,M}),
\]
where $V(g)$ is bounded on the output $L^2$ space, uniformly in
$\Lambda$. If the theta Gram operator on the selected space
satisfies $\|G_\Lambda^{-1}\|=O(\Lambda^a)$ for some fixed
$a\geq0$, then
\begin{equation}
 \label{eq:high-normalized-tail}
 \|T_{\Lambda,g}\Pi^0_{\Lambda,M}G_\Lambda^{-1/2}\|_{\rm HS}^2
 \ll\Lambda^{a+1+2\sigma-M(1-\sigma/s)},
                 \qquad \sigma<s<1/2.
\end{equation}
In particular the normalized tail tends to zero whenever
$M(1-\sigma/s)>a+1+2\sigma$. With the crude exponent $a=2$,
it is enough to take $M>(3+2\sigma)/(1-\sigma/s)$.
For example, the choices
\[
 \sigma=3/8,\qquad s=7/16,\qquad M=28,
\]
give an $O(\Lambda^{-1/4})$ bound in
\eqref{eq:high-normalized-tail} when $a=2$.
\end{corollary}
\begin{proof}
The assumed identity and \eqref{eq:high-graph-tail} bound the
operator norm of $T_{\Lambda,g}$ on the selected space by
$C\Lambda^{\sigma-M(1-\sigma/s)/2}$. Its dimension is
$O(\Lambda)$. Multiply the squared operator bound by this
dimension and by $\|G_\Lambda^{-1}\|$.
\end{proof}

The corollary isolates the arithmetic inputs: a fixed exponent
$\sigma<1/2$, an exact band-tail identity, and a polynomial inverse
Gram bound. It makes no cancellation assumption for ordinary
M\"obius sums. It proves no estimate for the fixed-threshold
$1/2$ selection, which is a different cutoff. That selection, and
every other fixed threshold in $(0,1)$, will be treated in
Theorem~\ref{nh:theorem}.


\section{Arithmetic estimates for a finite-order Eisenstein-field cutoff}
\label{sec:hecke-arithmetic-attempt}

This section treats the arithmetic part of a radial Hankel cutoff for
$K=\mathbb Q(\sqrt{-3})$.  It distinguishes the selection at concentration
level $1/2$ from the explicitly different selection at level
$1-\Lambda^{-M}$.  The latter admits a crude arithmetic Gram estimate
and a rapidly vanishing normalized theta tail.  These statements do not
by themselves identify the matching local upper limit.
Section~\ref{sec:hecke-original-upper} proves that remaining
local comparison.

Write $\mathcal O=\mathbb Z[\omega]$, with $\omega=e^{2\pi i/3}$,
and let $v=\sqrt3/2$ be its Euclidean covolume.  Fix a primitive
finite-order Hecke character $\chi$ of conductor $\mathfrak f$.  Put
\[
 Q_\chi=3N\mathfrak f,\qquad
 c=\frac{2\pi}{\sqrt{Q_\chi}},\qquad
 r_K=\frac{\pi}{3\sqrt3},\qquad
 \delta_\chi=\mathbf1_{\chi=1}.
\]
All constants may depend on this fixed character.  The raw Hilbert
space is $L^2((0,\infty),du)$, and
\[
 (H_c\phi)(y)=c\int_0^\infty J_0(2c\sqrt{uy})\phi(u)\,du,
 \qquad I\phi=\int_0^\infty\phi(u)\,du,
\]
\[
 E_\chi\phi(x)=\sqrt x\sum_{\mathfrak a}
       \chi(\mathfrak a)\phi(N\mathfrak a\,x)
       -\delta_\chi r_K x^{-1/2}I\phi.
\]
The multiplicative variable $x$ becomes logarithmic time $t=\log x$.
The raw dilation Sobolev space $H_D^s$ is ordinary $H^s(\mathbb R)$
under $U\phi(t)=e^{t/2}\phi(e^t)$.

\subsection{The lattice theta identity and a Riesz estimate}

\begin{proposition}[Finite-order lattice normalization]
\label{prop:hecke-lattice-poisson}
There is a constant $\varepsilon_\chi$ of modulus one such that
\[
 E_\chi\phi(x)=\varepsilon_\chi
                  E_{\bar\chi}(H_c\phi)(x^{-1}).
\]
For the principal character $\varepsilon_1=1$ and
\[
 H_c\phi(0)=cI\phi,\qquad I(H_c\phi)=\phi(0)/c.
\]
Define the norm-counting discrepancy
\[
 A_\chi(X)=\sum_{N\mathfrak a\leq X}\chi(\mathfrak a)
                          -\delta_\chi r_KX.
\]
Its first integral has the absolutely convergent expansion
\begin{equation}\label{eq:hecke-riesz}
 R_\chi(X):=\int_0^X A_\chi(y)\,dy
 =-\frac{\delta_\chi X}{6}
 +\frac{\varepsilon_\chi X}{c}
   \sum_{\mathfrak a}\frac{\bar\chi(\mathfrak a)}{N\mathfrak a}
                  J_2(2c\sqrt{XN\mathfrak a}).
\end{equation}
Consequently
\[
 R_\chi(X)=-\delta_\chi X/6+O_\chi(X^{3/4})\quad(X\geq1).
\]
Also $A_\chi(X)=O_\chi(X^{1/3})$ for $X\geq1$.
\end{proposition}

\begin{proof}
The class number is one, and each nonzero principal ideal has six
generators.  Write $\mathfrak f=(F)$.  The ideal character is represented
by a multiplicative residue character on $\mathcal O/(F)$ which is
trivial on the six units and zero on nonunits.  Hence the ideal theta
sum is one sixth of the corresponding weighted lattice sum.  The dual
lattice of $\mathcal O$ is $(i/v)\mathcal O$; the dual of
$F\mathcal O$ is $(i/(v\bar F))\mathcal O$.
Finite Fourier transformation of a primitive residue character gives
its conjugate times a Gauss sum of modulus $\sqrt{N\mathfrak f}$.
For a radial function $\phi(|z|^2x)$, Euclidean Fourier transformation
gives
\[
 \widehat{\phi(|\cdot|^2x)}(\xi)
   =\frac\pi x\int_0^\infty
             J_0(2\pi|\xi|\sqrt{u/x})\phi(u)\,du.
\]
Since $|\xi|=c\sqrt{N\mathfrak a}/\pi$ on the relevant dual lattice,
the Gauss factor and covolume combine to give
$\varepsilon_\chi x^{-1}H_c\phi(N\mathfrak a/x)$.
This proves the nonprincipal identity.  In the principal case the
two lattice origins give
\[
 6\sum_{\mathfrak a}\phi(N\mathfrak a x)+\phi(0)
 =x^{-1}\left(6\sum_{\mathfrak a}
 H_c\phi(N\mathfrak a/x)+cI\phi\right).
\]
Here $c/6=r_K$.  Subtracting the indicated mean proves the corrected
identity.  The two identities at zero follow from $H_c^2=I$.

Apply the same lattice formula to $(X-|z|^2)_+$.  At a nonzero
frequency $\xi$ its Fourier transform is
\[
 \frac{X}{\pi|\xi|^2}J_2(2\pi\sqrt X|\xi|),
\]
as follows by integrating $2\pi r(X-r^2)J_0(2\pi r|\xi|)$.
At zero the integral is $\pi X^2/2$.  The latter is the main term
$r_KX^2/2$ after division by the six units and the covolume.
Removing the original lattice origin contributes $-X/6$ in the
principal case.  Grouping the remaining dual points into ideals
proves \eqref{eq:hecke-riesz}.  Since
$J_2(z)=O(z^{-1/2})$ for $z\geq1$, each summand is bounded by a
constant times $X^{3/4}(N\mathfrak a)^{-5/4}$.  The latter series
converges, because the number of ideals of norm at most $Y$ is
$O_K(Y)$.

For completeness the weaker pointwise discrepancy estimate follows
from elementary smoothing.  Count points of each of the finitely
many residue-class lattices in a disk of radius $T=\sqrt X$.
Convolve indicators of the disks of radii $T\pm\eta$ with a
nonnegative smooth mollifier supported in a disk of radius $\eta$.
They respectively majorize and minorize the sharp disk indicator.
Their zero Fourier coefficients differ from the disk area by
$O(T\eta+\eta^2)$.  At nonzero dual frequencies the disk transform
is $O(T^{1/2}|\xi|^{-3/2})$.  The mollifier contributes
$O_N((1+\eta|\xi|)^{-N})$, so summation over a two-dimensional
dual lattice gives $O(T^{1/2}\eta^{-1/2})$.
Choose $\eta=T^{-1/3}$.  This proves discrepancy $O(T^{2/3})$
for each translated lattice.  Taking the finite character-weighted
combination proves $A_\chi(X)=O_\chi(X^{1/3})$.
\end{proof}

\begin{proposition}[Improper Mellin integrals at sharp endpoints]
\label{prop:hecke-improper-mellin}
Let $\phi$ be smooth on $[0,A]$ and zero outside that interval; a
jump at $A$ is allowed.  For $0<\Re s<1$,
\begin{equation}\label{eq:hecke-improper-mellin}
 \lim_{\eta\downarrow0}\int_\eta^\infty
 E_\chi\phi(x)x^{s-1/2}\frac{dx}{x}
 =L_K(s,\chi)\int_0^A\phi(u)u^{s-1}\,du.
\end{equation}
The convergence is locally uniform in this strip, and any fixed
number of derivatives in $s$ may be taken.  Absolute convergence at
zero is not asserted.
\end{proposition}

\begin{proof}
First take $\phi=1_{(0,a)}$.  Its corrected theta image is
$\sqrt x A_\chi(a/x)$.  For a nonprincipal character, integration
by parts against $R_\chi$ gives, uniformly on compact subsets of
$\Re s>-1/4$,
\[
 \int_Y^\infty A_\chi(y)y^{-s-1}\,dy
 =-R_\chi(Y)Y^{-s-1}
 +(s+1)\int_Y^\infty R_\chi(y)y^{-s-2}\,dy
 =O(Y^{-\Re s-1/4}).
\]
For the principal character separate the term $-y/6$ in $R_1$;
its contribution to the displayed integral is $O(Y^{-\Re s})$.
Thus the integral converges for every $\Re s>0$.  The same estimates
with fixed powers of $\log y$ prove locally uniform convergence of
its derivatives.  The change of variables $y=a/x$ is legitimate in
each truncated integral.  For a nonprincipal character,
Stieltjes summation first gives
$L_K(s,\chi)/s=\int_1^\infty A_\chi(y)y^{-s-1}\,dy$
when $\Re s>1$.  The preceding locally uniform convergence
continues this identity to $\Re s>0$.  Since $A_\chi(y)=0$
for $0<y<1$, multiplication by $a^s$ proves the disk identity.

For the principal character there is no common region of absolute
convergence for the whole corrected disk integral and the original
Dirichlet series.  Instead, Stieltjes summation for $\Re s>1$ gives
\[
 \frac{L_K(s,1)}s
 =\int_1^\infty A_1(y)y^{-s-1}\,dy+\frac{r_K}{s-1}.
\]
The integral on the right continues holomorphically to $\Re s>0$
by the Riesz estimate already proved.  On $0<\Re s<1$ one has
$A_1(y)=-r_Ky$ for $0<y<1$, and therefore
\[
 \frac{r_K}{s-1}=\int_0^1 A_1(y)y^{-s-1}\,dy.
\]
Combining these identities proves
$L_K(s,1)/s=\int_0^\infty A_1(y)y^{-s-1}\,dy$
in this strip, and hence the principal disk identity after multiplying
by $a^s$.  The restriction $\Re s<1$ is exactly the convergence
condition for the mean-subtraction tail at $x=\infty$.

For a general sharp profile use the exact representation
\[
 \phi(u)=\phi(A)1_{u<A}
                -\int_0^A\phi'(a)1_{u<a}\,da.
\]
The omitted lower Mellin integral for a nonprincipal disk is
$O(\eta^{\Re s+1/4}a^{-1/4})$; the same estimate remains valid
when $a<\eta$ by using the complete disk integral.  The principal
case has in addition $O(\eta^{\Re s})$.  Fixed derivatives only
insert fixed powers of $|\log\eta|$ and $|\log a|$.
Because $a^{-1/4}(1+|\log a|)^j$ is integrable at zero, these bounds
justify integration against $\phi'(a)$, the interchange of the
truncated integrals, and passage to their limits.  They also justify
all fixed Mellin jets.  This proves the proposition.
\end{proof}

The improper Mellin statement must not be replaced by a convergence
claim for the ordinary Dirichlet-series tail.  In fact the exact
analytically regularized tail is
\begin{equation}\label{eq:hecke-tail-endpoint}
 L_K(s,\chi)-\sum_{N\mathfrak a\leq X}
    \frac{\chi(\mathfrak a)}{(N\mathfrak a)^s}
 +\delta_\chi\frac{r_KX^{1-s}}{1-s}
 =-A_\chi(X)X^{-s}
           +s\int_X^\infty A_\chi(y)y^{-s-1}\,dy.
\end{equation}
The endpoint discrepancy survives, even though the integral on the
right has the stronger Riesz bound.

\subsection{A graph domain and the crude Gram estimate}

We use the explicit analytic input
\begin{equation}\label{eq:hecke-subconvex-input}
 |L_K(1/2+i\tau,\chi)|\leq C_\chi(2+|\tau|)^{3/8}.
\end{equation}
For a cuspidal finite-order theta series, including the cubic
character of conductor $(269)$, this follows from the fixed-form
Weyl bound of \citet[Theorem~1.1, p.~2]{booker2019}, which gives
$|\tau|^{1/3}\log|\tau|$ for a primitive holomorphic cusp form of
arbitrary level.  The theorem states no exclusion of weight one.
In a reducible theta sector the $L$-function is a product of two
fixed Dirichlet $L$-functions.  The all-conductor Weyl bound of
\citet[Theorem~1.1, p.~1]{petrowyoung2023} gives exponent
$1/6+\epsilon$ for each factor, hence $1/3+2\epsilon<3/8$
after choosing a fixed sufficiently small $\epsilon$.
Thus it gives the same conclusion.  A proof using this section must retain
this subconvexity input; degree-two convexity alone reaches the
unsuitable endpoint exponent $1/2$.

Plancherel and \eqref{eq:hecke-subconvex-input} give
\begin{equation}\label{eq:hecke-graph}
 \|E_\chi\phi\|_2\leq C_\chi\|\phi\|_{H_D^{3/8}}.
\end{equation}
This defines the closed theta operator on that Sobolev space.
For sharp time-supported smooth profiles it agrees almost
everywhere with the displayed finite theta sum: smooth the endpoint
in $H_D^{3/8}$, use \eqref{eq:hecke-graph} for $L^2$ convergence,
and observe convergence of the finite sum at every $x$ outside the
countable set of endpoint ratios.  The principal means converge by
ordinary $L^2$ approximation on a fixed compact interval.
The preceding improper Mellin proposition gives the full complex
Mellin identity and its jets for these same profiles.

Set
\[
 P_\Lambda=1_{(0,\Lambda)},\qquad
 B_\Lambda=H_cP_\Lambda H_c,\qquad
 A_\Lambda=P_\Lambda B_\Lambda P_\Lambda.
\]
Let $\mathcal R_\Lambda$ be any spectral subspace of $A_\Lambda$
with eigenvalues at least $\nu>0$, with the two constraints
$I\phi=\phi(0)=0$ imposed in the principal sector.  We use the
archimedean mass estimate
\begin{equation}\label{eq:hecke-core-mass}
 \int_{d/(c^2\Lambda)}^\Lambda|\phi(u)|^2du
                         \geq\tfrac12\|\phi\|_2^2
 \qquad(\phi\in\mathcal R_\Lambda).
\end{equation}
It follows directly by extending $\phi$ to a Hankel-band vector
of norm at most $\nu^{-1/2}\|\phi\|$ and using
$|w(u)|^2\leq c^2\Lambda\|w\|^2$ near zero.

\begin{proposition}[Crude Hecke Gram control]
\label{prop:hecke-crude-gram}
For a raw orthonormal basis of $\mathcal R_\Lambda$, let $G_\Lambda$
be the Gram matrix of its theta images.  Then
\[
 \|G_\Lambda^{-1}\|\leq C_{\chi,\nu}\Lambda^2.
\]
No cancellation estimate for the ideal M\"obius function is required.
\end{proposition}

\begin{proof}
Put $a=d/(c^2\Lambda)$ and $N=\lceil\Lambda/a\rceil$.
For $f=U\phi$ and $e=E_\chi\phi$, finite ideal M\"obius inversion
gives exactly, for $\log a\leq t\leq\log\Lambda$,
\[
 f(t)=\sum_{N\mathfrak a\leq N}
  \frac{\mu_K(\mathfrak a)\chi(\mathfrak a)}{\sqrt{N\mathfrak a}}
                         e(t+\log N\mathfrak a).
\]
Expansion groups terms by the product ideal $\mathfrak b$ and uses
$\sum_{\mathfrak a\mid\mathfrak b}\mu_K(\mathfrak a)
=\mathbf1_{\mathfrak b=\mathcal O}$.  The support of $\phi$ removes
all products beyond $\Lambda e^{-t}\leq N$.  In the principal sector
$I\phi=0$ removes the correction term, so the same finite identity
holds.  Its translation operator has $L^2$ norm at most
\[
 \sum_{N\mathfrak a\leq N}\frac1{\sqrt{N\mathfrak a}}
                         \leq C_K\sqrt N\leq C_\chi\Lambda,
\]
by the elementary ideal count $O_K(Y)$ and partial summation.
The core mass estimate therefore yields
$\|\phi\|_2\leq C_\chi\Lambda\|E_\chi\phi\|_2$.
Squaring is precisely the asserted Gram inequality.
\end{proof}

\subsection{Deep concentration and unconditional lower estimates}

For a fixed $M>0$ make the explicit choice
\[
 \Pi_{\Lambda,M}
       =\mathbf1_{[1-\Lambda^{-M},1]}(A_\Lambda).
\]
The principal raw column space is its intersection with
$\ker I\cap\ker\operatorname{ev}_0$.  This is a different cutoff
from $\mathbf1_{[1/2,1]}(A_\Lambda)$.  Let $E_\Lambda$ denote the
theta column operator, $Z_\Lambda=E_\Lambda G_\Lambda^{-1/2}$,
$Q_\Lambda=Z_\Lambda Z_\Lambda^*$, and
$S_L=1_{[-L,L]}$, $L=\log\Lambda$.

\begin{theorem}[Deep-cut theta leakage]
\label{thm:hecke-deep-leakage}
Assume the graph bound \eqref{eq:hecke-graph} and the archimedean
bound $\|\psi\|_{H_D^s}\leq C_s\Lambda^s\|\psi\|_2$ on the
selected raw space for $s<1/2$.  More generally let the graph exponent
be $\sigma<s<1/2$.  Then
\begin{equation}\label{eq:hecke-deep-leakage}
 \|(1-S_L)Z_\Lambda\|_2
 \leq C\Lambda^{3/2+\sigma-(M/2)(1-\sigma/s)}.
\end{equation}
In particular, $\sigma=3/8$, $s=7/16$, and $M=56$ give
$O(\Lambda^{-17/8})$.  For this choice,
\[
 \operatorname{Tr}((Q_\Lambda-S_LQ_\Lambda S_L)V(h))\longrightarrow0
 \qquad(h\in C_c^\infty(\mathbb R)).
\]
\end{theorem}

\begin{proof}
Write $F_\Lambda=P_\Lambda H_cP_\Lambda$, so $A_\Lambda=F_\Lambda^2$.
On the unconstrained deep selected span put
$W=B_\Lambda A_\Lambda^{-1}$, so $P_\Lambda W=I$.  Spectral calculus
gives
\[
 \|(1-P_\Lambda)W\psi\|_2^2
 =\langle\psi,(A_\Lambda^{-1}-I)\psi\rangle
 \leq C\Lambda^{-M}\|\psi\|_2^2.
\]
Furthermore $H_cW\psi=F_\Lambda^{-1}\psi$ belongs to the same
selected time-supported span.  Hankel transformation reverses the
dilation generator and preserves its Sobolev norms.  Thus both
$W\psi$ and $\psi$ have $H_D^s$ norm $O(\Lambda^s)\|\psi\|$.
Interpolation gives
\[
 \|(1-P_\Lambda)W\psi\|_{H_D^\sigma}
 \leq C\Lambda^{\sigma-(M/2)(1-\sigma/s)}\|\psi\|_2.
\]
The selected dimension is $O(\Lambda)$, since it is bounded by
twice $\operatorname{Tr}A_\Lambda$ for large $\Lambda$.  Sum the
squared estimate over a raw orthonormal basis and apply the graph
bound.  The identity $E_\chi H_c=\varepsilon_\chi J E_{\bar\chi}$,
where $(Jf)(t)=f(-t)$, extends to $H_D^\sigma$ because $H_c$
preserves that space and both theta maps satisfy the graph bound.
Apply it to $H_cW\psi=F_\Lambda^{-1}\psi$, which is time supported.
In the principal case its mean is
$I(H_cW\psi)=W\psi(0)/c=\psi(0)/c=0$.
It follows directly that the lower theta tail of $W\psi$ is zero.
The upper theta tail of $\psi$ is zero,
using $I\psi=0$ in the principal case.  Consequently
$\|(1-S_L)E_\Lambda\|_2$ is bounded by the theta norm just estimated.
Finally $\|G_\Lambda^{-1/2}\|=O(\Lambda)$ proves
\eqref{eq:hecke-deep-leakage}.

For the trace statement use
$Q-S_LQS_L=(1-S_L)Q+S_LQ(1-S_L)$ and Hilbert--Schmidt
Cauchy--Schwarz.  Since $\|Z_\Lambda\|_2=O(\Lambda^{1/2})$,
the absolute trace difference is
$O(\|V(h)\|\Lambda^{1/2}\|(1-S_L)Z_\Lambda\|_2)=o(1)$.
\end{proof}

\begin{proposition}[All zero jets for the deep cutoff]
\label{prop:hecke-deep-zero-jets}
Let $M>4/3$ and let $\mathcal V_L$ be the span of the truncated
exponential-polynomial jets associated with any fixed finite packet
of nontrivial zeros of $L_K(s,\chi)$, including their multiplicities.
Then, for some fixed exponent $C$ depending on the packet,
\[
 \|Q_\Lambda P_{\mathcal V_L}\|_2^2
                       =O(\Lambda^{4/3-M}(1+L)^C)=o(1).
\]
The resulting finite-packet lower bound is the harmonic-measure
bound: a zero $\rho=\beta+i\gamma$ contributes multiplicity times
the point mass at $\gamma$ when $\beta=1/2$, and otherwise the
Poisson measure of width $|\beta-1/2|$.  No zero-location theorem
is used.
\end{proposition}

\begin{proof}
For a zero $s_0=\beta+i\gamma$ and $j<m_{s_0}$, define the lower
Mellin tail with weight $(\log x+L)^j$.  Proposition
\ref{prop:hecke-improper-mellin} justifies this improper integral.
Its raw kernel is
\[
 H_{s_0,j,\Lambda}(u)=
 \left.(\partial_s+L)^j\left[
 u^{s-1}\left(L_K(s,\chi)
       -\sum_{N\mathfrak a\leq\Lambda u}
            \frac{\chi(\mathfrak a)}{(N\mathfrak a)^s}\right)
 +\delta_\chi\frac{r_K\Lambda^{1-s}}{1-s}
 \right]\right|_{s=s_0}.
\]
The formula follows on a truncated compact raw core by Stieltjes
summation; the Riesz bounds in the preceding proof justify its
limit and all fixed derivatives.  For $u<\Lambda^{-1}$, the zero
and its jets remove the $L$-term.  Thus the kernel is zero there
in a nonprincipal sector, and bounded there in the principal
sector.  For $u\geq\Lambda^{-1}$, apply
\eqref{eq:hecke-tail-endpoint}, using the pointwise bound
$A_\chi(X)=O(X^{1/3})$ for its endpoint term and the Riesz estimate
for its conditionally convergent integral.  The latter use is
necessary when $\beta\leq1/3$.  The derivatives involve scale-free
logarithms, so for every fixed jet
\[
 |H_{s_0,j,\Lambda}(u)|
       \leq C\Lambda^{1/3-\beta}u^{-2/3},
 \qquad
 \|1_{u>\Lambda}H_{s_0,j,\Lambda}\|_2^2
       \leq C\Lambda^{1/3-2\beta}.
\]
The global kernel belongs to raw $L^2$.  First apply its identity
to inverse Hankel transforms of smooth compact band profiles.
These are radial Schwartz functions; the disk-superposition proof
extends to them because their derivatives have rapid decay, so all
the same Riesz majorants remain integrable over $0<a<\infty$.
Poisson makes their lower theta tail zero, with the value constraint
in the principal case.  Such vectors are dense in the fixed band
space, and the principal value constraint can be retained because
evaluation there is continuous.  Continuity of the $L^2$ kernel
therefore proves annihilation on $W\psi$, without an assertion of
absolute Mellin convergence.
Therefore, for a raw orthonormal selected basis,
\[
 \sum_i|\ell_{s_0,j,\Lambda}(\psi_i)|^2
 \leq\|(1-P_\Lambda)W\|^2
         \|1_{u>\Lambda}H_{s_0,j,\Lambda}\|_2^2
 \leq C\Lambda^{-M+1/3-2\beta}.
\]
The inverse Gram bound changes this to
$O(\Lambda^{7/3-M-2\beta})$ for normalized theta columns.
The full Mellin jet vanishes and the upper theta tail is zero.
For $\alpha=\beta-1/2<0$, use endpoint vectors
$1_{[-L,L]}e^{(\alpha-i\gamma)(t+L)}(t+L)^j$; their extra squared
factor $\Lambda^{2\alpha}$ leaves $\Lambda^{4/3-M}$.
For $\alpha>0$, use the right-endpoint normalization and expand
$(t-L)^j$ in powers of $(t+L)$; the power is
$\Lambda^{10/3-M-4\beta}\leq\Lambda^{4/3-M}$, with fixed
logarithmic factors.  Central-line polynomial jets have their
usual polynomial $L$ normalization and satisfy the same bound.
The finite endpoint Gram matrices have nondegenerate limits, so
the estimate passes to the orthogonal projection onto their span.
The finite exponential-packet trace calculation then gives the
stated harmonic measures.  Increasing finite packets gives the
corresponding all-zero lower bound whenever the ordinary zero-count
bound is used to sum the resulting positive measures.
\end{proof}


\section{Local equality and uniform frequency bounds over the complex field}
\label{sec:hecke-original-upper}

Let $K=\mathbb Q(\sqrt{-3})$, let $\chi$ be a fixed primitive finite
Hecke character of conductor $\mathfrak f$, and put
\[
 c=\frac{2\pi}{\sqrt{3N\mathfrak f}},\qquad
 (\mathcal H_cf)(u)=c\int_0^\infty
                   J_0(2c\sqrt{uv})f(v)\,dv.
\]
The raw Hilbert space is $L^2((0,\infty),du)$.
The transform $\mathcal H_c$ is a unitary involution. Write
\[
 (Uf)(t)=e^{t/2}f(e^t),\qquad H_D^s=U^{-1}H^s(\mathbb R),
 \qquad (D_rf)(u)=r^{-1/2}f(u/r).
\]
Our Fourier convention on logarithmic space is
$\widehat g(\tau)=\int g(t)e^{i\tau t}dt$, with Plancherel measure
$d\tau/(2\pi)$. Let $C_a$ be the multiplier with symbol $a$.
The theta map has the closed multiplier realization
\[
 \widehat{E_\chi f}(\tau)
       =L_K(1/2+i\tau,\chi)\widehat{Uf}(\tau),
 \qquad
 E_\chi\mathcal H_c=\varepsilon_\chi J E_{\bar\chi},
 \quad (Jg)(t)=g(-t),\quad |\varepsilon_\chi|=1.
\]

Put $L=\log\Lambda$, $P_\Lambda=1_{[0,\Lambda]}$,
$B_\Lambda=\mathcal H_cP_\Lambda\mathcal H_c$, and
\[
 \Pi_\Lambda=
 1_{[1-\Lambda^{-56},1]}
       (P_\Lambda\mathcal H_cP_\Lambda\mathcal H_cP_\Lambda).
\]
For the principal character use
$\operatorname{Ran}\Pi_\Lambda\cap\ker I\cap\ker\operatorname{ev}_0$,
where $If=\int_0^\infty f(u)du$. Let $Q_\Lambda$ be the orthogonal
projection onto the corresponding theta image, with its full Gram
normalization, and set
\[
 S_L=1_{[-L,L]},\qquad D_L=S_L-S_LQ_\Lambda S_L.
\]
If $z_i$ is any orthonormal basis of that image, define
\[
 d\mu_\Lambda(\tau)=
 \left(2L-\sum_i|\widehat{S_Lz_i}(\tau)|^2\right)
                                      \frac{d\tau}{2\pi}.
\]

\paragraph{The estimates already established.}
The preceding sections give the following fixed-character estimates.
\begin{enumerate}
\item Proposition~\ref{prop:hankel-high-cutoff} gives selected rank
$O(\Lambda)$, and Proposition~\ref{prop:hecke-crude-gram} gives
arithmetic image Gram inverse $O(\Lambda^2)$.
\item Theorem~\ref{thm:hecke-deep-leakage} bounds the normalized theta
column outside $[-L,L]$ in Hilbert--Schmidt norm by $O(\Lambda^{-17/8})$.
\item Every fixed finite packet of nontrivial zero jets is asymptotically
annihilated by $Q_\Lambda$. More precisely, the consequence used here is
$\|Q_\Lambda P_{F,L}\|=O_{F,N}(L^{-N})$ for every fixed packet $F$
and every $N$, where $P_{F,L}$ projects onto its window jets. At a zero
$\rho=1/2+\alpha+i\gamma$, the jet generators are
\[
 \begin{cases}
 1_{[-L,L]}e^{(\alpha-i\gamma)(t-L)}(t-L)^j,&\alpha>0,\\
 1_{[-L,L]}e^{(\alpha-i\gamma)(t+L)}(t+L)^j,&\alpha<0,\\
 (2L)^{-1/2}1_{[-L,L]}e^{-i\gamma t}
                  ((t+L)/(2L))^j,&\alpha=0,
 \end{cases}
 \qquad 0\le j<m_\rho.
\]
Proposition~\ref{prop:hecke-deep-zero-jets} gives the squared
Hilbert--Schmidt bound $O(\Lambda^{4/3-56}(1+L)^C)$. Taking its
square root gives the assertion above, since every fixed negative
power of $\Lambda$ dominates every negative power of $L$.
\item For every fixed smooth strict-margin cutoff $p$, equal to one near
zero and supported in $[0,b]$ for some $b<1$, put
\[
 M_\Lambda=p(u/\Lambda),\qquad
 B_\Lambda^p=\mathcal H_cM_\Lambda\mathcal H_c,
 \qquad K_\Lambda^p=M_\Lambda B_\Lambda^pM_\Lambda.
\]
Proposition~\ref{prop:hankel-high-interior} gives, for every fixed
$N$, uniformly as an operator on the raw Hilbert space,
\begin{equation}\label{eq:hi-strict-margin-input}
 \|(1-\Pi_\Lambda)K_\Lambda^p\|_{L^2\to H_D^{3/8}}
                                      \le C_{p,N}\Lambda^{-N}.
\end{equation}
The two principal constraints can be restored exactly with the same
arbitrary-power graph accuracy. A concrete normalization of the correction
profiles and the elementary coordinate estimates are given below.
\item The Hankel transform preserves every $H_D^s$. On the critical line,
$|L_K(1/2+i\tau,\chi)|\ll(2+|\tau|)^{3/8}$, so
\begin{equation}\label{eq:hi-theta-graph}
       \|E_\chi f\|_2\le C_\chi\|f\|_{H_D^{3/8}}.
\end{equation}
The fixed-strip polynomial growth estimate, the Euler-product lower bound
at a fixed point to the right of one, and the zero count
$O_\chi(\log(2+|T|))$ per unit height are available.
\end{enumerate}
Only the strict-margin estimate, the fixed-packet annihilation, and the
analytic inputs are needed for the upper bounds. The rank, Gram, and
outside-window bounds explain their relation to the complete cutoff
construction and provide a trace-transfer corollary at the end.

\subsection{Support, scaling, and the principal constants}

For a nonprincipal character the raw time support gives
$E_\chi f(t)=0$ for $t>\log\Lambda$. The Poisson identity gives the
opposite support gap for a Hankel-band vector:
\[
 f\in\operatorname{Ran}B_\Lambda
       \quad\Longrightarrow\quad E_\chi f(t)=0\quad(t<-L).
\]
Indeed $\mathcal H_cf$ is time supported in $[0,\Lambda]$ and
$E_\chi f=\varepsilon_\chi J E_{\bar\chi}\mathcal H_cf$.
There is no conductor translation: the factor $c$ has already been
incorporated into the self-dual Hankel transform.

In the principal case let $r_K=\operatorname{Res}_{s=1}\zeta_K(s)$.
The theta multiplier normalization fixes the subtraction as
\[
 E_1f(t)=e^{t/2}\sum_{\mathfrak a\ne0}
                  f(N\mathfrak a\,e^t)-r_KIf\,e^{-t/2}.
\]
Here the ideal sum includes the unit ideal. The Euler factors give
$\zeta_K(s)=\zeta(s)L(s,\chi_{-3})$, and
\[
 L(1,\chi_{-3})
 =\int_0^1\frac{1-x}{1-x^3}\,dx
 =\int_0^1\frac{dx}{1+x+x^2}
 =\frac{\pi}{3\sqrt3}.
\]
Thus, for the primitive principal character, $c=2\pi/\sqrt3$ and
\begin{equation}\label{eq:hi-principal-residue}
                         r_K=\frac c6.
\end{equation}
The kernel at zero and involutivity give, on the smooth cores in use,
\begin{equation}\label{eq:hi-hankel-coordinates}
        (\mathcal H_cf)(0)=cIf,
        \qquad I(\mathcal H_cf)=\frac{f(0)}c.
\end{equation}
The principal root number is one. Consequently the two exact polar tails
are
\begin{align}
 E_1f(t)&=-r_KIf\,e^{-t/2} &&(t>L, P_\Lambda f=f),
 \label{eq:hi-upper-polar}\\
 E_1f(t)&=-\frac{r_K}{c}f(0)e^{t/2}
          =-\frac{f(0)}6e^{t/2}
                         &&(t<-L, B_\Lambda f=f).
 \label{eq:hi-lower-polar}
\end{align}
These support and polar identities extend to the graph domains by smooth
approximation preserving the indicated support and continuous coordinate.

The transform reverses raw scaling:
\[
 \mathcal H_cD_r=D_{1/r}\mathcal H_c,
 \qquad E_\chi D_r=\tau_{\log r}E_\chi,
 \quad (\tau_xg)(t)=g(t-x).
\]
Choose a fixed $\beta\in C_c^\infty((1/4,1/2))$ with $I\beta=1$,
and define the principal correction profiles
\begin{equation}\label{eq:hi-principal-profiles}
 b_{\Lambda,+}=D_\Lambda\beta,
 \qquad b_{\Lambda,-}=c^{-1}\mathcal H_cD_\Lambda\beta.
\end{equation}
For the normalized coordinates $I/\sqrt\Lambda$ and
$\operatorname{ev}_0/\sqrt\Lambda$, their coordinate matrix is exactly
the identity. Their theta images are fixed profiles after the appropriate
endpoint translation:
\[
 E_1b_{\Lambda,+}=\tau_LE_1\beta,
 \qquad E_1b_{\Lambda,-}=c^{-1}\tau_{-L}J E_1\beta.
\]
The factor $c^{-1}$ replaces the factor $1/2$ of the one-dimensional
even Fourier normalization.

\subsection{The Hankel cross-boundary estimate}

The only new localization calculation needed for the upper argument is
the replacement of ordinary Fourier tail estimates by Hankel tail
estimates. Set
\[
                  \mathcal D=-\partial_u(u\partial_u).
\]
The Bessel equation and integration by parts give
\begin{equation}\label{eq:hi-radial-diagonalization}
 \mathcal H_c\mathcal D f(v)=c^2v\mathcal H_cf(v),
 \qquad
 U\mathcal DU^{-1}=-e^{-t}(\partial_t-1/2)^2.
\end{equation}
Initially one may take $f\in C_c^\infty((0,\infty))$. The identity
extends whenever its displayed graph norms are finite, by truncation.
For the cutoff profiles flat near raw zero the boundary terms there also
vanish, so the same identity applies to them.

\begin{lemma}\label{lem:hi-hankel-cross}
Suppose a raw column operator $V$ has logarithmic output supported in
$t\ge-A$, and all fixed logarithmic output derivatives have finite
Hilbert--Schmidt norm. For every fixed $j,N$, its Hankel tail beyond a
spectral radius $Y$ satisfies
\[
 \|(1-p(v/Y))\mathcal H_cV\|_{2,H_D^j}
 \le C_{j,N,A}Y^{-N}
       \sum_{r\le j+2N}\|\partial_t^rUV\|_2.
\]
An enlargement of the finite derivative budget is harmless if $j$ is
not an integer. After translating a raw upper endpoint of radius
$\Lambda$ to radius one, the opposite cutoff has $Y=\Lambda^2$.
In particular, for the endpoint columns constructed below,
\begin{equation}\label{eq:hi-cross-errors}
 \|(1-B_\Lambda^p)(1-M_\Lambda)T\|_{2,H_D^{3/8}}
 +\|(1-M_\Lambda)(1-B_\Lambda^p)T\|_{2,H_D^{3/8}}
                         \le C_N\Lambda^{-N}.
\end{equation}
At height $H$, the right side becomes $H^{C_N}\Lambda^{-N}$ when
the endpoint derivative budgets are bounded by fixed powers of $H$.
\end{lemma}
\begin{proof}
For integer output derivative order, commute logarithmic differentiation
through the smooth spectral cutoff and use the reversal of the dilation
generator by $\mathcal H_c$. On the support of the spectral tail,
$v\gg Y$. Insert $v^{-N}$ and use
\eqref{eq:hi-radial-diagonalization} to transfer $v^N$ to
$c^{-2N}\mathcal D^N$ on the raw side. In logarithmic coordinates
$\mathcal D^N$ is $e^{-Nt}$ times a fixed differential polynomial of
order $2N$. The factors $e^{-Nt}$ are bounded on $t\ge-A$.
This proves the asserted estimate, column by column and then in
Hilbert--Schmidt norm. An integer derivative bound controls the required
fractional one.

By $\mathcal H_cD_\Lambda=D_{1/\Lambda}\mathcal H_c$, a cutoff at
the physical dual radius $\Lambda$ becomes a cutoff at
$Y=\Lambda^2$ after raw upper-endpoint scaling. This proves the first
term of \eqref{eq:hi-cross-errors}. Apply $\mathcal H_c$ to the second
term and use
\[
 \mathcal H_c(1-M_\Lambda)(1-B_\Lambda^p)
       =(1-B_\Lambda^p)(1-M_\Lambda)\mathcal H_c
\]
to reduce it to the same estimate for the Fourier-dual endpoint column.
\end{proof}

For later use, every fixed smooth raw profile supported strictly inside
$(0,1)$, dilated by $D_\Lambda$, is approximated in theta graph norm
by the selected image with arbitrary powers of $\Lambda^{-1}$.
Choose $p$ equal to one on its support. The preceding tail estimate
shows that $K_\Lambda^pD_\Lambda\phi-D_\Lambda\phi$ is arbitrarily
small in $H_D^{3/8}$; now use \eqref{eq:hi-strict-margin-input} and
\eqref{eq:hi-theta-graph}. The dual assertion follows by
$\mathcal H_c$ and the Poisson identity. It is enough to approximate
$\Pi_\Lambda K_\Lambda^pD_\Lambda\phi$; no estimate for
$\Pi_\Lambda$ on arbitrary fractional Sobolev vectors is required.

\paragraph{Restoring the principal coordinates.}
The profiles in \eqref{eq:hi-principal-profiles} admit the same graph
approximations. Their normalized coordinate matrix after selection tends
to the identity with arbitrary-power accuracy. Here point evaluation is
not inferred from $H_D^{3/8}$. For a selected vector $\psi$, its band
extension $w=B_\Lambda(P_\Lambda B_\Lambda P_\Lambda)^{-1}\psi$
satisfies $P_\Lambda w=\psi$ and $\|w\|\ll\|\psi\|$.
From the Bessel kernel and $|J_1(x)|\le x/2$,
\[
 |w(0)|\ll\Lambda^{1/2}\|\psi\|,
 \qquad \|w'\|_{L^2(0,1)}\ll\Lambda^{3/2}\|\psi\|.
\]
The smooth-cutoff columns have the same coarse derivative bound.
For their raw projection error $F$, use
\[
 |F(0)|^2\le C\|F\|_{L^2(0,1)}
                (\|F'\|_{L^2(0,1)}+\|F\|_{L^2(0,1)}).
\]
Cauchy--Schwarz over columns gives arbitrary-power smallness of the
normalized evaluation row. The mean row follows by Cauchy--Schwarz on
$[0,\Lambda]$; the exterior tail of the fixed dual profile is arbitrarily
small by \eqref{eq:hi-radial-diagonalization}. Subtracting the selected
profiles times the inverse of their two-by-two coordinate matrix restores
both constraints exactly. All constants depend on fixed derivative
budgets, not on the number of columns. This also records the uniformity
needed when the column norm grows with frequency height.

\subsection{Local equality}

For a nontrivial zero $\rho=\beta+i\gamma$ put
$\alpha=\beta-1/2$ and let $m_\rho$ be its multiplicity. Define
\begin{equation}\label{eq:hi-zero-measure}
 d\mathfrak b_\chi(\tau)=
 \sum_{\alpha=0}m_\rho\delta_\gamma(d\tau)
 +\sum_{\alpha\ne0}m_\rho
       \frac{|\alpha|\,d\tau}{\pi((\tau-\gamma)^2+\alpha^2)}.
\end{equation}
The unit-height zero count and $|\alpha|<1/2$ show that this is a
positive locally finite measure, with mass $O_\chi(\log(3+|T|))$ on
each unit interval. For distant ordinates, sum the inverse-square Cauchy
tails against the zero count. In particular it is integrable against
every Schwartz function.

\begin{lemma}[The one-sided Hecke closure]
\label{lem:hecke-halfline-closure}
The closure of the theta images of smooth raw profiles compactly
supported in $(0,1)$ is the inner multiple of the negative-half-line
Hardy space whose inner factor is the Blaschke product of the zeros
of $L_K(s,\chi)$ in $\Re s>1/2$, with their multiplicities.
For the principal character impose $I\psi=\psi(0)=0$.
The reflected Hankel-dual closure at the other endpoint has model
space given by the zeros in $\Re s<1/2$.
\end{lemma}
\begin{proof}
The theta--Poisson identity makes every smooth core image square
integrable, and raw support makes its logarithmic support nonpositive.
Raw contractions preserve the core and become left translations.
The closed image is therefore a nonzero translation-invariant subspace
of the half-line Hardy space. The invariant-subspace theorem writes it
as an inner multiple of that space, as in
Lemma~\ref{lem:local-halfline-closure}.

For a nonprincipal character the Mellin transforms of the generators
are $L_K(s,\chi)\mathcal M\phi(s)$. For the principal character take
$\psi=(u\partial_u+1)\phi$ with $\phi\in C_c^\infty((0,1))$.
Then $I\psi=\psi(0)=0$, and the transform is
$(1-s)\zeta_K(s)\mathcal M\phi(s)$. The factor at $s=1$ is regular
and has nonzero value $-r_K$. At any other prescribed point, a
sufficiently narrow bump has nonzero Mellin transform. Thus the common
zeros in the Hardy half-plane are exactly the nontrivial zeros, with
their orders.

Every generator continues analytically across each finite boundary
point. A common singular inner measure on that boundary is consequently
excluded: on a small interval where one generator is nonzero its inner
factor has no singular measure, while at its isolated boundary zeros
analytic finite-order vanishing also excludes a singular atom.
A remaining exponential inner factor would make every image vanish on
a fixed interval immediately below logarithmic zero. This is impossible.
Choose a bump supported arbitrarily close to raw $u=1$; for $e^t>1/2$
the theta sum has only the unit-ideal term. The same holds for the
principal differentiated bump, whose subtracted mean is zero.
There is no common delay, and the inner factor is precisely the stated
Blaschke product. Finally the functional equation and
$E_\chi\mathcal H_c=\varepsilon_\chi J E_{\bar\chi}$ identify the
dual right-half-plane divisor with the original left-half-plane divisor.
\end{proof}

\begin{theorem}\label{thm:hi-local-limit}
For the Hankel cutoff with the fixed exponent $M=56$, for every
$a\in C_c^\infty(\mathbb R)$,
\[
 \lim_{\Lambda\to\infty}\operatorname{Tr}(D_LC_aC_a^*)
                   =\int |a(\tau)|^2d\mathfrak b_\chi(\tau).
\]
The conclusion holds in the constrained principal sector, with the
constants in \eqref{eq:hi-principal-residue}--\eqref{eq:hi-lower-polar}.
No hypothesis on zero location, simplicity, or separation is imposed.
\end{theorem}
\begin{proof}
Write $F_\chi(\tau)=L_K(1/2+i\tau,\chi)$. Choose a fixed compact
interval containing $\operatorname{supp}a$ in its interior, with endpoints
away from critical zero ordinates. Let $P_{0,L}$ project onto the span of
the window jets of all critical-line zeros in that interval, and put
$W_L=S_L-P_{0,L}$.

Give every real pole of $b=a/F_\chi$ the negative-time prescription
$(\tau-\gamma-i0)^{-j}$. Its inverse transform $\kappa$ is smooth;
at positive infinity it and its derivatives are Schwartz, and at negative
infinity it is a finite exponential polynomial plus a Schwartz remainder.
This follows by subtracting the finitely many compactly cut off principal
parts: the inverse transform of a prescribed pole is a constant times
$1_{t<0}t^{j-1}e^{-i\gamma t}$.

Set $T_L=U^{-1}C_bW_L$. The input jet conditions annihilate every whole
exponential polynomial in its far negative output tail. The remaining
two exterior quadrant kernels are Hilbert--Schmidt, while the middle
kernel has polynomial growth in $L$. Thus for every fixed $j$,
\begin{equation}\label{eq:hi-fixed-bulk}
 \|UT_L\|_{\mathcal S_2(L^2,H^j)}\le C_j(1+L)^{d_j},
 \qquad E_\chi T_L=C_aW_L.
\end{equation}
The latter identity is exact in the closed multiplier realization.

The Mellin action of the Hankel transform is
\[
 \mathcal M(\mathcal H_cf)(s)
 =g_c(s)\mathcal Mf(1-s),
 \qquad g_c(s)=c^{1-2s}\frac{\Gamma(s)}{\Gamma(1-s)}.
\]
On the critical line the multiplier is smooth and unimodular.
Reflection changes the real-pole prescriptions; changing them back adds
only delta derivatives, which vanish on the reflected jet constraints.
Hence the dual column has the same kernel descriptions and polynomial
bounds. Equivalently one may use the exact Poisson identity and the dual
reciprocal multiplier.

After translating the upper endpoint to zero, the raw upper error has,
before its finite input-jet correction, kernel
\[
                 (1-p(e^t))\kappa(t-s)1_{[-2L,0]}(s).
\]
The output is bounded below by a fixed number and the input is
nonpositive. The positive Schwartz tail of $\kappa$ therefore gives a
fixed Hilbert--Schmidt limit, including every fixed logarithmic derivative
and output polynomial weight. The translated jet projections tend strongly
to zero. To check this directly, scale their normalized polynomial
amplitudes by $L^{-1/2}$; their fixed Gram matrices converge to nonsingular
Hilbert matrices, while distinct ordinates decouple by integration by
parts. Composition with the fixed Hilbert--Schmidt quadrant kernel then
removes their correction in Hilbert--Schmidt norm. The dual endpoint has
the identical property.

Lemma~\ref{lem:hi-hankel-cross} supplies the two cross errors with
arbitrary powers of $\Lambda^{-1}$. In a nonprincipal sector set
$A=M_\Lambda$, $B=B_\Lambda^p$, and $P_+=1_{t\le L}$. The exact
identity is
\begin{align*}
 S_LE_\chi T_L-E_\chi ABA T_L
   ={}&S_LE_\chi(1-A)T_L
      +S_LE_\chi(1-B)AT_L\\
     &+P_+E_\chi(1-A)BAT_L.
\end{align*}
Indeed expansion gives an extra $(P_+-S_L)E_\chi BAT_L$, which vanishes
by the band support gap, and the final time-supported image has no
component above $L$. The cross estimates replace $AT_L$ by $T_L$ in
the second term and $BAT_L$ by $AT_L$ in the last, with vanishing
Hilbert--Schmidt error after applying \eqref{eq:hi-theta-graph}.
The residual therefore has two fixed translated Hilbert--Schmidt limits.
Include also $S_LC_a(1-S_L)$, whose two outside-input quadrants have
fixed Hilbert--Schmidt limits because $\check a$ is Schwartz.

For the principal sector put $g=ABA T_L$,
$r_+=Ig/\sqrt\Lambda$, $r_-=g(0)/\sqrt\Lambda$, and
\[
 g^0=g-b_{\Lambda,+}r_+-b_{\Lambda,-}r_-.
\]
Its two coordinates vanish. Expansion of the preceding identity now has
the following exact form. If $R_0$ denotes its three-term right side,
\[
 R_0=S_LE_1T_L+(P_+-S_L)E_1BAT_L-P_+E_1g,
\]
so that
\[
 S_LE_1T_L-E_1g
 =R_0+r_KIg\,e^{-t/2}1_{t>L}
       +(r_K/c)g(0)e^{t/2}1_{t<-L}.
\]
The value of $BAT_L$ at zero equals that of $g$ because $A$ is
one there; the band-side polar tail is consequently expressed by
the same $g(0)$ as the constraint row.
Adding back the two correction profiles cancels these tails exactly,
leaving only
\begin{equation}\label{eq:hi-principal-boundaries}
 1_{t\le L}E_1b_{\Lambda,+}r_+
       +1_{t\ge-L}E_1b_{\Lambda,-}r_-.
\end{equation}

The two rows have fixed translated $L^2$ limits. For the mean row, move
$B-1$ to the normalized smooth integration test; its Hankel tail is
arbitrarily small. The main row is
$\Lambda^{-1/2}I(A^2T_L)$, whose translated kernel is the integral of
$\kappa(t-s)$ against $e^{t/2}p(e^t)^2$. For the value row use
$(\mathcal H_cf)(0)=cIf$ and move the first time cutoff to the smooth
Hankel integration test. The main dual row has weight
$c e^{t/2}p(e^t)$. These weights have exponential negative tails and
bounded upper support, so their row kernels have all fixed integrable
moments on $s\le0$. Their finite jet corrections vanish as above.
The output profiles in \eqref{eq:hi-principal-boundaries} are fixed after
translation and belong to $L^2$. Thus they are legitimate fixed
Hilbert--Schmidt boundary terms.

In both sectors we have constructed
\begin{equation}\label{eq:hi-endpoint-decomposition}
 S_LC_a(1-P_{0,L})
   =F_L+\tau_LB_+\tau_{-L}
          +\tau_{-L}B_-\tau_L+o_{\mathcal S_2}(1),
\end{equation}
where $B_+$ has range in $L^2(( -\infty,0])$, $B_-$ has range in
$L^2([0,\infty))$, and $\|F_L\|_2\le C(1+L)^d$.
The bulk is $E_\chi ABA T_L$ or $E_1g^0$, respectively.
The strict-margin input and exact constraint restoration give
\begin{equation}\label{eq:hi-bulk-distance}
                         \|(1-Q_\Lambda)F_L\|_2\longrightarrow0.
\end{equation}
This is a direct graph approximation by columns in the theta image;
it does not infer image-projection closeness from raw-projection closeness.

Let $M_+\subset L^2(( -\infty,0])$ be the one-sided theta closure and
$M_-\subset L^2([0,\infty))$ its dual reflected version.
Lemma~\ref{lem:hecke-halfline-closure} identifies their complements
$K_\pm$ as the pure Blaschke model spaces for the right and left
half-strip zeros of the original character.

The fixed-profile graph approximation proves
$\|(1-Q_\Lambda)\tau_{\pm L}v\|\to0$ for each $v\in M_\pm$.
Finite exponential-jet packets are dense in $K_\pm$. Their annihilation rate, and the exponentially small truncation of a fixed
off-line packet beyond the opposite endpoint, give for every fixed
packet projection $P_{K_{\pm,n}}$ and every $N$,
\[
 \|Q_\Lambda\tau_{\pm L}P_{K_{\pm,n}}\|
                                      =O_{n,N}(L^{-N}).
\]
Here $4/3-56=-164/3<0$ is more than sufficient; no height-uniform
packet Gram estimate is required.

For clarity, the endpoint Hilbert-space conclusion can be verified
directly from \eqref{eq:hi-endpoint-decomposition}. Strong convergence
on $M_\pm\oplus K_\pm$, followed by composition with the fixed
Hilbert--Schmidt boundaries, gives
\[
 (1-Q_\Lambda)S_LC_a(1-P_{0,L})
 =\tau_LP_{K_+}B_+\tau_{-L}
  +\tau_{-L}P_{K_-}B_-\tau_L+o_{\mathcal S_2}(1).
\]
Project the original decomposition onto a fixed endpoint packet.
Its bulk tends to zero since
\[
 \|P_{K_{+,n}}\tau_{-L}F_L\tau_L\|_2
 \le \|Q_\Lambda\tau_LP_{K_{+,n}}\|\,\|F_L\|_2
           +\|(1-Q_\Lambda)F_L\|_2\longrightarrow0.
\]
The opposite boundary tends to zero by compactness and weak convergence
of translations. The translated input projection tends strongly to zero.
Thus $P_{K_{+,n}}B_+=P_{K_{+,n}}C_a$, and exhaustion gives
$P_{K_+}B_+=P_{K_+}C_a$; the lower endpoint is identical.
The two translated boundaries are asymptotically Hilbert--Schmidt
orthogonal, by truncating their output ranges to fixed bounded intervals.

Lemma~\ref{lem:critical-window-packet} gives the critical contribution
$\sum_{\Re\rho=1/2}m_\rho|a(\Im\rho)|^2$ from the fixed-packet
annihilation. Its inverse-Gram calculation also accounts for repeated
critical zeros without commuting the multiplier and packet projection.
The critical input and its orthogonal complement add exactly in squared
Hilbert--Schmidt norm. Finally a Malmquist basis of each finite Blaschke
packet has squared Fourier diagonal equal to the sum
$2|\alpha|/((\tau-\gamma)^2+\alpha^2)$, including repeated zeros.
Increasing packets and the zero count give
\[
 \|P_{K_+}C_a\|_2^2+\|P_{K_-}C_a\|_2^2
 =\sum_{\alpha\ne0}m_\rho
       \int |a(\tau)|^2
             \frac{|\alpha|d\tau}{\pi((\tau-\gamma)^2+\alpha^2)}.
\]
The identity
$\|(1-Q_\Lambda)S_LC_a\|_2^2=\operatorname{Tr}(D_LC_aC_a^*)$
completes the proof.
\end{proof}

Polarization gives convergence on smooth compact frequency tests. A
smooth squared test equal to one on a fixed compact set bounds its masses;
uniform approximation then gives vague convergence of $\mu_\Lambda$ to
$\mathfrak b_\chi$. No noncompact test has yet been used.

\subsection{A reciprocal parametrix uniform in zero clusters}

Put $F_\chi(z)=L_K(1/2+iz,\chi)$ and $H=2+|T|$. A zero
$\rho=\beta+i\gamma$ corresponds to
$z_\rho=\gamma-i(\beta-1/2)$, so all nontrivial zeros lie in
$|\Im z|<1/2$.

\begin{lemma}\label{lem:hi-rectangle}
For sufficiently large $|T|$, there are a real bump
$a_T\in C_c^\infty((T-2,T+2))$, with $0\le a_T\le1$ and equal
to one on $[T-1/2,T+1/2]$, and a packet of $N\le C\log H$ zeros
$z_1,\ldots,z_N$, counted with multiplicity, with these properties.
Every zero projected into the support of $a_T$ belongs to the packet,
and $a_T$ is one near each such projection. If $R_{T,j}$ is the proper
rational sum of the principal parts of $z^j/F_\chi(z)$ at the packet
zeros, with negative-time prescriptions at real poles, then
\[
 r_{T,j}(\tau)=\frac{\tau^j a_T(\tau)}{F_\chi(\tau)}-R_{T,j}(\tau)
\]
is an ordinary smooth function and, for every fixed $j,m$,
\begin{equation}\label{eq:hi-remainder}
                  \|r_{T,j}^{(m)}\|_2\le H^{C_{j,m}}.
\end{equation}
If $R_{T,j,\pm}(t)$ are the inverse-transform exponential polynomials
on the positive and negative half-lines, extended as whole exponential
polynomials, then for each fixed $A$,
\begin{equation}\label{eq:hi-residue-difference}
 \sup_{-A\le t\le0}|R_{T,j,-}(t)-R_{T,j,+}(t)|
                                         \le H^{C_{j,A}}.
\end{equation}
The same assertions hold with an analytic numerator polynomially bounded
on a fixed neighborhood of the rectangle used in the proof. Constants
do not depend on distances between zeros or their distances to the axis.
\end{lemma}
\begin{proof}
Use a large fixed disk centered at $2+iT$ in the $s$-plane. The
fixed-strip bound is polynomial in $H$, the center value is bounded
below by the Euler product, and the number of zeros in a smaller disk
is $O(\log H)$. Divide by the disk Blaschke product of those zeros.
The quotient is zero free. The maximum principle bounds it above;
Harnack's inequality applied to a constant minus its logarithmic modulus
bounds its reciprocal on a fixed inner disk containing the required
rectangle in the $z$-coordinate. Fixed radii eight, between six and seven,
and five suffice. Large $|T|$ excludes trivial zeros and the principal
pole from these disks.

Let $\gamma_k$ be the real projections of the zeros in the Blaschke
disk, with multiplicity. On each of
$[T-3/2,T-1]$ and $[T+1,T+3/2]$, the integral of
$\log\min(1,|x-\gamma_k|)$ is bounded below uniformly. Summing and
averaging therefore gives $x_-,x_+$ with
\[
       \prod_k\min(1,|x_\pm-\gamma_k|)\ge e^{-CN_0}.
\]
Every factor is at most one, so each individual distance is at least
$e^{-CN_0}$. Choose a width $\varepsilon\ge H^{-C}$ small enough
that comparable product bounds hold within $4\varepsilon$ of the two
chosen lines. The vertical strips are zero free and have polynomial
reciprocal bounds. The horizontal sides at $\Im z=\pm3/4$ stay a
fixed distance from all nontrivial zeros. The resulting rectangle
$\Gamma$ has polynomial reciprocal bounds on its contour and suitable
side neighborhoods.

Take all zeros inside $\Gamma$ as the packet. Choose $a_T$ equal to
one between $x_-+\varepsilon$ and $x_+-\varepsilon$, and zero outside
$[x_--\varepsilon,x_++\varepsilon]$, with transitions in the zero-free
strips. Its fixed derivatives cost at most fixed powers of $H$.

Inside $\Gamma$, subtracting the principal parts leaves the Cauchy
integral of $z^j/F_\chi(z)$ on the contour. It and its fixed derivatives
are polynomially bounded away from a side. Within a transition strip,
move the vertical side to the other side of that strip without crossing
a zero. Cauchy's estimate, with distance comparable to $\varepsilon$,
gives the same bounds. This controls the real remainder and cancels every
real distributional pole, since the cutoff is one there.

For the tails, let $P(z)=\prod_{k=1}^N(z-z_k)$. Multiplying the contour
formula for the combined principal parts by $P$ gives a numerator of
degree less than $N$, with coefficients in powers of $z-T$ bounded by
$H^{C_j}C^N$. This follows by polynomial division on a fixed-size contour;
it involves no division by a root gap. Since $N=O(\log H)$, this is a
polynomial bound. Outside $|\tau-T|\le3$ the roots lie in a fixed
smaller rectangle, so
\[
 |R_{T,j}^{(m)}(\tau)|
       \le H^{C_{j,m}}(1+|\tau-T|)^{-m-1}.
\]
The remainder there is $-R_{T,j}$. This proves
\eqref{eq:hi-remainder}.

Closing the inverse-transform contour in the appropriate half-plane
shows
\[
 R_{T,j,-}(t)-R_{T,j,+}(t)
    =\frac1{2\pi}\int_\Gamma
                    \frac{z^j e^{-izt}}{F_\chi(z)}\,dz.
\]
On the fixed interval $-A\le t\le0$ the exponential is bounded by
$e^{3A/4}$, and the contour length is bounded. This proves
\eqref{eq:hi-residue-difference}. The analytic-numerator version uses
the same contour estimates.
\end{proof}

\begin{lemma}\label{lem:hi-allpole-boundaries}
Let $W_{L,T}$ be the orthogonal projection in $S_LL^2$ onto functions
whose entire Fourier transforms vanish at every packet zero $z_k$ to
its indicated multiplicity. Thus it removes the span of
$1_{[-L,L]}t^j e^{-i\overline z_k t}$. Put
\[
 b_T=a_T/F_\chi,
 \qquad T_{L,T}=U^{-1}C_{b_T}W_{L,T}.
\]
For every fixed $j$,
\begin{equation}\label{eq:hi-allpole-bulk}
 \|UT_{L,T}\|_{\mathcal S_2(L^2,H^j)}
             \le H^{C_j}(1+L)^{N+C_j},\qquad L\ge1.
\end{equation}
For a fixed smooth multiplier $m(t)$ supported in $t\ge-A$, with
bounded derivatives, the upper-translated boundary
$m(t)\tau_{-L}UT_{L,T}\tau_L$ has Hilbert--Schmidt norm at most
$H^C$, uniformly in $L$, after any fixed output derivative and polynomial
output weight. The dual endpoint has the same property.
\end{lemma}
\begin{proof}
The combined proper rational part admits a Newton expansion
\[
 R_{T,0}(z)=\sum_{j=1}^Nd_j\prod_{k=1}^j(z-z_k)^{-1},
 \qquad |d_j|\le H^C.
\]
The coefficient bound follows from the polynomial numerator in the
preceding proof, expressed in $z-T$; divided differences of a polynomial
use only bounded shifted roots and factors $C^N$, not root-gap
denominators.

On an input supported in $[-L,L]$ and with zero moment at $z_k$, divide
by $z-z_k$ using the decaying first-order exponential integration kernel.
Its modulus is at most one. The zero moment forces the result to vanish
beyond the opposite endpoint, so it remains supported in $[-L,L]$.
Its operator and Hilbert--Schmidt norms on that interval are at most
$2L$. Division removes one zero and preserves all remaining conditions,
including multiplicities. A product of $j$ inverse factors therefore
costs at most $(2L)^j$ in Hilbert--Schmidt norm. This also handles real
poles with the prescribed boundary value. The smooth remainder contributes
at most $(2L/(2\pi))^{1/2}\|r_{T,0}\|_2$. This proves
\eqref{eq:hi-allpole-bulk}; fixed derivatives cost only fixed powers of
$H$ because the output frequency support is $[T-2,T+2]$.

After upper-endpoint translation the input lies in $[-2L,0]$ and retains
all its exponential-jet moment conditions. For each output derivative,
use the decomposition with $R_{T,j}$ and $r_{T,j}$. The whole
exponential polynomial $R_{T,j,+}(t-s)$ annihilates the input projection
exactly. Subtract it first. The rational contribution left over is
\[
 1_{t-s<0}\bigl(R_{T,j,-}(t-s)-R_{T,j,+}(t-s)\bigr).
\]
On $t\ge-A$, $s\le0$, its support is the bounded triangle
$-A\le t\le s\le0$. Equation~\eqref{eq:hi-residue-difference}
therefore gives a polynomial Hilbert--Schmidt bound with every fixed
output weight. This is a bound on the combined residue polynomial,
even if individual residues become large.

For the smooth remainder, Plancherel and \eqref{eq:hi-remainder} give
all fixed weighted time norms. In particular
\[
 \int_{t\ge-A,\ s\le0}(1+|t|)^{2J}
              |\check r_{T,j}(t-s)|^2\,dt\,ds
 \le C_{J,A}\int_{u\ge-A}(1+|u|)^{2J+1}
              |\check r_{T,j}(u)|^2du\le H^{C_{J,j,A}}.
\]
Input restriction and right multiplication by an orthogonal projection
cannot increase the norm. Leibniz's rule handles $m$. Using the principal
part of $z^j/F_\chi$ at each derivative stage also includes correctly
any time-zero distribution produced by differentiation.

The exact Poisson identity identifies the dual inverse with the reflected
dual-character reciprocal. More explicitly,
$L_K(s,\bar\chi)g_c(s)=\varepsilon_{\bar\chi}L_K(1-s,\chi)$,
so
\[
        \frac{g_c(s)}{L_K(1-s,\chi)}
                    =\frac{\varepsilon_{\bar\chi}}{L_K(s,\bar\chi)}.
\]
Reflection sends the input zero conditions to exactly that dual packet.
Thus the dual boundary has the same estimates, with no gamma growth or
conductor displacement left over. Lemma~\ref{lem:hi-hankel-cross}
then gives the required ordinary radial derivative budgets as well.
\end{proof}

\subsection{Uniform local mass and Schwartz convergence}

\begin{theorem}\label{thm:hi-uniform-mass}
For the fixed exponent $M=56$ and the stated Hankel cutoff, there are
constants $C,A$ depending on the fixed character such that
\[
              \mu_\Lambda([T,T+1])\le C(2+|T|)^A
                     \qquad(\Lambda\ge2,\ T\in\mathbb R).
\]
The conclusion includes the constrained principal sector and requires no
zero-location or zero-separation assumption. Consequently
$\mu_\Lambda\to\mathfrak b_\chi$ against every Schwartz frequency
test.
\end{theorem}
\begin{proof}
Use the preceding bump centered at $T+1/2$. Let
$P_{L,T}=S_L-W_{L,T}$, which has rank at most $N\le C\log H$.
The local construction applies to $T_{L,T}$ with these additional
off-line zero conditions: the exact identity
$E_\chi T_{L,T}=C_{a_T}W_{L,T}$ is unchanged. By
Lemma~\ref{lem:hi-allpole-boundaries}, all boundary derivative and moment
budgets are bounded by fixed powers of $H$, uniformly in $L$. Their
theta images obey the same bounds by \eqref{eq:hi-theta-graph}.
The outside-input boundary $S_LC_{a_T}(1-S_L)$ is controlled by the
fixed derivative bounds of the bump.

In the principal sector the two constraint rows have the same polynomial
bound. Here are the details that avoid estimating individual residues.
For $k=1,2$, put $\theta_k(t)=e^{t/2}p(e^t)^k$. Then
$(\partial_t-1/2)\theta_k$ is smooth and compactly supported, and
\[
 \widehat\theta_k(z)
 =-\frac{\widehat{(\partial_t-1/2)\theta_k}(z)}{1/2+iz}.
\]
This continuation has only a possible pole at $z=i/2$. On the large-height
rectangles it is analytic with polynomial bounds. The normalized mean
row, up to an arbitrary-power cross error, is the restriction to $s\le0$
of
\[
       \check{(b_T\widehat\theta_2(-\cdot))}(-s)
\]
followed by the translated input projection. The whole reflected
positive-time residue polynomial is annihilated by that projection.
The analytic-numerator form of Lemma~\ref{lem:hi-rectangle} bounds the
remaining smooth part in $L^2$ by a fixed power of $H$. The value row
is the identical dual calculation with numerator
$c\widehat\theta_1(-\cdot)$. Only the unweighted $L^2$ norm of these
input rows is needed; orthogonal projection does not preserve uniform
polynomial input moments. The two output polar profiles in
\eqref{eq:hi-principal-boundaries} have fixed translated $L^2$ norms,
so the additional boundaries are bounded by $H^C$.

The strict-margin estimate is an operator estimate, with constants
independent of the column operator. Combining it with
\eqref{eq:hi-allpole-bulk}, and retaining the column norm explicitly,
bounds the bulk graph error after a fixed derivative budget by
\begin{equation}\label{eq:hi-uniform-bulk-error}
                 H^C(1+L)^{C\log H+C}e^{-\eta L}
\end{equation}
for a fixed $\eta>0$. Exact principal constraint restoration has the
same form by the coordinate estimates above. The cross errors are bounded
by this expression or by the stronger $H^Ce^{-\eta L}$.
Here the derivative budgets are fixed once; they do not grow with the
number of packet zeros. The fixed choice $M=56$ is unchanged.

Choose a fixed $D$ so large that all these errors are at most one when
$L\ge H^D$. Indeed the logarithm of \eqref{eq:hi-uniform-bulk-error}
is at most $C\log H+(C\log H+C)\log(1+L)-\eta L$, which is negative
uniformly above this threshold for sufficiently large $H$.
No approximation of the boundary ranges by the theta image is needed
for this upper bound; contractivity gives
\[
       \|(1-Q_\Lambda)S_LC_{a_T}(1-P_{L,T})\|_2\le H^C.
\]
The complementary input has squared norm at most
$\operatorname{rank}P_{L,T}\le C\log H$, since $\|a_T\|_\infty\le1$.
The input spaces are orthogonal, so their squared Hilbert--Schmidt norms
add exactly. Since $a_T=1$ on $[T,T+1]$, the distance identity proves the
desired mass bound for $L\ge H^D$.

For $L<H^D$, Bessel's inequality gives
\[
 0\le\frac{d\mu_\Lambda}{d\tau}\le\frac L\pi,
 \qquad \mu_\Lambda([T,T+1])\le H^D/\pi.
\]
Bounded heights follow from Theorem~\ref{thm:hi-local-limit} and the
same elementary estimate on the remaining bounded cutoff range.
This proves the uniform polynomial local mass bound.

For a Schwartz test, sum its rapidly decreasing supremum over unit
intervals against that polynomial mass bound. The tails are uniformly
summable. Vague convergence on the compact part therefore extends to
convergence against the full Schwartz test.
\end{proof}

\paragraph{Relation to the uncompressed trace.}
Theorem~\ref{thm:hecke-deep-leakage} also permits passage to the
uncompressed trace. If $Z_\Lambda$ is
the orthonormal theta column, then
$\|(1-S_L)Z_\Lambda\|_2=O(\Lambda^{-17/8})$ and its rank is
$O(\Lambda)$. Consequently
\[
 \|Q_\Lambda-S_LQ_\Lambda S_L\|_1
 \le2\sqrt{\operatorname{rank}Q_\Lambda}
                 \|(1-S_L)Z_\Lambda\|_2
                      =O(\Lambda^{-13/8}).
\]
Thus every bounded convolution operator has vanishing compression trace
error. For $h\in C_c^\infty(\mathbb R)$ and
$H(z)=\int h(t)e^{zt}dt$, the exact compressed identity and the preceding
Schwartz convergence give
\[
 2\log\Lambda\,h(0)-\operatorname{Tr}(Q_\Lambda V(h))
       \longrightarrow\int H(i\tau)\,d\mathfrak b_\chi(\tau).
\]
This proves the trace limit for the constrained principal image.
The next section supplies the two polar coordinates.

\section{The Hecke trace limit and the principal polar coordinates}
\label{sec:hecke-assembly}
The local limit and the uniform frequency bound now give the trace
formula for the deep Hankel selection. This section completes its
principal sector by proving that the two missing coordinates become
vertical in the enlarged Hilbert space.

\begin{proposition}[The principal Hecke graph]
\label{prop:hecke-principal-graph}
Let $\Pi_\Lambda=\Pi_{\Lambda,56}$ for the principal character of
$K=\Q(\sqrt{-3})$. Let $Q_\Lambda$ project onto
$E_1(\operatorname{Ran}\Pi_\Lambda\cap\ker I\cap\ker\operatorname{ev}_0)$.
On $L^2(\R)\oplus\C^2$ define
\[
 J\psi=(E_1\psi,\psi(0),I\psi),\qquad
 \mathcal Q_\Lambda=P_{J\operatorname{Ran}\Pi_\Lambda}.
\]
For all sufficiently large $\Lambda$, the constraint map has rank two,
and
\begin{equation}\label{eq:hecke-vertical-graph}
 \bigl\|\mathcal Q_\Lambda-(Q_\Lambda\oplus0)
                         -P_{\C^2}\bigr\|_1
                         =O(\Lambda^{-1/4}).
\end{equation}
For $H(z)=\int h(t)e^{zt}dt$ and
$R(h)=V(h)\oplus H(-1/2)\oplus H(1/2)$, it follows that
\[
 \Tr(\mathcal Q_\Lambda R(h))-
           \Tr(Q_\Lambda V(h))
                  \longrightarrow H(-1/2)+H(1/2).
\]
No hypothesis on the zeros is used.
\end{proposition}
\begin{proof}
Choose $q,r\in C_c^\infty([0,\infty))$ with
\[
 q(0)=1,\quad Iq=0,\qquad r(0)=0,\quad Ir=1.
\]
Take $r$ to be a normalized bump away from zero, and obtain $q$
by subtracting a multiple of $r$ from a bump equal to one near zero.
The functions $q(|z|^2)$ and $r(|z|^2)$ are smooth and compactly
supported on $\R^2$; hence their radial Fourier transforms, and thus
$\mathcal H_cq$ and $\mathcal H_cr$, decrease faster than every power.
Set
\[
 a=\Lambda^{-1/2},\qquad b=\Lambda^{1/2},\qquad
 q_\Lambda(u)=q(u/a),\qquad r_\Lambda(u)=b^{-1}r(u/b).
\]
For large $\Lambda$ these functions are supported inside $(0,\Lambda)$,
with their prescribed coordinates exactly equal to $(1,0)$ and $(0,1)$.
Their raw norms are $O(\Lambda^{-1/4})$. The exact scaling identities
\[
 E_1[f(\cdot/a)]=a^{1/2}\tau_{\log a}E_1f,
 \qquad
 \mathcal H_c[f(\cdot/a)](v)=a(\mathcal H_cf)(av)
\]
show that their theta norms are also $O(\Lambda^{-1/4})$ and that,
for every $N$,
\begin{equation}\label{eq:hecke-lift-tail}
 \|(1-P_\Lambda)\mathcal H_cq_\Lambda\|_2+
 \|(1-P_\Lambda)\mathcal H_cr_\Lambda\|_2
                                  =O_N(\Lambda^{-N}).
\end{equation}
The corresponding absolute tail integrals have the same property.
Indeed the scaled lower integration limits are respectively
$a\Lambda=\Lambda^{1/2}$ and $b\Lambda=\Lambda^{3/2}$.

Put $A=P_\Lambda\mathcal H_cP_\Lambda\mathcal H_cP_\Lambda$
and $F=P_\Lambda\mathcal H_cP_\Lambda$, so $A=F^2$. For either
seed $f=P_\Lambda f$, the spectral theorem gives
\[
 \|(1-\Pi_\Lambda)f\|_2^2
 \leq\Lambda^{56}\langle(1-A)f,f\rangle
 =\Lambda^{56}\|(1-P_\Lambda)\mathcal H_cf\|_2^2.
\]
The fixed loss $\Lambda^{28}$ in norm is absorbed by
\eqref{eq:hecke-lift-tail}. Thus
$\|\Pi_\Lambda f-f\|_2=O_N(\Lambda^{-N})$ for every $N$.
The mean error is bounded by $\sqrt\Lambda$ times this norm.

Point evaluation needs a separate estimate. Hankel inversion for the
smooth seed gives
\[
 f(0)-(Af)(0)=c\int_\Lambda^\infty(\mathcal H_cf)(v)dv
                                           =O_N(\Lambda^{-N}).
\]
Since $(Fg)(0)=c\int_0^\Lambda g(u)du$ and $F$ is invertible on
the selected subspace,
\[
 |(Af)(0)-(\Pi_\Lambda f)(0)|
 \leq c\sqrt\Lambda\,
                 \|Ff-F^{-1}\Pi_\Lambda f\|_2.
\]
Commutation of $F$, $A$, and $\Pi_\Lambda$ bounds the last norm by
\[
 \|(1-\Pi_\Lambda)f\|_2+
 (1-\Lambda^{-56})^{-1/2}\|(A-1)f\|_2.
\]
Also $(1-A)f=P_\Lambda\mathcal H_c(1-P_\Lambda)\mathcal H_cf$.
Consequently the value error, as well as the mean error, is
$O_N(\Lambda^{-N})$. This uses the explicit band kernel rather than
an invalid point-evaluation estimate in fractional logarithmic Sobolev
space.

For any fixed $3/8<s<1/2$, the projected seeds have at most polynomial
$H_D^s$ norm by Proposition~\ref{prop:hankel-raw-estimates}; the
unprojected seeds have $H_D^s$ norms $O(\Lambda^{-1/4})$ by scaling.
Interpolation with the arbitrary-power raw error gives
\[
 \|\Pi_\Lambda f-f\|_{H_D^{3/8}}=O_N(\Lambda^{-N}).
\]
The graph bound \eqref{eq:hecke-graph} gives the same theta error.
The two projected coordinate columns form a matrix
$M_\Lambda=I_2+O_N(\Lambda^{-N})$. Multiplication by its inverse
produces $w_{1,\Lambda},w_{2,\Lambda}\in\operatorname{Ran}\Pi_\Lambda$
with coordinates exactly $(1,0)$ and $(0,1)$ and
\[
                  \|E_1w_{j,\Lambda}\|_2=O(\Lambda^{-1/4}).
\]
The constraint map therefore has rank two.

Since $J(\operatorname{Ran}\Pi_\Lambda\cap\ker I\cap\ker\operatorname{ev}_0)
=\operatorname{Ran}Q_\Lambda\oplus0$, the difference
$\mathcal Q_\Lambda-(Q_\Lambda\oplus0)$ is an orthogonal projection
of rank two. Its range is the graph over $\C^2$ of the operator
$X_\Lambda$ whose columns are $(1-Q_\Lambda)E_1w_{j,\Lambda}$.
Hence $\|X_\Lambda\|=O(\Lambda^{-1/4})$. The graph-projection
formula
\[
 P_{\operatorname{graph}X}=
 \begin{pmatrix}X\\I_2\end{pmatrix}
 (I_2+X^*X)^{-1}
 \begin{pmatrix}X^*&I_2\end{pmatrix}
\]
gives \eqref{eq:hecke-vertical-graph}; the difference from the vertical
projection has rank at most four and operator norm $O(\|X_\Lambda\|)$.
Finally raw unitary dilation by $e^t$ acts on the two coordinates by
$e^{-t/2}$ and $e^{t/2}$. Its theta action is $V(e^t)$, so its
integrated action is exactly the stated $R(h)$. Trace-norm convergence
against this bounded operator proves the last assertion.
\end{proof}

\subsection{Assembly for the auxiliary deep selection}
For each fixed character, Theorem~\ref{thm:hi-local-limit} identifies
the limit on compact frequency tests, and
Theorem~\ref{thm:hi-uniform-mass} extends it to every Schwartz test.
For $h\in C_c^\infty(\R)$ the function $H(i\tau)$ is Schwartz.
The exact compressed trace identity therefore gives
\[
 2Lh(0)-\Tr(S_LQ_{\chi,\Lambda}S_LV(h))
                   \longrightarrow\mathfrak B_\chi(h).
\]
The rank bound and Theorem~\ref{thm:hecke-deep-leakage} imply
\[
 \|Q_{\chi,\Lambda}-S_LQ_{\chi,\Lambda}S_L\|_1
 \leq2\sqrt{\operatorname{rank}Q_{\chi,\Lambda}}
              \|(1-S_L)Z_{\chi,\Lambda}\|_2
                         =O(\Lambda^{-13/8}).
\]
Thus the same limit holds for the uncompressed image projection.
Proposition~\ref{prop:hecke-principal-graph} proves the principal
polar correction. This completes the auxiliary deep-cutoff trace
formula. Theorem~\ref{nh:theorem} uses its upper bound and uniform
mass estimate to prove the fixed-threshold statement.

For the finite-order test algebra of Section~\ref{sec:complex-place},
only finitely many character blocks occur. Summing the sector identities
gives
\[
 \lim_{\Lambda\to\infty}
 \left(2\log\Lambda\,h(1)-\Tr(\mathcal Q_\Lambda R(h))
       +\int_\R h_1(t)(e^{-t/2}+e^{t/2})dt\right)
                 =\sum_{\chi\in F}\mathfrak B_\chi(h_\chi).
\]
Here $\mathcal Q_\Lambda$ is the direct sum of the sector projections;
its product with a finite-order test has finite rank. The limit is a
positive distribution on convolution squares. By the Hecke form of
Weil's criterion, it equals Weil's distribution on every test in these
sectors if and only if every finite-order Hecke function of $K$ has
its nontrivial zeros on the critical line. This implication uses neither
simple zeros nor a faithful representation of their full nilpotent jets.

\subsection{Comparison with a fixed concentration threshold}
\begin{proposition}[Comparison with the threshold $1/2$]
\label{prop:hecke-natural-upper}
Let $Q_\Lambda^{\mathrm{nat}}$ be the full-Gram theta-image projection
for $1_{[1/2,1]}(A_\Lambda)$, with the same principal constraints,
and let $Q_\Lambda^{\mathrm{deep}}$ denote the selection used above.
Let $\mu_\Lambda^{\mathrm{nat}}$ and $\mu_\Lambda^{\mathrm{deep}}$
be their compressed defect measures. Then
\[
 0\leq\mu_\Lambda^{\mathrm{nat}}\leq
                     \mu_\Lambda^{\mathrm{deep}}.
\]
In particular the natural measures have a uniform polynomial local-mass
bound, and for every nonnegative Schwartz function $\varphi$,
\[
 \limsup_{\Lambda\to\infty}\int\varphi\,d\mu_\Lambda^{\mathrm{nat}}
                              \leq\int\varphi\,d\mathfrak b_\chi.
\]
Writing $R_\Lambda=Q_\Lambda^{\mathrm{nat}}-Q_\Lambda^{\mathrm{deep}}$,
the natural compressed measures converge to $\mathfrak b_\chi$ on
Schwartz tests if and only if
\begin{equation}\label{eq:hecke-natural-extra-image}
 \|R_\Lambda S_LC_a\|_2\longrightarrow0
                           \qquad(a\in C_c^\infty(\R)).
\end{equation}
\end{proposition}
\begin{proof}
The deep raw subspace is contained in the natural one, also after
intersection with the two principal kernels. Their theta-image
subspaces are therefore nested, so $R_\Lambda$ is an orthogonal
projection and
\[
 0\leq S_L-S_LQ_\Lambda^{\mathrm{nat}}S_L
       \leq S_L-S_LQ_\Lambda^{\mathrm{deep}}S_L.
\]
Taking the Fourier diagonal gives the measure inequalities. The
uniform mass bound and the stated upper limit follow from the deep
theorem. Furthermore the exact trace identity gives
\[
 \int |a|^2\,d(\mu_\Lambda^{\mathrm{deep}}-
                          \mu_\Lambda^{\mathrm{nat}})
                 =\|R_\Lambda S_LC_a\|_2^2.
\]
Thus equality of the limits implies \eqref{eq:hecke-natural-extra-image}.
Conversely, that condition gives convergence on squared smooth compact
tests. Polarization gives all smooth compact tests, and the common
polynomial local-mass bound then gives every Schwartz test.
\end{proof}

This identifies the contribution of the additional theta-image
directions. Passing from the compression to the uncompressed trace
also requires
\[
 \Tr\bigl((Q_\Lambda^{\mathrm{nat}}-
                 S_LQ_\Lambda^{\mathrm{nat}}S_L)V(h)\bigr)\longrightarrow0.
\]
Neither assertion follows merely from the equality of the leading raw
ranks. The concentration calculation before arithmetic Gram normalization
does not by itself control these two image-space quantities.
Section~\ref{nh:section} proves the needed estimates by arithmetic
duality and mean-square summation.


\section{Every fixed Hankel concentration threshold}
\label{nh:section}

Fix a primitive finite-order Hecke character $\chi$ of
$K=\mathbb Q(\sqrt{-3})$, and retain the normalization
$c=2\pi/\sqrt{3N\mathfrak f_\chi}$ and the corrected theta map
$E_\chi$ from Section~\ref{sec:hecke-arithmetic-attempt}.
In positive variables, $E_\chi\psi(x)$ denotes
$(E_\chi\psi)(\log x)$.
Every fixed concentration threshold $0<\nu<1$ gives the same
all-zero trace limit as the deep selection. The estimate needed
for this assertion concerns exterior linear functionals in the
theta norm; it does not require a stronger bound for the whole
inverse Gram matrix.

On the raw space $L^2((0,\infty),du)$ write
\[
 P=P_\Lambda=1_{(0,\Lambda)},\quad
 B=H_cPH_c,\quad A=PBP,\quad F=PH_cP,\quad
 \Pi=\mathbf1_{[\nu,1]}(A).
\]
Thus $A=F^2$. Inverses of $A$ and $F$ below act only on
$\operatorname{Ran}\Pi$. Let $\mathcal R_\Lambda^\circ$ be this
range in a nonprincipal sector, and let
\[
 \mathcal R_\Lambda^\circ
 =\operatorname{Ran}\Pi\cap\ker I\cap\ker\operatorname{ev}_0
\]
in the principal sector. Denote its raw isometric inclusion by
$j_\Lambda$, and set
\[
 G_\Lambda=(E_\chi j_\Lambda)^*E_\chi j_\Lambda,\qquad
 Z_\Lambda=E_\chi j_\Lambda G_\Lambda^{-1/2},\qquad
 Q_\Lambda=Z_\Lambda Z_\Lambda^*.
\]
The Gram matrix is positive definite by the Mellin identity.
The fixed-threshold versions of
Proposition~\ref{prop:hankel-raw-estimates} follow from its proof
with $2$ replaced by $\nu^{-1}$. In particular the selected
profiles belong to the theta graph domain, are smooth up to zero
before their sharp endpoint, and admit the band extension
\begin{equation}\label{nh:extension}
 W_\Lambda=H_cPF^{-1}\Pi=BA^{-1}\Pi,\qquad
 PW_\Lambda=\Pi,\qquad \|F^{-1}\Pi\|\leq\nu^{-1/2}.
\end{equation}
The rank also has the full leading asymptotic. Indeed
\[
 |\operatorname{rank}\Pi-\Tr A|
 \leq\max\{\nu^{-1},(1-\nu)^{-1}\}\Tr(A-A^2)
 =O_{c,\nu}(\log(2+\Lambda)),
\]
by Proposition~\ref{prop:hankel-selection-size} and its variance proof.
Thus $\operatorname{rank}\Pi\sim2c\Lambda/\pi$.
All constants in this section may depend on the fixed character,
$c$, and $\nu$. Inner products are linear in their first argument.

\subsection{The two classical arithmetic estimates}

\begin{lemma}[Fixed-character reciprocal coefficients]
\label{nh:mertens}
Put
\[
 b_\chi(n)=\sum_{N\mathfrak a=n}
               \mu_K(\mathfrak a)\chi(\mathfrak a).
\]
There are constants $a_\chi,C_\chi>0$ such that
\begin{equation}\label{nh:mertens-bound}
 \left|\sum_{n\leq X}b_\chi(n)\right|
 \leq C_\chi X\exp\{-a_\chi\sqrt{\log(2+X)}\}
 \qquad(X\geq1).
\end{equation}
\end{lemma}

\begin{proof}
Apply the classical region of
\citet[Theorem~3.1, p.~1051]{thornerzaman2019} to a fixed abelian
class field through which $\chi$ factors. Its Dedekind zeta
function is the product of the associated primitive Hecke
functions. A zero of the chosen factor cannot be cancelled by
another factor away from the principal pole at one. The possible
exceptional zero is real, simple, and strictly less than one.
Since the character is fixed, decreasing a positive constant
excludes that one point. We obtain a zero-free region
\[
 \sigma\geq1-\frac{a_0}{\log(3+|t|)}
 \quad(s=\sigma+it).
\]
Nonprincipal Hecke factors are nonzero at one. In the principal
case $1/L_K(s,1)$ is holomorphic and zero there.

We record the reciprocal bound needed on the contour. For
$|t|\geq2$, put $\ell=\log(3+|t|)$ and choose $a>0$ sufficiently
small compared with $a_0$. On the disk centered at
$1+2a/\ell+it$ with radius $4a/\ell$, the analytic branch of
$\log L_K(s,\chi)$ exists. Convexity in this shrinking strip,
obtained from the functional equation and the line
$1+1/\ell$, gives
$\log|L_K(s,\chi)|=O_\chi(\log\ell)$ from above.
At the center the Euler logarithm is bounded in modulus by
$\log\zeta_K(1+2a/\ell)=O_\chi(\log\ell)$.
Borel--Carath\'eodory on the concentric radius $3a/\ell$ disk,
and the Euler product further to the right, yield
\[
 |L_K(\sigma+it,\chi)^{-1}|
       \ll_\chi(\log(3+|t|))^{C_\chi},
 \qquad 1-\frac a{\log(3+|t|)}\leq\sigma\leq2.
\]
Bounded heights follow by compactness and removability of the
principal reciprocal. The use of fixed constants here makes
no conductor-uniform assertion.

The quadratic-field ideal count gives $|b_\chi(n)|\leq d(n)$.
The divisor hyperbola argument gives, for $1\leq H\leq X$,
\[
 \sum_{X<n\leq X+H}d(n)\ll H\log(2X)+\sqrt X.
\]
Indeed, in every factor pair at least one factor is at most
$\sqrt{2X}$; sum the bounds $H/a+1$ for the other factor.
The same estimate holds for an interval immediately to the left
of $X$.

Take $X$ to be a half-integer,
$\kappa=1+1/\log X$, and $2\leq T\leq\sqrt X$.
The truncated Perron error is bounded by
\[
 \sum_{n\geq1}d(n)(X/n)^\kappa
       \min\{1,(T|\log(X/n)|)^{-1}\}.
\]
Outside $[X/2,2X]$, this is $O(X(\log X)^2/T)$.
Within that interval, first use the preceding short-interval
bound for $|n-X|\leq X/T$, then sum it over the doubled ranges
$2^jX/T<|n-X|\leq2^{j+1}X/T$.
Their total is $O(X(\log X)^2/T+\sqrt X)$. Hence
\[
 \sum_{n\leq X}b_\chi(n)
 =\frac1{2\pi i}\int_{\kappa-iT}^{\kappa+iT}
          \frac{X^s}{sL_K(s,\chi)}\,ds
 +O_\chi\left(\frac{X(\log X)^2}{T}+\sqrt X\right).
\]
Move the contour to $\sigma_0=1-a/\log(3+T)$, decreasing $a$
once more if necessary. The whole rectangle is zero-free and
there are no residues. The vertical integral is
\[
 O_\chi\left(
 X^{1-a/\log(3+T)}(\log(3+T))^{C_\chi+1}\right),
\]
and the horizontal integrals are
$O_\chi(X(\log(3+T))^{C_\chi}/T)$.
Set $T=\exp(b\sqrt{\log X})$ with a fixed $b>0$, and absorb
the logarithmic factors by decreasing the final positive
constant. This proves \eqref{nh:mertens-bound}.
Replacing an arbitrary $X$ by $\lfloor X\rfloor+1/2$ preserves
the summatory function and changes the bound by a fixed factor.
\end{proof}

\begin{lemma}[Mean square of the ideal discrepancy]
\label{nh:mean-square}
For
\[
 A_\chi(X)=\sum_{N\mathfrak a\leq X}\chi(\mathfrak a)
                        -\delta_\chi r_KX,
\]
one has
\begin{equation}\label{nh:mean-square-bound}
 \int_1^X|A_\chi(y)|^2\,dy\ll_\chi X^{3/2},
 \qquad
 \int_Y^\infty\frac{|A_\chi(y)|^2}{y^2}\,dy
                                      \ll_\chi Y^{-1/2}.
\end{equation}
\end{lemma}

\begin{proof}
We apply the established theorem of
\citet[Theorem~1, p.~62, equation~(57)]{chandrasekharannar1964}.
Its setup on pp.~43--44 permits complex coefficient sequences.
In its notation put
\[
 \varphi(s)=c^{-s}L_K(s,\chi),\qquad
 \psi(s)=\varepsilon_\chi c^{-s}L_K(s,\bar\chi),\qquad
 \Delta(s)=\Gamma(s),\qquad \delta=1.
\]
The functional equation is then
$\Delta(s)\varphi(s)=\Delta(1-s)\psi(1-s)$.
Thus its gamma parameter is $A=1$, and its two Dirichlet
scales are $\lambda_n=\mu_n=cn$. Absolute convergence and
the completed vertical-strip decay have the required form.
The only singularities are the usual finite set of poles.

Writing $a_\chi(n)=\sum_{N\mathfrak a=n}\chi(\mathfrak a)$,
we have $|a_\chi(n)|\leq d(n)$ and
\[
 \sum_{n\leq X}|a_\chi(n)|^2
 \leq\sum_{n\leq X}d_4(n)\ll X\log^3(2X).
\]
Here $d(n)^2\leq d_4(n)$ follows at each prime from
$(j+1)^2\leq\binom{j+3}{3}$.
The dual coefficients obey the same estimate.
Choose the theorem's parameters $\beta=1$, $\beta'=3$.
Its first branch explicitly includes
$2\beta-\delta-1/A=0$; the resulting parameter is
$\omega_0=1/4$. It gives
\[
 \int_1^X|E(y)|^2\,dy=C X^{3/2}+O(X\log^5X).
\]
Its residue convention gives
$E(y)=A_\chi(y/c)-L_K(0,\chi)$ away from the summation
jumps. The constant is bounded, and half weighting at the
discrete jumps changes no Lebesgue integral. Scaling back
and applying Cauchy--Schwarz proves the first assertion of
\eqref{nh:mean-square-bound}. Partial integration proves the
second assertion.
\end{proof}

\subsection{An exterior functional in the full theta norm}

\begin{lemma}[Exterior row estimate]\label{nh:row}
Let $q\in L^2((\Lambda,\infty),du)$ be extended by zero,
and put
\[
 k_q=F^{-1}\Pi PH_cq.
\]
Then $\|k_q\|_2\leq\nu^{-1/2}\|q\|_2$. For
$0<u\leq\Lambda/2$,
\begin{align}
 |k_q(u)|&\leq C\|q\|_2(1+\Lambda u)^{-1/4},
       \label{nh:pointwise}\\
 |u k_q'(u)|&\leq C\|q\|_2(1+\Lambda u)^{1/4}.
       \label{nh:derivative}
\end{align}
\end{lemma}

\begin{proof}
Let $J_u(v)=1_{(0,\Lambda)}(v)J_0(2c\sqrt{uv})$.
Since $F^{-1}\Pi=FA^{-1}\Pi$, evaluation of the finite
Hankel transform, followed by Cauchy--Schwarz, gives
\[
 |k_q(u)|^2
 \leq c^2\|q\|_2^2
       \langle J_u,A^{-2}\Pi(1-A)J_u\rangle
 \leq c^2\nu^{-2}\|q\|_2^2
       \langle J_u,(1-A)J_u\rangle.
\]
The use of $1-A$ is exact: $q$ is outside $P$ and
$P H_c(1-P)H_cP=P-A$.
Since $H_cPJ_u(v)=B_\Lambda(v,u)/c$,
\begin{equation}\label{nh:row-energy}
 |k_q(u)|^2\leq\nu^{-2}\|q\|_2^2
                 \int_\Lambda^\infty|B_\Lambda(v,u)|^2\,dv.
\end{equation}
Replacing $J_u$ by $u\partial_uJ_u$ gives the analogous
bound with $u\partial_uB_\Lambda(v,u)$.

The exact Lommel formula is
\[
 B_\Lambda(v,u)=\frac{c\sqrt\Lambda}{v-u}
 \left[
 \sqrt v J_1(2c\sqrt{\Lambda v})J_0(2c\sqrt{\Lambda u})
 -\sqrt u J_1(2c\sqrt{\Lambda u})J_0(2c\sqrt{\Lambda v})
 \right].
\]
For $v>\Lambda$ and $u\leq\Lambda/2$, the denominator is
at least $v/2$. The standard bounds for $J_0,J_1$ give
\[
 |B_\Lambda(v,u)|
 \leq C_c\Lambda^{1/4}v^{-3/4}(1+\Lambda u)^{-1/4}.
\]
Use
\[
 u\partial_uJ_0(2c\sqrt{\Lambda u})
       =-c\sqrt{\Lambda u}\,J_1(2c\sqrt{\Lambda u}),\qquad
 u\partial_u\{\sqrt u J_1(2c\sqrt{\Lambda u})\}
       =c\sqrt\Lambda\,uJ_0(2c\sqrt{\Lambda u})
\]
to obtain
\[
 |u\partial_uB_\Lambda(v,u)|
 \leq C_c\Lambda^{1/4}v^{-3/4}(1+\Lambda u)^{1/4}.
\]
The integral of $\Lambda^{1/2}v^{-3/2}$ over
$v>\Lambda$ is bounded. Substitution in
\eqref{nh:row-energy} proves both assertions.
\end{proof}

\begin{proposition}[Exterior functionals after Gram normalization]
\label{nh:functional}
For every $N>0$ and every $q$ supported in $(\Lambda,\infty)$,
\begin{equation}\label{nh:functional-bound}
 \sup_{0\ne\psi\in\mathcal R_\Lambda^\circ}
 \frac{|\langle W_\Lambda\psi,q\rangle|^2}
      {\|E_\chi\psi\|_2^2}
 \leq C_N\Lambda(\log\Lambda)^{-N}\|q\|_2^2.
\end{equation}
Equivalently, after relabelling $N$,
\begin{equation}\label{nh:normalized-extension}
 \|(1-P)W_\Lambda j_\Lambda G_\Lambda^{-1/2}\|
          \leq C_N\Lambda^{1/2}(\log\Lambda)^{-N}.
\end{equation}
\end{proposition}

\begin{proof}
Put $k=k_q$ from Lemma~\ref{nh:row}.
For $\psi\in\mathcal R_\Lambda^\circ$,
$\langle W_\Lambda\psi,q\rangle=\langle\psi,k\rangle$.
Write $b(\mathfrak a)=\mu_K(\mathfrak a)\chi(\mathfrak a)$.
Finite ideal inversion gives the exact pointwise identity
\[
 \psi(u)=\sum_{N\mathfrak a\leq\Lambda/u}
       \frac{b(\mathfrak a)}{\sqrt{N\mathfrak a\,u}}
                    (E_\chi\psi)(N\mathfrak a\,u).
\]
In the principal case the imposed mean condition makes the
corrected theta map equal to the ordinary theta sum here.
Define, for $0<v<\Lambda$,
\begin{equation}\label{nh:dual-representative}
 \beta_k(v)=\sqrt v\sum_{\mathfrak a}
       \frac{\overline{b(\mathfrak a)}}{N\mathfrak a}
                                      k(v/N\mathfrak a),
\end{equation}
and set $\beta_k=0$ above $\Lambda$.
We will prove
\begin{equation}\label{nh:dual-bound}
 \|\beta_k\|_{L^2(d^\times v)}^2
       \leq C_N\Lambda(\log\Lambda)^{-N}\|q\|_2^2.
\end{equation}

Separate the ideal of norm one. Its term is $\sqrt v\,k(v)$
and has squared $L^2(d^\times v)$ norm $\|k\|_2^2$.
For every remaining ideal, $v/N\mathfrak a\leq\Lambda/2$,
so \eqref{nh:pointwise}--\eqref{nh:derivative} apply.
Put $X=\Lambda v$. For large $X$ split the remaining sum at
$Y=X/(\log X)^D$, with $D$ a fixed number to be chosen.
The absolute ideal count is $O(y)$, and hence
\[
 \sum_{2\leq N\mathfrak a\leq Y}
     \frac{|k(v/N\mathfrak a)|}{N\mathfrak a}
 \ll\|q\|_2 X^{-1/4}Y^{1/4}
 =\|q\|_2(\log X)^{-D/4}.
\]
For the other terms apply Lemma~\ref{nh:mertens} and Abel
summation to $w(y)=y^{-1}k(v/y)$. Its derivative satisfies
\[
 |w'(y)|\leq C\|q\|_2y^{-2}(1+X/y)^{1/4}.
\]
The integral contribution is therefore bounded by
\[
 C\|q\|_2\int_Y^\infty
 e^{-a_\chi\sqrt{\log(2+y)}}(1+X/y)^{1/4}\frac{dy}{y},
\]
with a smaller boundary term. Since $X/Y=(\log X)^D$,
this is smaller than every fixed power of $(\log X)^{-1}$.
For bounded $X$, the same Abel argument gives $O(\|q\|_2)$.
Consequently the part of \eqref{nh:dual-representative}
from ideals of norm at least two obeys
\[
 |\beta_{k,\geq2}(v)|
 \leq C_N\|q\|_2\sqrt v\,
               [\log(2+\Lambda v)]^{-N}.
\]
Integrate its square against $dv/v$. The part
$v\leq\sqrt\Lambda$ is $O(\sqrt\Lambda\|q\|_2^2)$;
on the rest, $\log(2+\Lambda v)\asymp\log\Lambda$.
Adding the norm-one term and relabelling $N$ proves
\eqref{nh:dual-bound}.

We justify the dual identity without an exchange of
conditionally convergent sums. First integrate
$\langle\psi,k\rangle$ only over $u\geq\epsilon$.
Finite inversion and the substitution $v=N\mathfrak a\,u$
then give
\[
 \int_\epsilon^\Lambda\psi(u)\overline{k(u)}\,du
 =\int_0^\Lambda E_\chi\psi(v)\,
       \overline{\sqrt v
          \sum_{N\mathfrak a\leq v/\epsilon}
          \frac{\overline{b(\mathfrak a)}}{N\mathfrak a}
                         k(v/N\mathfrak a)}\,\frac{dv}{v}.
\]
The sum is finite. For fixed $\Lambda$, the truncated dual
functions converge to $\beta_k$ in $L^2(d^\times v)$.
Indeed, on $v\geq\delta>0$, Abel summation and
\eqref{nh:mertens-bound} make the reciprocal-coefficient
tail tend uniformly to zero. On $v<\delta$, smoothness of
$k$ at zero and the same Abel estimate give a uniform
$O_\Lambda(\sqrt v)$ bound. Its squared integral tends to
zero with $\delta$. The raw left side converges by
Cauchy--Schwarz. We have proved
\[
 \langle\psi,k\rangle_{L^2(du)}
       =\langle E_\chi\psi,\beta_k\rangle_{L^2(d^\times v)}.
\]
The estimate \eqref{nh:functional-bound} now follows from
\eqref{nh:dual-bound}. Duality gives
\eqref{nh:normalized-extension}.
\end{proof}

\subsection{The full smoothed theta tail}

\begin{lemma}[A raw Hilbert--Schmidt kernel]\label{nh:raw-tail}
Let $L=\log\Lambda$ and $g\in C_c^\infty(\mathbb R)$.
The operator
\[
 T_{\Lambda,g}=V(g)1_{t<-L}E_\chi(1-P)
\]
extends from the raw graph core to a Hilbert--Schmidt operator,
and
\begin{equation}\label{nh:raw-tail-bound}
 \|T_{\Lambda,g}\|_{\mathrm{HS}}
             \leq C_\chi\Lambda^{-1/2}
                             (\|g\|_2+\|g'\|_2).
\end{equation}
The same estimate holds for $g\in L^1(\mathbb R)\cap H^1(\mathbb R)$.
\end{lemma}

\begin{proof}
Use output coordinate $r=t+L$. For $u>\Lambda$, the raw
kernel is the Stieltjes integral
\[
 K_g(r,u)=u^{-1/2}\int_{\Lambda u}^\infty
        y^{-1/2}g(r+\log(y/(\Lambda u)))\,dA_\chi(y).
\]
For compact smooth $g$, the original sum is finite at each
fixed $(r,u)$. The principal density subtraction gives exactly
the discrepancy $A_\chi$ in this formula.
Integration by parts gives an endpoint
\[
 -\Lambda^{-1/2}u^{-1}A_\chi(\Lambda u)g(r)
\]
and an integral term
\[
 -\Lambda^{-1/2}u^{-1}
       \int_1^\infty A_\chi(\Lambda u z)z^{-3/2}
                    (g'-g/2)(r+\log z)\,dz.
\]
By Lemma~\ref{nh:mean-square}, the squared norm of the
endpoint is
\[
 \|g\|_2^2\int_{\Lambda^2}^\infty
                  |A_\chi(y)|^2y^{-2}\,dy
                       \ll_\chi\Lambda^{-1}\|g\|_2^2.
\]
Minkowski's inequality bounds the norm of the integral term by
\[
 \|g'-g/2\|_2
 \int_1^\infty z^{-3/2}
 \left(z\int_{z\Lambda^2}^\infty
                       |A_\chi(y)|^2y^{-2}\,dy\right)^{1/2}dz
 \ll_\chi\Lambda^{-1/2}\|g'-g/2\|_2
                         \int_1^\infty z^{-5/4}dz.
\]
This proves \eqref{nh:raw-tail-bound}.
The raw-kernel extension agrees with the theta graph
operator by approximation in $H_D^{3/8}$ on the exterior
half-line and the graph bound \eqref{eq:hecke-graph}.
Approximation of $g$ in $L^1\cap H^1$ gives the last assertion.
The $L^1$ condition ensures that convolution is bounded on
theta space; no such assertion is needed for an arbitrary
$H^1$ function.
\end{proof}

\begin{proposition}[Smoothed leakage at a fixed threshold]
\label{nh:leakage}
For every $N>0$ and every
$g\in C_c^\infty(\mathbb R)+\mathbb C\,(\tfrac12e^{-|t|})$,
\begin{equation}\label{nh:leakage-bound}
 \|V(g)(1-S_L)Z_\Lambda\|_{\mathrm{HS}}
                         \ll_{N,g}L^{-N}.
\end{equation}
\end{proposition}

\begin{proof}
The upper theta tail of a selected raw profile is zero.
In the principal case this uses $I\psi=0$.
By \eqref{nh:extension}, $H_cW_\Lambda\psi=F^{-1}\psi$
is supported in $(0,\Lambda)$; in the principal case
\[
 I(H_cW_\Lambda\psi)=W_\Lambda\psi(0)/c=\psi(0)/c=0.
\]
The graph-domain Poisson identity
$E_\chi H_c=\varepsilon_\chi J E_{\bar\chi}$ therefore
makes the lower theta tail of $W_\Lambda\psi$ zero.
Consequently
\[
 V(g)1_{t<-L}E_\chi\psi
       =-T_{\Lambda,g}(1-P)W_\Lambda\psi.
\]
Apply the Hilbert--Schmidt bound \eqref{nh:raw-tail-bound}
and the operator bound \eqref{nh:normalized-extension}.
Their factors $\Lambda^{-1/2}$ and $\Lambda^{1/2}$ cancel.
The stated Green kernel belongs to $L^1\cap H^1$,
so Lemma~\ref{nh:raw-tail} covers it as well.
\end{proof}

\subsection{All zero jets and the trace limit}

\begin{proposition}[Fixed zero packets]\label{nh:zero-packets}
Let $\mathcal V_{\mathcal F,L}$ be the truncated
exponential-polynomial space associated with a fixed finite
packet $\mathcal F$ of nontrivial zeros of $L_K(s,\chi)$,
including their full multiplicities. Then, for every $N>0$,
\begin{equation}\label{nh:zero-packet-bound}
 \|Q_\Lambda P_{\mathcal V_{\mathcal F,L}}\|_{\mathrm{HS}}^2
                                  \ll_{\mathcal F,N}L^{-N}.
\end{equation}
\end{proposition}

\begin{proof}
Fix $\rho=\beta+i\gamma$ in the packet and $0\leq j<m_\rho$.
The lower Mellin-tail kernel is
\[
 H_{\rho,j,\Lambda}(u)=
 \left.(\partial_s+L)^j\left[
 u^{s-1}\left(L_K(s,\chi)
       -\sum_{N\mathfrak a\leq\Lambda u}
                   \chi(\mathfrak a)(N\mathfrak a)^{-s}\right)
 +\delta_\chi\frac{r_K\Lambda^{1-s}}{1-s}
 \right]\right|_{s=\rho}.
\]
Proposition~\ref{prop:hecke-improper-mellin} justifies its
improper integral and all fixed derivatives. The kernel is in
raw $L^2$: below $\Lambda^{-1}$ its $L$-terms vanish at the
zero and its jets; above that point the endpoint discrepancy
and the Riesz estimate give the bounds used in
Proposition~\ref{prop:hecke-deep-zero-jets}.

We improve its exterior norm using
Lemma~\ref{nh:mean-square}. At $j=0$,
\[
 H_{\rho,0,\Lambda}(u)
 =-\Lambda^{-\rho}u^{-1}A_\chi(\Lambda u)
 +\rho u^{\rho-1}\int_{\Lambda u}^\infty
                             A_\chi(y)y^{-\rho-1}\,dy.
\]
The first term has squared exterior norm
\[
 \Lambda^{1-2\beta}\int_{\Lambda^2}^\infty
                         |A_\chi(y)|^2y^{-2}\,dy
                                      \ll\Lambda^{-2\beta}.
\]
For the integral use \eqref{eq:hecke-riesz}. One integration
by parts, valid for every fixed $\beta>0$, gives
\[
 \int_X^\infty A_\chi(y)y^{-\rho-1}\,dy
 =-\frac{\delta_\chi}{6\rho}X^{-\rho}
                         +O_{\chi,\rho}(X^{-\beta-1/4}).
\]
After multiplication by $\rho u^{\rho-1}$, the principal
term has squared exterior norm $O(\Lambda^{-2\beta-1})$.
The remainder has squared norm
\[
 O\left(\Lambda^{-2\beta-1/2}
                      \int_\Lambda^\infty u^{-5/2}\,du\right)
                         =O(\Lambda^{-2\beta-2}).
\]
For $j\geq1$, the operator $(\partial_s+L)^j$ kills the
endpoint factor $\Lambda^{-s}$ and the principal constant
term. Substitute $y=(\Lambda u)z$ in the Riesz remainder;
the remaining derivatives introduce only fixed powers of
$\log z$, integrable since $\beta+1/4>0$. In particular
\begin{equation}\label{nh:zero-tail-norm}
 \|1_{u>\Lambda}H_{\rho,j,\Lambda}\|_2^2
                                  \ll_{\rho,j}\Lambda^{-2\beta}.
\end{equation}

These kernels annihilate the fixed Hankel band, with value
zero at the origin in the principal case. Indeed, first use
inverse Hankel transforms of smooth compact band profiles;
Poisson makes their lower theta tail zero, and the disk
superposition proof of the improper Mellin formula applies
to these rapidly decreasing profiles. Such vectors are dense
in the band space. Evaluation at zero is continuous on that
space, so the principal value condition can be preserved.
The raw $L^2$ kernel extends the annihilation by continuity.
In particular its pairing with $W_\Lambda\psi$ is zero.
Its pairing with $\psi$ is therefore the negative of its
exterior pairing with $(1-P)W_\Lambda\psi$.

Apply Proposition~\ref{nh:functional} with
$q=1_{u>\Lambda}H_{\rho,j,\Lambda}$, taking conjugates
according to the Mellin pairing. The squared theta-dual norm
of the lower Mellin jet is
\[
 O_{\rho,j,N}\bigl(\Lambda^{1-2\beta}L^{-N}\bigr).
\]
The full Mellin jet vanishes at $\rho$, and the upper theta
tail is zero. Thus the same functional, with opposite sign,
is its window integral.

Put $\alpha=\beta-1/2$. For $\alpha<0$, use
$1_{[-L,L]}e^{(\alpha-i\gamma)(t+L)}(t+L)^j$.
The additional squared factor $\Lambda^{2\beta-1}$
cancels the displayed power. For $\alpha>0$, use the
right-endpoint normalization
$1_{[-L,L]}e^{(\alpha-i\gamma)(t-L)}(t-L)^j$.
Expand $(t-L)^j=((t+L)-2L)^j$; the additional factor is
$\Lambda^{1-2\beta}$, giving
$\Lambda^{2-4\beta}L^{O(j)-N}$.
On the central line, normalized polynomial jets introduce
only powers of $L$.
The finite endpoint Gram matrices have nonsingular limits;
opposite endpoints decouple, and the central polynomial blocks
have their nonsingular limiting Gram matrices.
Since $N$ is arbitrary, these fixed changes of basis prove
\eqref{nh:zero-packet-bound}.
\end{proof}

\begin{proposition}[Transfer to the uncompressed trace]
\label{nh:trace-transfer}
For every $h\in C_c^\infty(\mathbb R)$,
\[
 \Tr((Q_\Lambda-S_LQ_\Lambda S_L)V(h))\longrightarrow0.
\]
\end{proposition}

\begin{proof}
Let $\kappa(t)=\tfrac12e^{-|t|}$ and
$b=(1-\partial_t^2)h$. Then $h=b*\kappa$.
The isometry of $Z_\Lambda$ gives, for $g\in L^2$,
\[
 \|V(g)S_LZ_\Lambda\|_{\mathrm{HS}}
 \leq\|V(g)S_L\|_{\mathrm{HS}}
                         =\sqrt{2L}\,\|g\|_2.
\]
For $g$ in the class of Proposition~\ref{nh:leakage},
the same bound for $V(g)Z_\Lambda$ has an additional
$O_N(L^{-N})$ term.
Write $P_o=1-S_L$ and $b^*(t)=\overline{b(-t)}$.
Cyclicity of the finite-rank trace and
$Q-S_LQS_L=P_oQ+S_LQP_o$ give
\begin{align*}
 |\Tr((Q-S_LQS_L)V(h))|
 &\leq\|V(b^*)Z_\Lambda\|_{\mathrm{HS}}\,
                  \|V(\kappa)P_oZ_\Lambda\|_{\mathrm{HS}}\\
 &\quad+\|V(b^*)P_oZ_\Lambda\|_{\mathrm{HS}}\,
                  \|V(\kappa)S_LZ_\Lambda\|_{\mathrm{HS}}\\
 &\ll_{N,h}L^{1/2-N}+L^{-2N}.
\end{align*}
Here the two trace pairings are respectively
$Z_\Lambda^*V(b)V(\kappa)P_oZ_\Lambda$ and
$Z_\Lambda^*P_oV(b)V(\kappa)S_LZ_\Lambda$.
This proves the assertion without a square-root rank loss.
\end{proof}

\begin{theorem}[The fixed-threshold Hecke trace formula]
\label{nh:theorem}
For every fixed $0<\nu<1$, the compressed defect measures
associated with $\Pi=\mathbf1_{[\nu,1]}(A_\Lambda)$ converge
on all Schwartz frequency tests to the all-zero harmonic
measure $\mathfrak b_\chi$, with full zero multiplicities.
For every $h\in C_c^\infty(\mathbb R)$,
\begin{equation}\label{nh:trace-limit}
 2\log\Lambda\,h(0)-\Tr(Q_\Lambda V(h))
                              \longrightarrow\mathfrak B_\chi(h).
\end{equation}
The natural threshold $\nu=1/2$ is included.
In the principal full completion, the same rank-two polar
correction as in Proposition~\ref{prop:hecke-principal-graph}
must be added. No hypothesis on the positions or
multiplicities of the nontrivial zeros is used.
\end{theorem}

\begin{proof}
For a fixed finite zero packet put
$P_V=P_{\mathcal V_{\mathcal F,L}}$ and
$D_\Lambda=S_L-S_LQ_\Lambda S_L$.
Then $0\leq D_\Lambda\leq1$, $P_V\leq S_L$, and
\[
 \|(1-D_\Lambda^{1/2})P_V\|_{\mathrm{HS}}^2
 \leq\Tr(P_V(1-D_\Lambda)P_V)
                    =\|Q_\Lambda P_V\|_{\mathrm{HS}}^2=o(1)
\]
by Proposition~\ref{nh:zero-packets}.
For any nonnegative compactly supported continuous frequency test $\varphi$,
\[
 \Tr(D_\Lambda C_\varphi)
 =\|D_\Lambda^{1/2}C_\varphi^{1/2}\|_{\mathrm{HS}}^2
 \geq\|P_VD_\Lambda^{1/2}C_\varphi^{1/2}\|_{\mathrm{HS}}^2
 =\Tr(P_VC_\varphi)+o(1).
\]
Theorem~\ref{thm:finite-packet} identifies the last limit
with the harmonic measure of that packet. Increase the finite
packet and use positivity and the zero-count summability to
obtain the all-zero lower bound.

For large $\Lambda$, the deep raw selected space is contained
in the present fixed-threshold space, also after the principal
constraints. Their theta-image projections are therefore
nested, and
\[
 0\leq D_\Lambda^{(\nu)}\leq D_\Lambda^{\mathrm{deep}}.
\]
Theorem~\ref{thm:hi-local-limit} gives the matching compact
frequency upper bound. The uniform polynomial mass estimate
of Theorem~\ref{thm:hi-uniform-mass} passes through the same
order inequality. Equality of the compact limits, polarization,
and this uniform mass bound give convergence on every Schwartz
frequency test.
The compressed trace identity and
Proposition~\ref{nh:trace-transfer} prove
\eqref{nh:trace-limit}.

Finally take the two raw principal coordinate-repair vectors
already used for the deep cutoff in
Proposition~\ref{prop:hecke-principal-graph}.
They also belong to the present raw selected space. Let
$Y_\Lambda:\mathbb C^2\to L^2$ be their normalized theta
columns. The two graph maps are
\[
 X_\Lambda^{\mathrm{deep}}
       =(1-Q_\Lambda^{\mathrm{deep}})Y_\Lambda,\qquad
 X_\Lambda^{(\nu)}=(1-Q_\Lambda^{(\nu)})Y_\Lambda.
\]
Nesting gives
$\|X_\Lambda^{(\nu)}\|\leq
\|X_\Lambda^{\mathrm{deep}}\|\to0$.
The rank-two graph-projection formula consequently gives
trace-norm convergence of the full-principal correction to
the vertical $\mathbb C^2$ projection. Its integrated trace is
\[
 \int_{\mathbb R}h(t)(e^{-t/2}+e^{t/2})\,dt
\]
in the stated principal coordinate normalization. This proves
the full-sector assertion.
\end{proof}


Theorem~\ref{nh:theorem}, followed by the finite-character sum
above, proves the trace and measure assertions of
Theorem~\ref{thm:main-hecke}. The following comparison proves
its half-plane assertion and makes the conductor dependence explicit.

\subsection{A Hecke half-plane and the cubic example}
\begin{corollary}[What a Hecke half-plane gives]
\label{cor:hecke-band}
Fix a finite-order Hecke character $\chi$ and assume that
$L_K(s,\chi)$ has no zeros for $\Re s>\theta$, where
$1/2\leq\theta<1$. Put $\delta=\theta-1/2$ and
\[
 Z_{2,\chi}=\sum_\rho\frac{m_\rho}{1+(\Im\rho)^2}.
\]
For every logarithmic test $h_\chi\in C_c^\infty(\R)$,
\begin{equation}\label{eq:hecke-band}
 |\mathfrak B_\chi(h_\chi)-\mathfrak W_\chi(h_\chi)|
     \leq2\delta\mathcal N_\delta(h_\chi)Z_{2,\chi}.
\end{equation}
The sum is finite and, for finite-order characters of the fixed field,
$Z_{2,\chi}\ll_K\log(2+Q_\chi)$. By
Theorem~\ref{thm:main-hecke}, $\mathfrak B_\chi$ is the limit of
the trace defect of every fixed-threshold Hankel selection, so the same
inequality bounds that trace limit's deviation from Weil's distribution.
\end{corollary}

\begin{proof}
Conjugation of coefficients followed by the functional equation
reflects the zero divisor by $\rho\mapsto1-\overline\rho$. The
half-plane therefore confines it to $|\Re\rho-1/2|\leq\delta$.
Here is the conductor-uniform zero count needed for the second assertion.
Put $F_\chi(s)=(s-1)^{\mathbf1_{\chi=1}}L_K(s,\chi)$, which is
entire. The absolutely convergent Euler product on $\Re s=2$ gives
$|F_\chi(2+iT)|\geq\zeta_K(2)^{-1}$. The same product bounds
$L_K$ uniformly to the right of this line. The functional equation,
Stirling's formula on the left boundary of a fixed strip, and the
Phragm\'en--Lindel\"of theorem between those boundaries give
\[
 \max_{|s-(2+iT)|\leq4}|F_\chi(s)|
                    \leq C_K\bigl(Q_\chi(3+|T|)^2\bigr)^C,
\]
with a fixed exponent $C$. The removal of the principal pole makes
the same bound valid for bounded $T$; the gamma factors and degree
are fixed throughout this family. Jensen's formula between the
concentric radii three and four therefore bounds the number of zeros
with $0<\Re\rho<1$ and $|\Im\rho-T|\leq1/2$ by
$C_K\log(Q_\chi(3+|T|)^2)$, counting multiplicities. Summing these
unit-interval bounds gives
$N_\chi(T)\ll_K(T+1)\log(Q_\chi(T+3)^2)$. Summation against
$(1+T^2)^{-1}$ now gives
$Z_{2,\chi}\ll_K\log(2+Q_\chi)$.
Apply Theorem~\ref{thm:band-bound} to this divisor and the trace
limit proved in Theorem~\ref{nh:theorem}.
\end{proof}

The external choices of $\theta$ have distinct provenance.
\citet[Theorem~7.33, p.~66]{liu2026} states
\[
 .8733326374459049352<\theta_G<.8733326374459051924
\]
for every continuous unitary Hecke character; \citet[Theorems~1.1--1.2,
pp.~1--2]{hu2026} gives $\theta=7/8$. Formula \eqref{eq:hecke-band}
is a deduction from whichever half-plane is assumed. It involves
$\delta\approx.373332637445905$ in the former case and $3/8$ in
the latter. Both the explicit norm $\mathcal N_\delta$ and the factor
$\delta$ decrease when the band narrows. A conductor-free bound
cannot be inferred from band confinement alone: repeated copies of
an off-line reflected pair scale its contribution arbitrarily.
This observation concerns the logical strength of the hypothesis,
not the existence of arithmetic off-line zeros.

For the cubic character defining
$K(\sqrt[3]{269})/K$, the conductor is $(269)$ and
$Q_\chi=3\cdot269^2=217083$. Its two conjugate characters give the
same $L$-function, namely
$\zeta_{\Q(\sqrt[3]{269})}(s)/\zeta(s)$, by induction and
conjugation of ideals. Thus \eqref{eq:hecke-band} applies to this
particular degree-two function. In an arithmetic reduction whose
content exponent is
\[
 \gamma(\beta,\rho)
       =\frac{\beta-1/2}{2\beta-\rho/2},
\]
the parameters mean that the incomplete cubic Hecke function is
zero-free in $\Re s>\beta$ and zeta is zero-free in $\Re s>\rho$.
The reduction is the ``Joint-content tail'' theorem, with its
``Exact remaining square-root criterion'', in \citet{cubiccontentnotes}.
A Hecke boundary $\beta$ and rational zeta RH give
$\gamma(\beta,1/2)=(\beta-1/2)/(2\beta-1/4)$.
The choices $\beta=7/8$ and $\beta=\theta_G$ yield, respectively,
$1/4$ and $0.249442974129865\ldots$. The derivative with respect
to $\beta$ is $\tfrac34(2\beta-1/4)^{-2}>0$.
These substitutions use the parameterized reduction as an explicit
input; they do not re-prove its arithmetic cancellation assertion.
A balanced ideal convolution and the harmonic trace deviation are
different quantities. Equivalence of their vanishing criteria, when
all comparison theorems are available, supplies no identity between
their numerical values.
For this fixed cubic character, equality of the Hankel cutoff limit
with $\mathfrak W_\chi$ on every compact smooth test is equivalent to
RH by the positivity criterion. The half-plane bound controls its
possible deviation but does not force that deviation to vanish.

\section{Cutoffs for primitive holomorphic newforms}
\label{hnew:section}

The angular Fourier transform gives the corresponding cutoff for every
fixed primitive holomorphic cuspidal newform.  Weight one then gives
a cutoff for every odd irreducible two-dimensional Artin representation
over $\mathbb Q$ through the classical modularity theorem.

Fix a normalized primitive holomorphic cuspidal newform
$f\in S_k(\Gamma_0(N),\xi)$, where $k\ge1$ is an integer,
$N\ge1$, and $\xi$ is the nebentypus, extended by zero at integers
not coprime to $N$.  Write
\[
 f(z)=\sum_{n\ge1}\lambda_f(n)n^{(k-1)/2}e^{2\pi inz},
 \qquad \lambda_f(1)=1,\qquad
 a=\frac{k-1}{2},\quad \nu=k-1,\quad c=\frac{2\pi}{\sqrt N}.
\]
The coefficients may be complex.  Denote by $\bar f$ the primitive
dual form, with $\lambda_{\bar f}(n)=\overline{\lambda_f(n)}$,
the same weight and level, and nebentypus $\bar\xi$.
All constants below may depend on the fixed form and, where applicable,
the fixed concentration threshold; no uniformity
in its weight, level, or nebentypus is asserted.  Put
\begin{equation}\label{hnew:completed}
 L_f(s)=\sum_{n\ge1}\lambda_f(n)n^{-s},\qquad
 \Lambda_f(s)=c^{-s}\Gamma(s+a)L_f(s).
\end{equation}
Both $L_f$ and the completed function are entire, and
\[
 \Lambda_f(s)=\varepsilon_f\Lambda_{\bar f}(1-s),\qquad
 |\varepsilon_f|=1,\qquad
 \varepsilon_{\bar f}=\overline{\varepsilon_f}.
\]
The standard completed normalization
$N^{s/2}\Gamma_{\mathbb C}(s+a)L_f(s)$ differs from
$\Lambda_f(s)$ by the constant $2(2\pi)^{-a}$, which cancels
from the functional equation; see
\citet[pp.~1--3 of the author preprint]{booker2019}.

On the raw space $\mathcal H=L^2((0,\infty),du)$ define
\[
 (\mathcal H_{c,\nu}\psi)(v)
   =c\int_0^\infty J_\nu(2c\sqrt{uv})\psi(u)\,du,\qquad
 (U\psi)(t)=e^{t/2}\psi(e^t),\qquad H_D^s=U^{-1}H^s(\mathbb R).
\]
The transform is a unitary involution, as verified below.
Set
\[
 P_\Lambda=1_{(0,\Lambda)},\quad
 B_\Lambda=\mathcal H_{c,\nu}P_\Lambda\mathcal H_{c,\nu},\quad
 A_\Lambda=P_\Lambda B_\Lambda P_\Lambda,\quad L=\log\Lambda.
\]
For a fixed $0<\vartheta<1$ and for the auxiliary deep cutoff put
\begin{equation}\label{hnew:selections}
 \Pi_\Lambda^{(\vartheta)}=1_{[\vartheta,1]}(A_\Lambda),
 \qquad
 \Pi_\Lambda^{\mathrm d}=1_{[1-\Lambda^{-56},1]}(A_\Lambda).
\end{equation}
Define the theta operator first on smooth compact profiles by
\begin{equation}\label{hnew:theta}
 (E_f\psi)(t)=e^{t/2}\sum_{n\ge1}\lambda_f(n)\psi(ne^t).
\end{equation}
Its closed multiplier realization on the graph domain used below is
$\widehat{E_f\psi}(\tau)=L_f(1/2+i\tau)\widehat{U\psi}(\tau)$.
Here $\widehat g(\tau)=\int_{\mathbb R}g(t)e^{i\tau t}dt$,
with Plancherel measure $d\tau/(2\pi)$.

For either selection let $j_\Lambda:\mathbb C^{r_\Lambda}\to\mathcal H$
be any raw orthonormal column onto its range.  The full image Gram
matrix, orthonormal theta column, and image projection are
\begin{equation}\label{hnew:gram-definition}
 G_\Lambda=(E_fj_\Lambda)^*E_fj_\Lambda,\qquad
 Z_\Lambda=E_fj_\Lambda G_\Lambda^{-1/2},\qquad
 Q_\Lambda=Z_\Lambda Z_\Lambda^*.
\end{equation}
Their existence and independence of the raw basis follow from the
Gram bound below.  There is no principal sector or polar-coordinate
completion, including when $k=1$ and $\nu=0$.
In particular neither $\psi(0)=0$ nor $\int\psi=0$ is imposed.
Put $S_L=1_{[-L,L]}$ and
$D_\Lambda=S_L-S_LQ_\Lambda S_L$.  For an orthonormal image basis
$z_1,\ldots,z_{r_\Lambda}$ define the positive measure
\begin{equation}\label{hnew:measure}
 d\mu_\Lambda(\tau)=
 \left(2L-\sum_i|\widehat{S_Lz_i}(\tau)|^2\right)\frac{d\tau}{2\pi}.
\end{equation}
Its positivity is Bessel's inequality.

For a nontrivial zero $\rho=\beta+i\gamma$ of $L_f$, with its
analytic multiplicity $m_\rho$, let $\mu_\rho$ be the Cauchy
probability measure
\[
 d\mu_\rho(\tau)=
 \frac{|\beta-1/2|\,d\tau}
 {\pi((\tau-\gamma)^2+|\beta-1/2|^2)}
 \quad(\beta\ne1/2),\qquad
 \mu_\rho=\delta_\gamma\quad(\beta=1/2).
\]
Write $\mathfrak b_f=\sum_\rho m_\rho\mu_\rho$ and
$\mathfrak B_f(h)=\int H(i\tau)d\mathfrak b_f(\tau)$, where
$H(z)=\int_{\mathbb R}h(t)e^{zt}dt$.

\begin{theorem}[The holomorphic newform cutoff]
\label{hnew:theorem}
For every fixed primitive form $f$ as above and every fixed
$0<\vartheta<1$, the full-Gram cutoff
$Q_\Lambda^{(\vartheta)}$ has the following properties.
\begin{enumerate}
\item Its raw selected rank is
$4\Lambda/\sqrt N+O_{f,\vartheta}(\log(2+\Lambda))$, and
$G_\Lambda^{-1}=O_{f,\vartheta}(\Lambda^2\log^2(2+\Lambda))$.
\item The measures \eqref{hnew:measure} converge to $\mathfrak b_f$
against every Schwartz frequency test.  For some fixed $C_{f,\vartheta},A_f>0$,
\[
 \mu_\Lambda([T,T+1])\le C_{f,\vartheta}(2+|T|)^{A_f}
       \qquad(\Lambda\ge2,\ T\in\mathbb R).
\]
\item If $(V(h)g)(t)=\int_{\mathbb R}h(r)g(t-r)dr$, then for
every $h\in C_c^\infty(\mathbb R)$,
\begin{equation}\label{hnew:trace-theorem}
 2\log\Lambda\,h(0)-\operatorname{Tr}(Q_\Lambda V(h))
             \longrightarrow\mathfrak B_f(h).
\end{equation}
\end{enumerate}
The measure and trace assertions also hold for the deep selection
in \eqref{hnew:selections}.  No hypothesis on the positions,
multiplicities, or separation of the nontrivial zeros is imposed.
If $L_f$ is zero free in $\Re s>\sigma_*$, with
$1/2\le\sigma_*<1$ and $\delta=\sigma_*-1/2$, then
\[
 |\mathfrak B_f(h)-\mathfrak W_f(h)|
 \le2\delta\mathcal N_\delta(h)
           \sum_\rho\frac{m_\rho}{1+\gamma^2},\qquad
 \mathfrak W_f(h)=\sum_\rho m_\rho H(\rho-1/2),
\]
where
\[
 \mathcal N_\delta(h)=|h(0)|+
 \int_{\mathbb R}e^{\delta|t|}
 \bigl(|t||h''(t)|+2|h'(t)|+(1+\delta^2)|t||h(t)|\bigr)dt.
\]
\end{theorem}

The leading term is $2\log\Lambda$, despite the degree two
$L$-function and rank $4\Lambda/\sqrt N+O(\log\Lambda)$: the logarithmic
window has length $2\log\Lambda$ and one scalar theta channel.
There are no terms $H(1/2)$ or $H(-1/2)$.
We prove the theorem after establishing the required inputs.

\subsection{The angular transform and its sharp cutoff estimates}

\begin{lemma}[Normalization and domain]\label{hnew:transform}
The transform $\mathcal H_{c,\nu}$ is a unitary involution.  With
$\mathcal M\psi(s)=\int_0^\infty\psi(u)u^{s-1}du$, one has
\begin{equation}\label{hnew:mellin-hankel}
 \mathcal M(\mathcal H_{c,\nu}\psi)(s)
  =g_{c,\nu}(s)\mathcal M\psi(1-s),\qquad
 g_{c,\nu}(s)=c^{1-2s}
      \frac{\Gamma(s+a)}{\Gamma(1-s+a)}.
\end{equation}
Consequently
\begin{equation}\label{hnew:poisson}
 E_f\mathcal H_{c,\nu}=\varepsilon_f J E_{\bar f},\qquad
 (Jg)(t)=g(-t).
\end{equation}
The transform reverses raw dilations and preserves every $H_D^s$.
Its smooth angular core consists of profiles $u^ab(u)$ near zero,
with $b$ smooth, and with Schwartz decay at infinity.
\end{lemma}
\begin{proof}
The lift $\psi(u)\mapsto e^{i\nu\theta}\psi(r^2)$ has norm
$\sqrt\pi$ times the raw norm and identifies the raw space with
the angular-frequency-$\nu$ subspace of $L^2(\mathbb R^2)$.
Euclidean Fourier transformation preserves that subspace.  Removing
its constant phase $(-i)^\nu$ and making the fixed radial rescaling
gives $\mathcal H_{c,\nu}$ and proves unitarity and involutivity.
On a convergence strip, substitute $z=2c\sqrt{uv}$ and use
\[
 \int_0^\infty z^{2s-1}J_\nu(z)\,dz
     =2^{2s-1}\frac{\Gamma(s+\nu/2)}{\Gamma(1-s+\nu/2)}
\]
to obtain \eqref{hnew:mellin-hankel}.  Continuation and the unitary
Mellin realization give it on the critical line.
Equation \eqref{hnew:completed} gives
$L_f(s)g_{c,\nu}(s)=\varepsilon_fL_{\bar f}(1-s)$.
Mellin inversion proves \eqref{hnew:poisson} on the angular core:
its Mellin transforms decrease rapidly on closed strips to the
right of $-a$, permitting a contour shift from $\Re s>1$ to
the critical line.  A shift slightly to the left also proves
square integrability at the lower end.  The graph bound in
Lemma~\ref{hnew:arithmetic} extends the identity to the domains used.
Finally, for $(D_r\psi)(u)=r^{-1/2}\psi(u/r)$, substitution gives
$\mathcal H_{c,\nu}D_r=D_{1/r}\mathcal H_{c,\nu}$.
The dilation generator is therefore reversed, proving the Sobolev
assertion.  Smoothness of the angular lift at the origin gives
the stated form of its core.
\end{proof}
\begin{proposition}[Fixed-order geometric estimates]\label{hnew:geometry}
For fixed $\nu$ and $c$,
\begin{align}
 \operatorname{Tr}A_\Lambda
    &=\frac{2c\Lambda}{\pi}+O_{c,\nu}(\log(2+\Lambda)),\label{hnew:trace}\\
 \operatorname{Tr}(A_\Lambda-A_\Lambda^2)
    &=O_{c,\nu}(\log(2+\Lambda)).\label{hnew:variance}
\end{align}
Every fixed-threshold selected vector has a band extension
\[
 W_\Lambda\psi
 =B_\Lambda A_\Lambda^{-1}\psi
 =\mathcal H_{c,\nu}P_\Lambda F_\Lambda^{-1}\psi,\qquad
 F_\Lambda=P_\Lambda\mathcal H_{c,\nu}P_\Lambda,
\]
such that $P_\Lambda W_\Lambda\psi=\psi$ and
$\|W_\Lambda\|\le\vartheta^{-1/2}$.  For some fixed $d>0$,
\begin{equation}\label{hnew:raw-mass-graph}
 \int_{d/\Lambda}^\Lambda|\psi(u)|^2du
   \ge\tfrac12\|\psi\|_2^2,\qquad
 \|\psi\|_{H_D^s}\le C_s\Lambda^s\|\psi\|_2
       \quad(0<s<1/2).
\end{equation}
For a smooth strict-margin $p$, equal to one near zero and supported
in $[0,b]$, $b<1$, write
$M_\Lambda=p(u/\Lambda)$,
$B_\Lambda^p=\mathcal H_{c,\nu}M_\Lambda\mathcal H_{c,\nu}$, and
$K_\Lambda^p=M_\Lambda B_\Lambda^pM_\Lambda$.  For either selection,
every $0\le\sigma<1/2$, and every $R$,
\begin{equation}\label{hnew:interior}
 \|(1-\Pi_\Lambda)K_\Lambda^p\|_{L^2\to H_D^\sigma}
                  =O_{p,\sigma,R}(\Lambda^{-R}).
\end{equation}
For the deep selection and $0\le\sigma<s<1/2$,
\begin{equation}\label{hnew:deep-raw-tail}
 \|(1-P_\Lambda)W_\Lambda\psi\|_{H_D^\sigma}
 \le C\Lambda^{\sigma-28(1-\sigma/s)}\|\psi\|_2.
\end{equation}
\end{proposition}
\begin{proof}
The change $r=2c\sqrt{\Lambda u}$ identifies the band kernel with
\[
 K_\nu(r,s)=\sqrt{rs}\int_0^1vJ_\nu(rv)J_\nu(sv)dv.
\]
Lommel's identity and its diagonal limit give
\[
 K_\nu(r,s)=\frac{\sqrt{rs}
 [sJ_{\nu+1}(s)J_\nu(r)-rJ_{\nu+1}(r)J_\nu(s)]}{s^2-r^2},
 \quad
 K_\nu(r,r)=\frac r2[J_\nu(r)^2-J_{\nu-1}(r)J_{\nu+1}(r)].
\]
The fixed-order asymptotic gives $K_\nu(r,r)=\pi^{-1}+O(r^{-1})$
for $r\ge1$.  The series at zero gives $K_\nu(r,r)=O(r^{2\nu+1})$.
The off-diagonal formula and Cauchy--Schwarz imply
\[
 |K_\nu(r,s)|\le C\min(1,|r-s|^{-1})\quad(r,s\ge1),\qquad
 |K_\nu(r,s)|\le Cr^{\nu+1/2}s^{-1}\quad(r<1,\ s\ge2).
\]
Integrating the diagonal to $Z=2c\Lambda$ proves \eqref{hnew:trace}.
Integrating $|K_\nu(r,s)|^2$ over $r<Z<s$ proves
\eqref{hnew:variance}; the large-variable integral is
$O(\int_0^Z(1+Z-r)^{-1}dr)$.
The scalar inequality
$|1_{x\ge\vartheta}-x|\le C_\vartheta x(1-x)$ proves the
fixed-threshold rank assertion.

Since $F_\Lambda^2=A_\Lambda$ and $F_\Lambda$ commutes with its
spectral selections, the two displayed formulas for $W_\Lambda$
agree.  Moreover
$\|W_\Lambda\psi\|_2^2=\langle\psi,A_\Lambda^{-1}\psi\rangle$.
For a band vector $w$, Cauchy--Schwarz in the Hankel integral gives
$|w(u)|\le C\sqrt\Lambda\|w\|_2$.  This proves the mass assertion
by choosing $d$ sufficiently small.  Termwise integration of the
Bessel series gives $w(u)=u^ab(u)$ with $b$ analytic near zero.
In particular $w'(u)=O(u^{a-1})$ when $a>0$.
When $a=0$, $w=b$ is analytic and $w'$ is bounded near zero.
For positive half-integral $a$, bounded ordinary derivatives of
every order at zero are neither claimed nor needed.

For the angular lift, the Fourier band radius is
$(c/\pi)\sqrt\Lambda$.  Plancherel therefore bounds its gradient
by $2c\sqrt\Lambda$ times its norm.  Since $u\partial_u$
corresponds to one half the radial dilation derivative, its norm
on $u<\Lambda$ is $O(\Lambda)\|w\|_2$.
The logarithmic profile tends to zero at its lower end and has
one jump at $L$.  The trace inequality
$|U\psi(L-)|^2\le2\|U\psi\|_2\|(U\psi)'\|_2$ and the Fourier
formula for that zero extension give \eqref{hnew:raw-mass-graph}
exactly for $s<1/2$.

The soft multipliers lift to multiplication by $p(|z|^2/\Lambda)$
and its Fourier counterpart; both preserve the angular sector.
The fixed gap between the soft and hard Fourier disks, and
Schwartz decay of the Fourier transform of the first multiplier,
give a Schur bound $O_R(\Lambda^{-R})$ on the hard band tail.
For a time-supported vector $v$,
\[
 \|(1-\Pi_\Lambda)v\|_2^2
 \le\begin{cases}
 (1-\vartheta)^{-1}\|(1-B_\Lambda)v\|_2^2,
       &\Pi_\Lambda=\Pi_\Lambda^{(\vartheta)},\\
 \Lambda^{56}\|(1-B_\Lambda)v\|_2^2,
       &\Pi_\Lambda=\Pi_\Lambda^{\mathrm d}.
 \end{cases}
\]
The arbitrary power absorbs this fixed deep loss.  Interpolating
with \eqref{hnew:raw-mass-graph} proves \eqref{hnew:interior}.
This is the precise order-$\nu$ replacement for
Propositions~\ref{prop:hankel-interior-approximation}
and~\ref{prop:hankel-high-interior}.

For the deep extension,
$\|(1-P_\Lambda)W_\Lambda\psi\|_2^2
\le\Lambda^{-56}(1-\Lambda^{-56})^{-1}\|\psi\|_2^2$.
The identity involving $F_\Lambda^{-1}$, dilation invariance,
and the uniformly bounded half-line restriction on $H^s$,
$s<1/2$, give a coarse $O(\Lambda^s)$ bound for this same tail.
Interpolation gives \eqref{hnew:deep-raw-tail}, the
order-$\nu$ version of \eqref{eq:high-graph-tail}.
The deep rank is $O(\Lambda)$ directly from its selected eigenvalue
bound and \eqref{hnew:trace}.  If its leading asymptotic is desired,
the positive soft-compression argument in
Proposition~\ref{prop:hankel-high-cutoff} applies using
\eqref{hnew:trace}--\eqref{hnew:variance} and the preceding Schur bound;
it gives rank asymptotic to $2c\Lambda/\pi$.
\end{proof}

\begin{lemma}[The order-dependent cross operator]\label{hnew:cross}
Let a raw column $V$ have logarithmic output supported in $t\ge-A$,
with all fixed logarithmic derivatives Hilbert--Schmidt.
For every fixed $j,R$,
\[
 \|(1-p(v/Y))\mathcal H_{c,\nu}V\|_{\mathcal S_2(L^2,H_D^j)}
 \le C_{j,R,A}Y^{-R}\sum_{r\le j+2R}\|\partial_t^rUV\|_2.
\]
Integer $j$ suffice to control the fractional output norms in use.
After endpoint scaling the opposite spectral radius is $Y=\Lambda^2$.
\end{lemma}
\begin{proof}
The Bessel equation gives the exact diagonalization
\begin{equation}\label{hnew:differential}
 \mathcal D_\nu=-\partial_u(u\partial_u)+\frac{\nu^2}{4u},
 \quad \mathcal H_{c,\nu}\mathcal D_\nu=c^2v\mathcal H_{c,\nu},
 \quad
 U\mathcal D_\nu U^{-1}
 =e^{-t}\bigl[-(\partial_t-1/2)^2+\nu^2/4\bigr].
\end{equation}
Its self-adjoint realization is supplied by the angular Fourier
model.  On the tail $v\gg Y$, insert $v^{-R}$ and transfer $v^R$
to $c^{-2R}\mathcal D_\nu^R$.  In logarithmic coordinates this
is $e^{-Rt}$ times a fixed differential polynomial of degree $2R$,
whose coefficients are bounded on the support $t\ge-A$.
Commuting logarithmic derivatives through the smooth tail cutoff
and using dilation reversal proves the estimate.
Approximation by compact cores justifies the diagonalization on
these columns.  Finally
$\mathcal H_{c,\nu}D_\Lambda=D_{1/\Lambda}\mathcal H_{c,\nu}$
turns the dual radius $\Lambda$ into $\Lambda^2$.
This proves the replacement for Lemma~\ref{lem:hi-hankel-cross},
including its polynomial-height form whenever the displayed
derivative budgets are polynomial in height.
\end{proof}

\subsection{Arithmetic inputs and the sharp Mellin passage}

\begin{lemma}[Fixed-form analytic estimates]\label{hnew:arithmetic}
Let $\mu_f$ be the Dirichlet inverse coefficients of $\lambda_f$.
For fixed $f$ there are $b_f>0$ and constants such that
\begin{align}
 |\lambda_f(n)|+|\mu_f(n)|&\le2d(n),\label{hnew:divisor}\\
 \sum_{n\le X}\mu_f(n)&\ll_f Xe^{-b_f\sqrt{\log(2+X)}},\label{hnew:mertens}\\
 |L_f(1/2+i\tau)|&\ll_f(2+|\tau|)^{3/8},
 \qquad \|E_f\psi\|_2\ll_f\|\psi\|_{H_D^{3/8}}.\label{hnew:graph}
\end{align}
For $A_f(X)=\sum_{n\le X}\lambda_f(n)$,
\begin{equation}\label{hnew:mean-square}
 \int_1^X|A_f(y)|^2dy\ll_f X^{3/2},\qquad
 \int_Y^\infty|A_f(y)|^2y^{-2}dy\ll_f Y^{-1/2}.
\end{equation}
The function $L_f$ has polynomial growth on each fixed vertical
strip, a uniform Euler-product lower bound at $2+iT$, and
$O_f(\log(2+|T|))$ zeros in any fixed-radius disk near height $T$.
\end{lemma}
\begin{proof}
For $k\ge2$ and $p\nmid N$, Deligne's primary statement
\citep[Theorem~8.2, p.~302]{deligne1974}
gives unit-modulus normalized Hecke roots.
At $p\mid N$, the newform Euler factor is linear.
The ramified coefficient formulas of
\citet[\S\S1.7--1.8, pp.~509--510]{deligneserre1974}
give $|\lambda_f(p)|\in\{0,1,p^{-1/2}\}$.
For $k=1$, the finite-image representation and the sharper
coefficient bound in
\citet[Theorem~9.1 and equation~(9.3), p.~527]{deligneserre1974}
give $|\lambda_f(n)|\le d(n)$, including the ramified factors.
Thus, in every case, the complete Euler product is
\[
 L_f(s)=\prod_p
     (1-\lambda_f(p)p^{-s}+\xi(p)p^{-2s})^{-1}.
\]
At good primes $|\lambda_f(p^r)|\le r+1$; at bad primes
$\lambda_f(p^r)=\lambda_f(p)^r$ and $|\lambda_f(p)|\le1$.
The inverse coefficients satisfy
$\mu_f(p)=-\lambda_f(p)$, $\mu_f(p^2)=\xi(p)$, and
$\mu_f(p^r)=0$ for $r\ge3$.
In particular $\mu_f(p^2)=0$ at every bad prime.
Multiplicativity proves \eqref{hnew:divisor}, for $f$ and $\bar f$.
The fixed-form Weyl bound of
\citet[Theorem~1.1, p.~2 of the author preprint]{booker2019}
is $O_f(|\tau|^{1/3}\log|\tau|)$, at arbitrary level and
including weight one.  The latter cases are also explicit in
the discussion of dihedral, tetrahedral, octahedral, and
icosahedral forms in footnote~1 on p.~4 of that preprint.
It implies \eqref{hnew:graph} by Plancherel.
Convexity with exponent $1/2+\epsilon$ would not suffice.

Hoffstein and Ramakrishnan's primary theorem
\citep[Theorem~C(3), p.~299]{hoffsteinram1995}
excludes zeros in
\[
 \sigma\ge1-\frac{c_2}{\log(\max(1,|t|)+M_f)}.
\]
For fixed $f$, shrink this to $1-b/\log(3+|t|)$.
On a disk of radius $O(1/\log(3+|t|))$ centered just to the
right of 1, an analytic logarithm exists.
Convexity in that shrinking strip bounds $\log|L_f|$ above by
$O_f(\log\log(3+|t|))$; the Euler logarithm at the center has
that bound in modulus.  Borel--Carath\'eodory on a smaller disk gives
$|L_f(s)^{-1}|\ll_f(\log(3+|t|))^C$ in a narrower region.
Bounded heights follow by compactness, without an exceptional zero.

The divisor hyperbola argument gives
$\sum_{X<n\le X+H}d(n)\ll H\log(2X)+\sqrt X$.
For a half-integer $X$, truncated Perron at $1+1/\log X$ and
height $2\le T\le\sqrt X$ therefore has error
$O(X\log^2X/T+\sqrt X)$: split the near range into
$|n-X|\le X/T$ and its successive doublings.
Moving the contour to $1-b/\log(3+T)$ uses the reciprocal bound.
Choose $T=\exp(d\sqrt{\log X})$ with fixed $d>0$ and absorb the
logarithms into a smaller positive exponent.  This proves
\eqref{hnew:mertens} for half-integers and hence all $X$.
No height-uniform twisted inverse-coefficient assertion is used.

Apply \citet[Theorem~1, p.~62, equation~(57), with setup on
pp.~43--44]{chandrasekharannar1964} to
\[
 \varphi(s)=c^{-s}L_f(s),\quad
 \psi(s)=\varepsilon_fc^{-s}L_{\bar f}(s),\quad
 \Delta(s)=\Gamma(s+a),\quad\delta=1,\quad A_{\rm CN}=1.
\]
The gamma shift, including $a=0$, and complex coefficient
sequences are allowed.  Both Dirichlet scales are $cn$, and
the completed function has the required strip decay.
The inequality $d(n)^2\le d_4(n)$ gives
$\sum_{n\le X}|\lambda_f(n)|^2\ll X\log^3(2X)$.
Their parameters $\beta=1,\beta'=3$ satisfy the included equality
case $2\beta-\delta-1/A_{\rm CN}=0$ and give $\omega_0=1/4$.
After undoing the fixed scale, their error differs from $A_f$
only by the residue constant $L_f(0)$.
The estimate is $O(X^{3/2})+O(X\log^5X)$ and proves the first
part of \eqref{hnew:mean-square}.  Partial integration gives the second.

Finally the functional equation and absolute Euler convergence give
the fixed-strip polynomial bound.  At $2+iT$ both Euler factors
and their reciprocals are bounded uniformly.
Jensen's formula in a larger fixed disk now gives the local zero
count.  Nontrivial zeros lie in $0<\Re s<1$: the Euler product,
the zero-free line at 1, and the functional equation exclude the
other points apart from the usual trivial zeros at $-a-\mathbb N_0$.
\end{proof}

\subsection{A weighted-tent identity and all zero jets}

\begin{proposition}[The required Riesz estimate]\label{hnew:riesz}
For $A_f(X)=\sum_{n\le X}\lambda_f(n)$,
\begin{equation}\label{hnew:riesz-bound}
 R_f(X):=\int_0^XA_f(y)dy=L_f(0)X+O_f(X^{3/4}).
\end{equation}
When $k=1$, $L_f(0)=0$, so the bound is $R_f(X)=O_f(X^{3/4})$.
In particular $A_f(X)=O_{f,\epsilon}(X^{3/8+\epsilon})$ for
every $\epsilon>0$.
\end{proposition}
\begin{proof}
Set
\[
 \phi_X(u)=u^a(X-u)_+,\qquad
 B_a(X)=\sum_{n\le X}\lambda_f(n)n^a(X-n).
\]
The angular lift of $\phi_X$ is $z^\nu(X-|z|^2)_+$ and is smooth
at zero.  Substituting $u=r^2$ and twice using
$\frac d{dr}(r^{\nu+1}J_{\nu+1}(br))
 =br^{\nu+1}J_\nu(br)$ gives Sonine's identity in the required
normalization:
\begin{equation}\label{hnew:sonine}
 (\mathcal H_{c,\nu}\phi_X)(v)
       =\frac{X^{a+1}}{cv}J_{\nu+2}(2c\sqrt{Xv}).
\end{equation}
The outer endpoint is a continuous tent and requires a limiting
argument.  Convolve its compact angular lift with a nonnegative
radial smooth approximate identity on $\mathbb R^2$.
This preserves angular frequency and gives smooth compact angular
profiles, to which \eqref{hnew:poisson} applies.
On the Fourier side it multiplies the transform by a radial
function of modulus at most one, converging pointwise to one.
Equation \eqref{hnew:sonine} and the fixed-order Bessel asymptotic
give the summable majorant $C_{f,X}d(n)n^{-5/4}$ for the dual sum.
Dominated convergence applies there.  The original supports stay
in a fixed compact disk, so their finite theta sums converge,
including the zero tent value at a possible boundary integer.
At logarithmic time zero the identity is therefore
\begin{equation}\label{hnew:weighted-tent}
 B_a(X)=\frac{\varepsilon_fX^{a+1}}c
       \sum_{n\ge1}\frac{\lambda_{\bar f}(n)}n
               J_{\nu+2}(2c\sqrt{Xn})
       =O_f(X^{a+3/4})\quad(X\ge1).
\end{equation}

If $a=0$, then $\nu=0$ and $B_0=R_f$.
Equation \eqref{hnew:weighted-tent} directly proves
$R_f(X)=O_f(X^{3/4})$; no Abel division by a positive weight is
needed.  Completed entireness of
$c^{-s}\Gamma(s)L_f(s)$ implies $L_f(0)=0$.
This proves \eqref{hnew:riesz-bound} in weight one.
The final pointwise argument below applies in this case as well.
For the intervening Abel calculation assume $a>0$.

Twice integrating by parts in the measure with second primitive
$B_a$, against $g(y)=(X-y)y^{-a}$, yields
\begin{equation}\label{hnew:abel-twice}
 \begin{split}
 R_f(X)={}&X^{-a}B_a(X)
  +a(a+1)X\int_0^X B_a(y)y^{-a-2}dy\\
 &+a(1-a)\int_0^X B_a(y)y^{-a-1}dy.
 \end{split}
\end{equation}
There is no boundary at zero because $B_a$ vanishes below 1;
at $X$, $g(X)=0$ and $g'(X)=-X^{-a}$.
The first integral has a finite limit at infinity, since its
integrand is $O(y^{-5/4})$.  Its omitted tail times $X$, and both
other terms, are $O(X^{3/4})$.  Thus
$R_f(X)=CX+O_f(X^{3/4})$.
For $\Re s>1$,
\[
 L_f(s)=s(s+1)\int_1^\infty R_f(y)y^{-s-2}dy
       =C(s+1)+s(s+1)\int_1^\infty(R_f(y)-Cy)y^{-s-2}dy.
\]
The last integral is holomorphic on $\Re s>-1/4$.
Continuing and evaluating at zero gives $C=L_f(0)$.

Finally average $A_f$ on $[X,X+H]$.
Equation \eqref{hnew:riesz-bound} bounds this average by
$O_f(1+X^{3/4}/H)$.  Its difference from $A_f(X)$ is at most
$O_{f,\epsilon}((H+1)X^\epsilon)$, since
$|\lambda_f(n)|\ll_\epsilon n^\epsilon$.
Taking $H=X^{3/8}$ proves the stated pointwise bound.
\end{proof}

For $k\ge2$, the term $L_f(0)X$ cannot in general be discarded.
It is a Riesz residue, not a pole of $L_f$, and creates no extra
theta coordinate.  The direct substitution $J_2\mapsto J_{\nu+2}$
in the order-zero unweighted tent formula would be invalid:
the weight $u^a$ in \eqref{hnew:sonine} is essential.

\begin{proposition}[Sharp Mellin identity and raw tail kernels]
\label{hnew:sharp-mellin}
Let $\phi$ be a selected sharp profile, or an angular Schwartz
profile.  Then for $0<\Re s<1$,
\begin{equation}\label{hnew:improper}
 \lim_{\eta\downarrow0}
 \int_\eta^\infty (E_f\phi)(\log x)x^{s-1/2}\frac{dx}{x}
                =L_f(s)\mathcal M\phi(s).
\end{equation}
The convergence is locally uniform and admits every fixed
$s$-derivative.  For a nontrivial zero $\rho=\beta+i\gamma$ and
$j<m_\rho$, put
\begin{equation}\label{hnew:jet-kernel}
 K_{\rho,j,\Lambda}(u)=
 \left.(\partial_s+L)^j\left[
 u^{s-1}\left(L_f(s)-\sum_{n\le\Lambda u}\lambda_f(n)n^{-s}\right)
                         \right]\right|_{s=\rho}.
\end{equation}
This is a raw $L^2$ kernel annihilating the Hankel band, and
\begin{equation}\label{hnew:jet-exterior}
 \|1_{u>\Lambda}K_{\rho,j,\Lambda}\|_2^2
                           \ll_{f,\rho,j}\Lambda^{-2\beta}.
\end{equation}
\end{proposition}
\begin{proof}
Integration by parts against \eqref{hnew:riesz-bound} proves,
for $\Re s=\beta>0$,
\begin{equation}\label{hnew:integrated-discrepancy}
 \int_X^\infty A_f(y)y^{-s-1}dy
 =\frac{L_f(0)}sX^{-s}+O_{f,s}(X^{-\beta-1/4}).
\end{equation}
The same estimate with fixed logarithmic powers proves local
uniformity of derivatives.
It continues $L_f(s)/s=\int_1^\infty A_f(y)y^{-s-1}dy$
from $\Re s>1$ to $\Re s>0$.
For the disk $1_{(0,b)}$ the theta function is
$\sqrt x A_f(b/x)$, so its omitted lower Mellin integral is
\[
 O_{f,s}(\eta^\beta)
      +O_{f,s}(\eta^{\beta+1/4}b^{-1/4}).
\]
For $b<\eta$, the complete disk integral gives the same majorant.
For a sharp selected profile use
\[
 \phi(u)=\phi(\Lambda-)1_{u<\Lambda}
                      -\int_0^\Lambda\phi'(b)1_{u<b}db.
\]
For $k\ge2$, Proposition~\ref{hnew:geometry} gives
$\phi'(b)=O(b^{a-1})$, where $a\ge1/2>1/4$.
For $k=1$, it gives the stronger $\phi'(b)=O(1)$.
In either case $|\phi'(b)|b^{-1/4}(1+|\log b|)^j$
is integrable at zero.  This proves passage to the limit and
every fixed derivative.  In weight one the $O(\eta^\beta)$
term can be omitted because $L_f(0)=0$; the same bound suffices.
For angular Schwartz profiles use the
same representation over the full positive axis; rapid decay
handles infinity.  This proves \eqref{hnew:improper}.

The kernel \eqref{hnew:jet-kernel} vanishes below $\Lambda^{-1}$,
because the indicated derivatives of $L_f$ vanish at $\rho$.
For $j=0$, Stieltjes summation and the preceding continuation give
\begin{equation}\label{hnew:tail-decomposition}
 K_{\rho,0,\Lambda}(u)
  =-\Lambda^{-\rho}u^{-1}A_f(\Lambda u)
    +\rho u^{\rho-1}\int_{\Lambda u}^\infty
                                A_f(y)y^{-\rho-1}dy.
\end{equation}
The endpoint term is in raw $L^2$ by \eqref{hnew:mean-square}.
Its squared exterior norm is
\[
 \Lambda^{1-2\beta}\int_{\Lambda^2}^\infty
                       |A_f(y)|^2y^{-2}dy\ll_f\Lambda^{-2\beta}.
\]
Equation \eqref{hnew:integrated-discrepancy} makes the second term
$L_f(0)\Lambda^{-\rho}u^{-1}
+O(\Lambda^{-\beta-1/4}u^{-5/4})$.
Its squared exterior norm is
$O(\Lambda^{-2\beta-1}+\Lambda^{-2\beta-2})$.
For $j\ge1$, $\partial_s+L$ cancels the endpoint factor
$\Lambda^{-s}$ and the displayed $L_f(0)$ term.
After the substitution $y=(\Lambda u)z$, the remaining derivatives
insert only fixed powers of $\log z$ in an integrable remainder.
This proves the global $L^2$ assertion and \eqref{hnew:jet-exterior}.

First test the kernels on inverse Hankel transforms of smooth
compact spectral profiles supported strictly inside $(0,\Lambda)$.
They are angular Schwartz profiles.  Equation \eqref{hnew:poisson}
gives their lower theta support gap, and \eqref{hnew:improper}
identifies the lower-tail integral and its jets.
These therefore vanish.  Such profiles are dense in the band;
the raw $L^2$ kernel extends annihilation by continuity.
There is no principal value-at-zero constraint.
\end{proof}

\subsection{Gram control and the deep endpoint inputs}

\begin{proposition}[Full Gram and deep leakage]\label{hnew:gram-deep}
For either selection,
\begin{equation}\label{hnew:gram-bound}
 \|G_\Lambda^{-1}\|\ll_f\Lambda^2\log^2(2+\Lambda).
\end{equation}
For the deep selection,
\begin{equation}\label{hnew:deep-leakage}
 \|(1-S_L)Z_\Lambda^{\mathrm d}\|_{\mathrm{HS}}
                       \ll_f\Lambda^{-17/8}\log(2+\Lambda).
\end{equation}
Every fixed complete window zero packet $\mathcal V_{\mathcal F,L}$
satisfies
\begin{equation}\label{hnew:deep-packets}
 \|Q_\Lambda^{\mathrm d}P_{\mathcal V_{\mathcal F,L}}\|_{\mathrm{HS}}^2
            \ll_{f,\mathcal F}\Lambda^{1-56}(1+L)^C.
\end{equation}
In particular this is smaller than every fixed negative power of $L$.
\end{proposition}
\begin{proof}
The coefficient convolution identity gives finite pointwise inversion
\[
 \psi(u)=\sum_{m\le\Lambda/u}
       \frac{\mu_f(m)}{\sqrt{mu}}(E_f\psi)(\log(mu)).
\]
On $u\ge d/\Lambda$, Minkowski and the change to logarithmic
measure bound its norm by
$\sum_{m\le C\Lambda^2}|\mu_f(m)|m^{-1/2}\|E_f\psi\|_2$.
Use \eqref{hnew:raw-mass-graph} and
$\sum_{m\le Y}d(m)m^{-1/2}\ll\sqrt Y\log(2Y)$.
The result is
$\|\psi\|_2\ll_f\Lambda\log(2+\Lambda)\|E_f\psi\|_2$,
which proves injectivity and \eqref{hnew:gram-bound}.

Time support makes the upper theta tail zero.
Since $\mathcal H_{c,\nu}W_\Lambda\psi=F_\Lambda^{-1}\psi$
is time supported, \eqref{hnew:poisson} makes the lower theta
tail of $W_\Lambda\psi$ zero.  Apply the graph bound to its raw
exterior difference.  Equation \eqref{hnew:deep-raw-tail} with
$\sigma=3/8$, $s=7/16$ has $1-\sigma/s=1/7$.
Summing columns and applying the full Gram bound gives
\[
 \|(1-S_L)Z_\Lambda^{\mathrm d}\|_{\mathrm{HS}}^2
 \ll_f
 \Lambda^{2+1+2(3/8)-56/7}\log^2(2+\Lambda)
 =\Lambda^{-17/4}\log^2(2+\Lambda).
\]
This is the direct substitution in the abstract tail calculation
of Corollary~\ref{cor:hankel-high-theta-tail}, retaining its
additional Gram logarithm.

For a jet in \eqref{hnew:jet-kernel}, band annihilation makes its
pairing with $\psi$ the negative of its exterior pairing with
$(1-P_\Lambda)W_\Lambda\psi$.
The latter has squared raw operator norm $O(\Lambda^{-56})$.
Equations \eqref{hnew:gram-bound} and \eqref{hnew:jet-exterior}
therefore give the squared normalized theta-dual bound
$O(\Lambda^{2-56-2\beta}(1+L)^C)$.
This is a single row functional on the full normalized column;
there is no extra factor equal to the selected rank.
The full Mellin jet vanishes at $\rho$, and the upper theta tail
is zero, so it controls the window jet.
For $\beta<1/2$, left normalization contributes
$\Lambda^{2\beta-1}$; for $\beta>1/2$, right normalization
contributes $\Lambda^{1-2\beta}$ and fixed polynomial factors in $L$.
The respective powers are $\Lambda^{1-56}$ and
$\Lambda^{3-56-4\beta}\le\Lambda^{1-56}$.
Critical polynomial normalization costs only powers of $L$.
The finite packet Gram calculation of
Theorem~\ref{thm:finite-packet}, including repeated zeros and
decoupling of the two endpoints, proves \eqref{hnew:deep-packets}.
\end{proof}

\begin{lemma}[The two one-sided closures]\label{hnew:closure}
The closure of $E_fC_c^\infty((0,1))$ in the negative-half-line
Hardy space is its inner multiple by the pure Blaschke product
of the zeros of $L_f$ in $\Re s>1/2$, with their orders.
The other closure is the reflection by $J$ of the closed
$E_{\bar f}$ core image.  Its model-space divisor is the left
divisor of $L_f$.  No singular inner factor or common exponential
delay is present.
\end{lemma}
\begin{proof}
The core images are square integrable by \eqref{hnew:graph} and
have nonpositive logarithmic support.  Raw contractions give
left translations, so the invariant-subspace theorem applies.
Their Mellin transforms are $L_f(s)\mathcal M\phi(s)$.
At each prescribed point a sufficiently narrow raw bump has
nonzero Mellin transform, so the common zero divisor is exactly
the indicated divisor of $L_f$.
Every generator continues analytically through each finite
boundary point.  An interval where one generator is nonzero
cannot support its singular inner measure; at an isolated
finite-order boundary zero an atom is likewise excluded.
Covering the boundary excludes a common singular measure.
A common exponential inner factor would force all images to
vanish on one fixed interval just below zero.
Take a bump arbitrarily close to raw $1$: when $e^t>1/2$,
only $n=1$ can contribute, and its coefficient is one.
This excludes a common delay.
Finally apply the same proof to $E_{\bar f}$ and reflect its
closure using \eqref{hnew:poisson}.
If $\rho=\beta+i\gamma$ is a left-half zero of $L_f$,
then $1-\rho=1-\beta-i\gamma$ is a right-half zero of
$L_{\bar f}$, with the same order.
The negative-half-line jet
$e^{(1/2-\beta+i\gamma)t}t^j$ for that zero becomes, under $J$,
the positive-half-line jet $e^{(\beta-1/2-i\gamma)t}(-t)^j$.
Thus its Cauchy center is the original $\gamma$, as required.
This uses the dual closure and does not assert that the
finite image projection for $f$ commutes with $J$ or with
$F_\Lambda$.  This proves the holomorphic specialization
of the closure argument in Lemma~\ref{lem:hecke-halfline-closure};
no Hecke-character closure theorem is assumed for $f$.
\end{proof}

\begin{proposition}[Deep local equality and polynomial mass]
\label{hnew:deep-upper}
The deep measure converges vaguely to $\mathfrak b_f$, and
there are fixed $C,A>0$ such that
\begin{equation}\label{hnew:uniform-mass}
 \mu_\Lambda^{\mathrm d}([T,T+1])
             \le C(2+|T|)^A\qquad(\Lambda\ge2,\ T\in\mathbb R).
\end{equation}
It therefore converges on every Schwartz frequency test.
\end{proposition}
\begin{proof}
We verify the hypotheses of
Theorem~\ref{thm:endpoint-local-criterion} explicitly and then
identify which parts of the polynomial-height construction are
unchanged.  Fix $\eta\in C_c^\infty(\mathbb R)$ and set
$F_f(z)=L_f(1/2+iz)$.  Choose a compact frequency interval
containing $\supp\eta$, with endpoints away from critical zeros.
Let $P_{0,L}$ project onto all critical-line window jets in that
interval, with full multiplicities, and put $W_L=S_L-P_{0,L}$.
Give the finitely many real poles of $b=\eta/F_f$ the
negative-time prescriptions $(\tau-\gamma-i0)^{-j}$.
With $T_L=U^{-1}C_bW_L$ one has
\begin{equation}\label{hnew:inverse-column}
 E_fT_L=C_\eta W_L,\qquad
 \|UT_L\|_{\mathcal S_2(L^2,H^j)}\le C_j(1+L)^{d_j}.
\end{equation}
Indeed subtract the compactly cut off principal parts of $b$.
Its inverse transform $\kappa$ has a finite exponential-polynomial
negative tail and a Schwartz remainder, and is Schwartz at
positive infinity.  The input zero moments annihilate the whole
exponential polynomials outside the input interval.
The exterior quadrant kernels are Hilbert--Schmidt, and the
middle kernel has polynomial growth in $L$.
These are the one-dimensional calculations in
\eqref{eq:hi-fixed-bulk}; they involve $F_f$, not the Hankel order.
The exact multiplier identity proves the first assertion of
\eqref{hnew:inverse-column}.

The dual column has these same bounds with $\bar f$: by
\eqref{hnew:mellin-hankel}--\eqref{hnew:poisson},
\begin{equation}\label{hnew:dual-reciprocal}
 \frac{g_{c,\nu}(s)}{L_f(1-s)}
                        =\frac{\varepsilon_{\bar f}}{L_{\bar f}(s)}.
\end{equation}
Reflecting a real-pole prescription and restoring the chosen one
adds only delta derivatives, annihilated by the reflected moments.
Reflection sends the moment conditions precisely to the packet
of the dual reciprocal; both $f$ and $\bar f$ obey
Lemma~\ref{hnew:arithmetic}.  Thus no estimate for a large gamma
quotient is being substituted for this cancellation, and no
Gram-norm symmetry between the two theta images is assumed.

After upper-endpoint translation, the first raw boundary column
has kernel
\[
 (1-p(e^t))\kappa(t-s)1_{[-2L,0]}(s)
\]
before its finite input-jet correction.  Its output is bounded
below and its input nonpositive, so the positive Schwartz tail
of $\kappa$ gives a fixed Hilbert--Schmidt limit, with every
fixed derivative and output polynomial weight.
Lemma~\ref{lem:critical-window-packet} gives strong convergence
to zero of the translated input projections.  Right composition
by these projections therefore removes their corrections in
Hilbert--Schmidt norm.  The dual endpoint follows from
\eqref{hnew:dual-reciprocal}.  Lemma~\ref{hnew:cross} controls both
opposite-endpoint errors with arbitrary powers of $\Lambda^{-1}$.

Here is the exact cutoff algebra.  With $M=M_\Lambda$,
$B=B_\Lambda^p$, and $P_+=1_{t\le L}$,
\begin{align}\label{hnew:cutoff-algebra}
 S_LE_fT_L-E_fMBMT_L
   ={}&S_LE_f(1-M)T_L+S_LE_f(1-B)MT_L\\
      &+P_+E_f(1-M)BMT_L.\nonumber
\end{align}
The term $(P_+-S_L)E_fBMT_L$ is zero by the
band support gap, and $E_fMBMT_L$ has no upper tail.
The cross estimates replace $MT_L$ by $T_L$ in the second
term and $BMT_L$ by $MT_L$ in the last, with a vanishing
Hilbert--Schmidt error after \eqref{hnew:graph}.
Including the two outside-input quadrants of
$S_LC_\eta(1-S_L)$, whose kernel is Schwartz, gives
\begin{equation}\label{hnew:endpoint-decomposition}
 S_LC_\eta(1-P_{0,L})
   =\mathcal F_L+\tau_L\mathcal B_+\tau_{-L}
           +\tau_{-L}\mathcal B_-\tau_L+o_{\mathcal S_2}(1).
\end{equation}
Here $\mathcal B_+$ and $\mathcal B_-$ are fixed Hilbert--Schmidt
operators with ranges in the negative and positive half-lines.
The bulk is $\mathcal F_L=E_fK_\Lambda^pT_L$.
Equation \eqref{hnew:interior}, with
\eqref{hnew:inverse-column}, gives
$\|(1-Q_\Lambda^{\mathrm d})\mathcal F_L\|_2\to0$ directly
in the theta graph norm.  Also $\|\mathcal F_L\|_2$ is polynomial
in $L$: this follows from \eqref{hnew:cutoff-algebra}, its bounded
boundary terms, and the window convolution norm, rather than
from a crude raw bound with a power of $\Lambda$.

For the closure spaces in Lemma~\ref{hnew:closure}, each fixed
core profile, dilated to an endpoint, is approximated in theta
graph norm by the deep selection.  To see this, its angular
Hankel transform is Schwartz; dilation therefore makes the
opposite raw tail arbitrarily small, and
\eqref{hnew:interior} completes the approximation.
Contractivity and density extend it to every fixed vector in
the closure.  Finite exponential-jet packets are dense in the
two Blaschke complements.  Equation \eqref{hnew:deep-packets}
and exponential decay beyond the opposite endpoint give the
required $O_R(L^{-R})$ annihilation on each fixed such packet.
The finite-packet diagonal calculation of
Theorem~\ref{thm:finite-packet} is
\[
 \|P_{K_+}C_\eta\|_2^2+\|P_{K_-}C_\eta\|_2^2
 =\sum_{\beta\ne1/2}m_\rho
 \int|\eta(\tau)|^2
 \frac{|\beta-1/2|\,d\tau}
 {\pi((\tau-\gamma)^2+|\beta-1/2|^2)}.
\]
The local zero count makes these Hilbert--Schmidt norms finite
and justifies exhaustion by finite packets.
Lemma~\ref{lem:critical-window-packet} supplies the critical
contribution $\sum_{\beta=1/2}m_\rho|\eta(\gamma)|^2$.
In particular $C_\eta$ has not been commuted with $P_{0,L}$:
that lemma proves its approximate scalar action on each jet.
All four hypotheses of
Theorem~\ref{thm:endpoint-local-criterion} are now verified.
It gives
\[
 \|(1-Q_\Lambda^{\mathrm d})S_LC_\eta\|_2^2
  =\operatorname{Tr}(D_\Lambda^{\mathrm d}C_\eta C_\eta^*)
       \longrightarrow\int|\eta|^2d\mathfrak b_f.
\]
Polarization, compact mass bounds, and uniform approximation
give vague convergence.

For the uniform estimate, let $H=2+|T|$.
The analytic construction of Lemma~\ref{lem:hi-rectangle}
uses only the fixed-strip bound, Euler-center lower bound,
and $O(\log H)$ local zero count, all proved in
Lemma~\ref{hnew:arithmetic}.  It therefore supplies a bump
$\eta_T=1$ on $[T,T+1]$, a packet of $N_T=O(\log H)$ zeros,
and the proper rational sums of the principal parts of
$z^j/F_f(z)$, with the bounds
\[
 \|\partial_\tau^m(\tau^j\eta_T/F_f-R_{T,j})\|_2
       \le H^{C_{j,m}},\qquad
 \sup_{-A\le t\le0}|R_{T,j,-}(t)-R_{T,j,+}(t)|
       \le H^{C_{j,A}}.
\]
For precision, this construction divides by a disk Blaschke
product and uses Harnack to control the zero-free quotient.
Vertical sides are chosen by averaging logarithmic distances
from zero projections; the horizontal sides have imaginary
parts $\pm3/4$.  Cauchy integrals bound the combined principal
parts, without division by root gaps.  These scalar analytic
steps are independent of the underlying raw transform.

Let $W_{L,T}$ remove the complete packet of window jets, including
off-line zeros.  The one-dimensional division and combined-residue
argument of Lemma~\ref{lem:hi-allpole-boundaries} then gives
\[
 \|U^{-1}C_{\eta_T/F_f}W_{L,T}\|_{\mathcal S_2(L^2,H_D^j)}
       \le H^{C_j}(1+L)^{N_T+C_j},
\]
and polynomially bounded fixed translated endpoint columns.
Each inverse linear factor has a decaying integration kernel
of modulus at most one on a moment-zero input; its cost is
at most $2L$.  At an endpoint, subtract the whole opposite
residue polynomial, which is annihilated exactly.  What remains
is supported in a fixed triangle and is controlled by the
displayed residue-difference bound.
For the dual endpoint use \eqref{hnew:dual-reciprocal}, and for
its raw cross errors use Lemma~\ref{hnew:cross}.
These are the only transform-dependent steps.

Applying \eqref{hnew:interior} to these columns produces a bulk
error of the form
\begin{equation}\label{hnew:height-error}
 H^C(1+L)^{C\log H+C}e^{-\eta_0L},
\end{equation}
with fixed $\eta_0>0$.  The derivative budgets are fixed,
not proportional to $N_T$, and $M=56$ is unchanged.
For $L\ge H^D$, with a fixed sufficiently large $D$, the error
is bounded by one.  Contractivity bounds the boundary residual
by a fixed power of $H$.  The complementary input projection
has rank at most $N_T$ and contributes at most that rank.
Orthogonality of the two input subspaces adds their squared
Hilbert--Schmidt norms exactly.  The distance identity above
therefore gives \eqref{hnew:uniform-mass} for $L\ge H^D$.
For $L<H^D$ use $0\le d\mu_\Lambda^{\mathrm d}/d\tau\le L/\pi$.
Bounded heights follow from the compact limit and the elementary
bound on the remaining bounded cutoff range.
This is the explicit specialization of the argument of
Theorem~\ref{thm:hi-uniform-mass}; no raw strong-convergence
claim replaces an image-distance estimate.
Summing a Schwartz supremum over unit intervals against
\eqref{hnew:uniform-mass} proves Schwartz convergence.
\end{proof}

\subsection{The fixed-threshold functional and the final transfer}

\begin{proposition}[An exterior functional in the theta norm]
\label{hnew:natural-functional}
For the fixed threshold, every $R>0$, and
$q\in L^2((\Lambda,\infty),du)$ extended by zero,
\begin{equation}\label{hnew:exterior-functional}
 \sup_{0\ne\psi\in\operatorname{Ran}\Pi_\Lambda}
  \frac{|\langle W_\Lambda\psi,q\rangle|^2}{\|E_f\psi\|_2^2}
       \le C_{f,R}\Lambda L^{-R}\|q\|_2^2
       \qquad(\supp q\subset(\Lambda,\infty)).
\end{equation}
Equivalently, after relabelling $R$,
\[
 \|(1-P_\Lambda)W_\Lambda j_\Lambda G_\Lambda^{-1/2}\|
                       \le C_{f,R}\Lambda^{1/2}L^{-R}.
\]
\end{proposition}
\begin{proof}
Put $k_q=F_\Lambda^{-1}\Pi_\Lambda P_\Lambda\mathcal H_{c,\nu}q$.
Then $\langle W_\Lambda\psi,q\rangle=\langle\psi,k_q\rangle$ and
$\|k_q\|_2\le\vartheta^{-1/2}\|q\|_2$.
For $J_u(v)=1_{(0,\Lambda)}(v)J_\nu(2c\sqrt{uv})$, evaluation
using $F_\Lambda^{-1}\Pi_\Lambda=F_\Lambda A_\Lambda^{-1}\Pi_\Lambda$
and the exact identity
$P_\Lambda\mathcal H(1-P_\Lambda)\mathcal H P_\Lambda=P_\Lambda-A_\Lambda$
give
\begin{align}\label{hnew:row-energy}
 |k_q(u)|^2
 &\le c^2\|q\|_2^2
 \langle J_u,A_\Lambda^{-2}\Pi_\Lambda(1-A_\Lambda)J_u\rangle\\
 &\le\vartheta^{-2}\|q\|_2^2
                 \int_\Lambda^\infty|B_\Lambda(v,u)|^2dv.\nonumber
\end{align}
Commutation here is only that of spectral functions of $A_\Lambda$.
Replacing $J_u$ by $u\partial_uJ_u$ gives the differentiated bound.
The raw Lommel kernel is
\[
 B_\Lambda(v,u)=\frac{c\sqrt\Lambda}{v-u}
 \left[\sqrt v J_{\nu+1}(2c\sqrt{\Lambda v})J_\nu(2c\sqrt{\Lambda u})
 -\sqrt u J_{\nu+1}(2c\sqrt{\Lambda u})J_\nu(2c\sqrt{\Lambda v})\right].
\]
For $v>\Lambda$, $u\le\Lambda/2$, its modulus is at most
$C\Lambda^{1/4}v^{-3/4}(1+\Lambda u)^{-1/4}$.
Using $zJ_\nu'(z)=\nu J_\nu(z)-zJ_{\nu+1}(z)$ changes the last
exponent to $1/4$ after $u\partial_u$.
The integral of $\Lambda^{1/2}v^{-3/2}$ over $v>\Lambda$ is bounded.
Thus \eqref{hnew:row-energy} proves
\begin{equation}\label{hnew:row-bounds}
 |k_q(u)|\ll\|q\|_2(1+\Lambda u)^{-1/4},\qquad
 |u k_q'(u)|\ll\|q\|_2(1+\Lambda u)^{1/4}.
\end{equation}
These estimates, including the exact energy step, replace
Lemma~\ref{nh:row}; a bounded raw inverse alone would not imply them.

Finite coefficient inversion gives the theta-dual representative
\[
 \beta_k(v)=\sqrt v\sum_{m\ge1}\frac{\overline{\mu_f(m)}}m k(v/m),
       \quad 0<v<\Lambda,
\]
extended by zero above $\Lambda$.
Separate $m=1$, whose squared $L^2(dv/v)$ norm is $\|k\|_2^2$.
Put $X=\Lambda v$, and for large $X$ split the remaining terms
at $Y=X/(\log X)^D$.  The small part is bounded by
\[
 C\|q\|_2X^{-1/4}\sum_{m\le Y}d(m)m^{-3/4}
       \ll\|q\|_2(\log X)^{1-D/4}.
\]
The extra logarithm is retained and absorbed by choosing $D$
sufficiently large.  For the rest use \eqref{hnew:mertens} and
Abel summation with
\[
 \left|\frac d{dy}\{y^{-1}k(v/y)\}\right|
       \ll\|q\|_2y^{-2}(1+X/y)^{1/4}.
\]
The resulting integral is bounded by
$C\|q\|_2\int_Y^\infty e^{-b_f\sqrt{\log(2+y)}}
(1+X/y)^{1/4}dy/y$, and is smaller than every fixed
negative power of $\log X$.
For bounded $X$ it is bounded.
Consequently
$|\beta_{k,\ge2}(v)|\ll_R\|q\|_2\sqrt v
[\log(2+\Lambda v)]^{-R}$.
Integrating first over $v\le\sqrt\Lambda$ and then over its
complement gives
\[
 \|\beta_k\|_{L^2(dv/v)}^2\ll_R\Lambda L^{-R}\|q\|_2^2.
\]

To justify the dual identity, first integrate
$\langle\psi,k\rangle$ over $u\ge\epsilon$.
Finite inversion and $v=mu$ give the same expression with the
sum in $\beta_k$ truncated at $m\le v/\epsilon$.
It converges uniformly on each $v\ge\delta>0$ by Abel summation.
Near zero, if $a>0$, angular regularity gives
$k(v/m)=O_\Lambda((v/m)^a)$, so the truncated representatives
are bounded by $C_\Lambda v^{a+1/2}\sum_m d(m)m^{-1-a}$.
For $a=0$ this absolute series would diverge and is not used.
Instead $k(v/m)=k(0)+O_\Lambda(v/m)$.
Abel summation in \eqref{hnew:mertens} gives
\[
 \sup_{M\ge1}\left|\sum_{m\le M}\frac{\mu_f(m)}m\right|<\infty,
 \qquad \sum_{m\ge1}\frac{d(m)}{m^2}<\infty.
\]
The truncated representatives are therefore
$O_\Lambda(\sqrt v)$ in weight one, uniformly in their
truncation.  Both bounds give integrable squared
$L^2(dv/v)$ majorants near zero.
Thus they converge in that space and
\[
 \langle\psi,k\rangle
   =\langle (E_f\psi)(\log v),\beta_k(v)\rangle_{L^2(dv/v)}.
\]
Cauchy--Schwarz proves \eqref{hnew:exterior-functional}.
\end{proof}

\begin{proposition}[Natural packets and smoothed leakage]
\label{hnew:natural-transfer-inputs}
For every fixed complete zero packet and every $R$,
\begin{equation}\label{hnew:natural-packets}
 \|Q_\Lambda^{(\vartheta)}P_{\mathcal V_{\mathcal F,L}}\|_{\mathrm{HS}}^2
                                             =O_{f,\mathcal F,R}(L^{-R}).
\end{equation}
For $g\in C_c^\infty(\mathbb R)+\mathbb C(\tfrac12e^{-|t|})$,
\begin{equation}\label{hnew:smoothed-leakage}
 \|V(g)(1-S_L)Z_\Lambda^{(\vartheta)}\|_{\mathrm{HS}}
                                                =O_{f,g,R}(L^{-R}).
\end{equation}
For every $h\in C_c^\infty(\mathbb R)$,
\begin{equation}\label{hnew:trace-transfer}
 \operatorname{Tr}((Q_\Lambda-S_LQ_\Lambda S_L)V(h))\longrightarrow0.
\end{equation}
\end{proposition}
\begin{proof}
For a jet kernel from \eqref{hnew:jet-kernel}, annihilation of the
band replaces its pairing with $\psi$ by its exterior pairing
with $W_\Lambda\psi$.
Equations \eqref{hnew:exterior-functional} and
\eqref{hnew:jet-exterior} give the squared theta-dual bound
$O(\Lambda^{1-2\beta}L^{-R})$.
The full Mellin jet vanishes and the upper theta tail is zero.
Left-endpoint normalization contributes $\Lambda^{2\beta-1}$;
right-endpoint normalization contributes $\Lambda^{1-2\beta}$,
giving $\Lambda^{2-4\beta}\le1$ for $\beta>1/2$.
Central jets cost only powers of $L$.
The finite Gram calculation of Theorem~\ref{thm:finite-packet}
therefore proves \eqref{hnew:natural-packets}.

For completeness the transform-independent raw kernel used in
Lemma~\ref{nh:raw-tail} is, in output coordinate $r=t+L$,
\[
 K_g(r,u)=u^{-1/2}\int_{\Lambda u}^\infty
   y^{-1/2}g(r+\log(y/(\Lambda u)))\,dA_f(y),\qquad u>\Lambda.
\]
Integration by parts gives an endpoint
$-\Lambda^{-1/2}u^{-1}A_f(\Lambda u)g(r)$ and the integral
\[
 -\Lambda^{-1/2}u^{-1}\int_1^\infty A_f(\Lambda uz)z^{-3/2}
                    (g'-g/2)(r+\log z)dz.
\]
Equation \eqref{hnew:mean-square} bounds their combined
Hilbert--Schmidt norm by
\begin{equation}\label{hnew:raw-smoothed-tail}
 \|V(g)1_{t<-L}E_f(1-P_\Lambda)\|_{\mathrm{HS}}
       \le C_f\Lambda^{-1/2}(\|g\|_2+\|g'\|_2).
\end{equation}
The Minkowski integral is a constant times
$\int_1^\infty z^{-5/4}dz$.
First use compact smooth $g$, then approximate in $L^1\cap H^1$.
Approximation on the raw graph core identifies the kernel
extension with the displayed theta operator.

The upper tail of $E_f\psi$ and the lower tail of $E_fW_\Lambda\psi$
are zero.  Thus the normalized lower tail in
\eqref{hnew:smoothed-leakage} is the composition of
\eqref{hnew:raw-smoothed-tail} with
$(1-P_\Lambda)W_\Lambda j_\Lambda G_\Lambda^{-1/2}$.
Its factors $\Lambda^{-1/2}$ and $\Lambda^{1/2}$ cancel, proving
the asserted arbitrary logarithmic rate.
The Green kernel $\kappa(t)=\tfrac12e^{-|t|}$ lies in
$L^1\cap H^1$, so it is covered.

The purely operator-theoretic proof of
Proposition~\ref{nh:trace-transfer} now applies with these verified
leakage hypotheses.  Explicitly, put $b=(1-\partial_t^2)h$, so
$h=b*\kappa$, and $P_o=1-S_L$.
The isometry of $Z_\Lambda$ gives
$\|V(g)S_LZ_\Lambda\|_{\mathrm{HS}}\le\sqrt{2L}\|g\|_2$.
Cyclicity and $Q-S_LQS_L=P_oQ+S_LQP_o$ yield
\begin{align*}
 |\operatorname{Tr}((Q-S_LQS_L)V(h))|
 &\le\|V(b^*)Z_\Lambda\|_2\|V(\kappa)P_oZ_\Lambda\|_2\\
 &\quad+\|V(b^*)P_oZ_\Lambda\|_2
                         \|V(\kappa)S_LZ_\Lambda\|_2\\
 &\ll_h L^{1/2-R}+L^{-2R}.
\end{align*}
The norms in this last display are Hilbert--Schmidt norms.
This proves \eqref{hnew:trace-transfer} without a rank loss.
\end{proof}

\begin{proof}[Proof of Theorem~\ref{hnew:theorem}]
The rank and Gram assertions are
Propositions~\ref{hnew:geometry} and~\ref{hnew:gram-deep}.
Proposition~\ref{hnew:deep-upper} gives the full deep measure limit
and its uniform polynomial mass bound.  Its uncompressed trace
follows already from
\[
 \|Q_\Lambda^{\mathrm d}-S_LQ_\Lambda^{\mathrm d}S_L\|_1
 \le2\sqrt{\operatorname{rank}Q_\Lambda^{\mathrm d}}\,
       \|(1-S_L)Z_\Lambda^{\mathrm d}\|_2
       =O_f(\Lambda^{-13/8}\log(2+\Lambda)).
\]

For the fixed threshold and a fixed finite zero packet, put
$P_V=P_{\mathcal V_{\mathcal F,L}}$.
Positivity, $0\le D_\Lambda\le S_L$, and
\eqref{hnew:natural-packets} give
\[
 \|(1-D_\Lambda^{1/2})P_V\|_{\mathrm{HS}}^2
       \le\operatorname{Tr}(P_V(1-D_\Lambda)P_V)
       =\|Q_\Lambda P_V\|_{\mathrm{HS}}^2=o(1).
\]
For a nonnegative compact frequency test $\varphi$, it follows that
\[
 \operatorname{Tr}(D_\Lambda C_\varphi)
 \ge\|P_VD_\Lambda^{1/2}C_\varphi^{1/2}\|_{\mathrm{HS}}^2
 =\operatorname{Tr}(P_VC_\varphi)+o(1).
\]
Theorem~\ref{thm:finite-packet} identifies the last limit.
Increasing the finite packet gives the all-zero lower bound;
the local zero count gives the necessary Cauchy-tail summability.
For large $\Lambda$ the deep raw space is contained in the
fixed-threshold space, hence their theta-image projections are
nested and
\[
 0\le D_\Lambda^{(\vartheta)}\le D_\Lambda^{\mathrm d}.
\]
The deep local limit gives the matching upper bound and its
polynomial local mass passes to the fixed threshold.
Compact convergence and this mass bound imply Schwartz convergence.
This is the positive sandwich and exhaustion argument of
Theorem~\ref{nh:theorem}, now with its lower-bound and upper-bound
inputs explicitly established for $f$.
The compressed Fourier trace identity and
\eqref{hnew:trace-transfer} prove \eqref{hnew:trace-theorem}.

The weighted zero count
$\sum_\rho m_\rho/(1+\gamma^2)<\infty$ follows from the local
zero bound.  The functional equation and
$L_{\bar f}(s)=\overline{L_f(\overline s)}$
reflect $\rho$ to $1-\overline\rho$ in the same $L_f$ divisor,
with the same order and ordinate.  Reality of the coefficients
is not assumed.
The assumed zero-free half-plane therefore confines the divisor
to $|\beta-1/2|\le\delta$.
Theorem~\ref{thm:band-bound}, whose hypothesis is precisely this
band and weighted count, proves the final inequality.
This last deduction changes neither its constants nor its
test-function norm.
\end{proof}

\begin{corollary}[Odd two-dimensional Artin functions]
\label{hnew:odd-artin}
Let $\rho:G_{\mathbb Q}\to GL_2(\mathbb C)$ be a continuous
irreducible odd Artin representation, with Artin conductor $N$.
Let $\lambda_\rho(n)$ be the Dirichlet coefficients of its
complete finite-prime Artin Euler product.
Use the raw Hankel transform
$\mathcal H_{2\pi/\sqrt N,0}$ and the theta operator
\[
 (E_\rho\psi)(t)
       =e^{t/2}\sum_{n\ge1}\lambda_\rho(n)\psi(ne^t).
\]
Then all the fixed-threshold and deep measure, trace, uniform-mass,
and half-plane assertions of Theorem~\ref{hnew:theorem} hold
for $L(s,\rho)$, with leading term $2\log\Lambda$ and no polar
coordinates.  This includes nonsolvable projective image $A_5$.
\end{corollary}
\begin{proof}
The classical modularity theorem of
\citet[Theorem~10.1(ii) and Corollary~10.2(ii), author-preprint
pp.~20--21]{kharewintenberger2009}
associates $\rho$ to a primitive holomorphic newform $f$ of
weight one.  By
\citet[Theorem~4.6, p.~514; proof pp.~515--516]{deligneserre1974},
its level is exactly the Artin conductor and
\begin{equation}\label{hnew:artin-all-primes}
 L(s,\rho)=L_f(s)
 =\prod_p\det(1-\rho(\operatorname{Frob}_p)p^{-s}
                       \mid V_\rho^{I_p})^{-1},
\end{equation}
including every ramified prime.
The representations are identified by their good Frobenius traces
and semisimplicity; this permits the full conductor and local-factor
statement, not merely an equality after deleting bad primes.
The same source's Corollary~4.8 gives entireness.

Oddness gives eigenvalues $1,-1$ at complex conjugation, so the
archimedean factor is
$\Gamma_{\mathbb R}(s)\Gamma_{\mathbb R}(s+1)
=2(2\pi)^{-s}\Gamma(s)$.
It matches the weight-one normalization up to the irrelevant
constant 2, and fixes $c=2\pi/\sqrt N$, $\nu=0$.
Equality \eqref{hnew:artin-all-primes} identifies every Dirichlet
coefficient.  Therefore the raw concentration operator, selected
space, theta map, full Gram matrix, and image projection are
literally those of $f$.  Theorem~\ref{hnew:theorem} applies.
The dual representation supplies the opposite boundary through
the functional equation; it is not added as a second theta channel.
The Artin--modular identification is classical.  The corollary
identifies its finite-cutoff realization.
\end{proof}

The theorem fixes the weight, level, and nebentypus.  Its constants
need not be uniform as the form varies.  The deep rank is asymptotic
to $4\Lambda/\sqrt N$; the logarithmic rank remainder applies to a
fixed threshold.  Even two-dimensional Artin representations and
general higher-degree functions require further cutoff estimates.


\section{Multiplicity, positive traces, and finite compression moments}
\label{sec:extensions-multiplicity}
Positivity of a scalar zero distribution and faithfulness of a jet
representation impose different conditions, and only the latter
excludes nilpotent zero blocks.

\begin{proposition}[Multiplicity and the jet radical]
\label{prop:multiplicity-positive}
Let real ordinates $\gamma$ have positive integer multiplicities
$m_\gamma$, with $\sum m_\gamma(1+\gamma^2)^{-1}<\infty$.
Then the scalar distribution
\[
 f*f^*\longmapsto\sum_\gamma
                 m_\gamma|\widehat f(\gamma)|^2,
 \qquad f\in C_c^\infty(\R),
\]
is nonnegative, without a simplicity hypothesis. It has a unitary
diagonal realization with $m_\gamma$ copies of each ordinate.
For $m>1$, however, the algebra $\C[\varepsilon]/(\varepsilon^m)$
with $\varepsilon^*=\varepsilon$ has no faithful bounded Hilbert
$*$-representation.
\end{proposition}

\begin{proof}
The series converges by Fourier decay and is a sum of nonnegative
terms. On $\bigoplus_\gamma\C^{m_\gamma}$ let $U(t)$ act by
$e^{i\gamma t}$ on each summand. Integration against a test function
is diagonal, and its trace is the desired multiplicity-weighted sum;
absolute summability proves that it is trace class. In a Hilbert
$*$-representation the image of $\varepsilon$ would be a bounded
self-adjoint nilpotent. The spectral theorem forces it to be zero,
so the representation cannot be faithful. Nevertheless
$a\mapsto m\,a(0)$ is a positive functional on the jet algebra:
it evaluates $a^*a$ as $m|a(0)|^2$ and annihilates its radical.
\end{proof}

Consequently repeated Dedekind zeros do not obstruct scalar Weil
positivity under GRH. Passing to individual irreducible Artin
representations may separate arithmetic factors, but multiplicity
does not force that passage. Corollary~\ref{hnew:odd-artin} covers
all odd irreducible degree-two Artin functions. The following
specialization identifies directly the cutoff inherited from the
fixed-field Hecke theorem.

\subsection{An Artin transfer from the Eisenstein field}
\label{sec:artin-induced-cutoff}
The following is a corollary of the fixed-field Hecke theorem.
Artin--Hecke reciprocity, induction of $L$-functions, and the
conductor formula are classical; no additional analytic estimate
for a nonabelian representation is claimed.

Let $K=\Q(\sqrt{-3})$, let $G_F=\operatorname{Gal}(\overline{\Q}/F)$,
and let $\sigma$ be the nontrivial automorphism of $K$.
For a continuous finite-order character $\chi$ of $G_K$, also
write $\chi$ for its primitive Hecke character, using arithmetic
Frobenius so that $\chi(\operatorname{Frob}_{\mathfrak p})
=\chi(\mathfrak p)$ at an unramified prime.
This correspondence and its local-factor compatibility are stated in
\citet[Theorems~4.2--4.3 and Proposition~4.4, pp.~52--53;
Proposition~5.2, p.~56]{horawaLfunctions}.
Extend ideal values by zero at the conductor primes.
Choose complex conjugation $j$ lifting $\sigma$, put
$\chi^\sigma(g)=\chi(jgj^{-1})$, and suppose
\[
 \chi\ne\chi^\sigma,\qquad
 \rho=\operatorname{Ind}_{G_K}^{G_\Q}\chi,\qquad
 q_\rho=3N\mathfrak f_\chi,\qquad c_\rho=2\pi/\sqrt{q_\rho}.
\]
In particular $\chi$ is nonprincipal.

On $L^2((0,\infty),du)$ use the Hankel transform
$\mathcal H_{c_\rho}$ defined above and, for fixed $0<\nu<1$, set
\[
 \Pi_{\rho,\Lambda}^{(\nu)}
  =\mathbf1_{[\nu,1]}
   (P_\Lambda\mathcal H_{c_\rho}P_\Lambda
                     \mathcal H_{c_\rho}P_\Lambda),
 \qquad P_\Lambda=\mathbf1_{(0,\Lambda)}.
\]
Summing over nonzero integral ideals of $K$, define the norm-line
theta map and its full Gram image projection by
\[
 E_\rho\psi(t)=e^{t/2}\sum_{\mathfrak a\ne0}
                     \chi(\mathfrak a)\psi(N\mathfrak a\,e^t),
 \qquad
 Q_{\rho,\Lambda}^{(\nu)}
     =P_{E_\rho\operatorname{Ran}\Pi_{\rho,\Lambda}^{(\nu)}}.
\]
There is no principal subtraction or constraint in this family.
For the nontrivial zeros $\omega=\beta_\omega+i\gamma_\omega$
of $L_\Q(s,\rho)$, with multiplicities $m_\omega$, put
\[
 \mathfrak b_\rho
   =\sum_\omega m_\omega\mu_{\beta_\omega-1/2,\gamma_\omega},
 \quad
 \mathfrak B_\rho(h)=\int H(i\tau)\,d\mathfrak b_\rho(\tau),
 \quad
 \mathfrak W_\rho(h)=\sum_\omega m_\omega H(\omega-1/2).
\]
Here $H(z)=\int h(t)e^{zt}dt$ and $\mu_{a,\gamma}$ is the
probability measure of Section~\ref{sec:harmonic}: its density is
$|a|/(\pi((\tau-\gamma)^2+a^2))$ when $a\ne0$, and it is
$\delta_\gamma$ when $a=0$.

\begin{corollary}[The cutoff for quadratic-induced Artin functions]
\label{cor:artin-induced-cutoff}
The representation $\rho$ has finite image, is irreducible of
dimension two, has Artin conductor $q_\rho$, and satisfies
\begin{equation}\label{eq:artin-induced-L}
                   L_\Q(s,\rho)=L_K(s,\chi)
\end{equation}
with every finite Euler factor, including the ramified factors.
This function is entire. Its real-place gamma factor is
$\Gamma_\R(s)\Gamma_\R(s+1)=\Gamma_\C(s)$, where
$\Gamma_\R(s)=\pi^{-s/2}\Gamma(s/2)$ and
$\Gamma_\C(s)=2(2\pi)^{-s}\Gamma(s)$.

Replacing $\chi$ by $\chi^\sigma$ leaves the raw selection,
theta map, and projection $Q_{\rho,\Lambda}^{(\nu)}$ unchanged.
For every fixed $\rho$ in this family, fixed $0<\nu<1$, and
$h\in C_c^\infty(\R)$,
\begin{equation}\label{eq:artin-induced-trace}
 \lim_{\Lambda\to\infty}
 \left(2\log\Lambda\,h(0)
      -\Tr(Q_{\rho,\Lambda}^{(\nu)}V(h))\right)
                         =\mathfrak B_\rho(h).
\end{equation}
The threshold $\nu=1/2$ is included, and
$\operatorname{rank}\Pi_{\rho,\Lambda}^{(\nu)}
\sim2c_\rho\Lambda/\pi$.

Put $L=\log\Lambda$, $S_L=\mathbf1_{[-L,L]}$, and let
$(z_{i,\Lambda})$ be an orthonormal basis of the theta image.
The nonnegative measures
\begin{equation}\label{eq:artin-induced-measures}
 d\mu_{\rho,\Lambda}(\tau)=
 \left(2L-\sum_i|\widehat{S_Lz_{i,\Lambda}}(\tau)|^2\right)
                                      \frac{d\tau}{2\pi}
\end{equation}
converge to $\mathfrak b_\rho$ on every Schwartz test.
For some $C,A$ depending on $\rho,\nu$, their masses on
$[T-1,T+1]$ are at most $C(1+|T|)^A$, uniformly for all
sufficiently large $\Lambda$ and all real $T$.
No zero-location, simplicity, or separation hypothesis is used.

If $L_\Q(s,\rho)$ is zero-free for $\Re s>\theta$,
$1/2\leq\theta<1$, put $\delta=\theta-1/2$. Then
\begin{equation}\label{eq:artin-induced-band}
 |\mathfrak B_\rho(h)-\mathfrak W_\rho(h)|
       \leq2\delta\mathcal N_\delta(h)Z_{2,\rho},
 \qquad
 Z_{2,\rho}=\sum_\omega\frac{m_\omega}{1+\gamma_\omega^2}
                  \ll_K\log(2+q_\rho),
\end{equation}
with $\mathcal N_\delta$ from Theorem~\ref{thm:band-bound}.
By \eqref{eq:artin-induced-trace}, this bounds the deviation of
the uncompressed trace limit from $\mathfrak W_\rho$.
\end{corollary}

\begin{proof}
The intersection $\ker\chi\cap j(\ker\chi)j^{-1}$ is open and
normal in $G_\Q$, so the induction factors through a finite group.
In the two coset lines its restriction to $G_K$ is
$\chi\oplus\chi^\sigma$, and $j$ exchanges these lines.
Since the characters differ, some element has distinct
eigenvalues on the two lines. Any invariant line must be one
of them, which contradicts the exchange. This proves
irreducibility and dimension two.

Artin induction gives \eqref{eq:artin-induced-L} prime by prime
\citep[Theorem~5.10 and Lemma~5.12, pp.~59--61]{horawaLfunctions}.
Explicitly, write $W_\rho$ for the representation space,
$I_p$ for inertia at a rational prime $p$, and
$\operatorname{Frob}_p$ for arithmetic Frobenius modulo inertia.
The local assertion is
\[
 \det(1-p^{-s}\rho(\operatorname{Frob}_p)\mid W_\rho^{I_p})^{-1}
   =\prod_{\mathfrak p\mid p}
                 (1-\chi(\mathfrak p)(N\mathfrak p)^{-s})^{-1}.
\]
Inertia invariants remove a factor exactly when the inducing
character is ramified there. An unramified character at the
prime above three still contributes its degree-one factor;
ramification of $K/\Q$ does not delete it.
The conductor induction formula
\citep[Proposition~5.58(2), p.~78]{horawaLfunctions} gives
$\operatorname{cond}(\rho)=|D_K|N\mathfrak f_\chi=q_\rho$,
where $D_K=-3$ is the field discriminant.
The nonprincipal Hecke function is entire.

The two eigenvalues of $j$ are $+1,-1$.
Gamma duplication consequently gives the standard completion
\[
 q_\rho^{s/2}\Gamma_\R(s)\Gamma_\R(s+1)L_\Q(s,\rho)
                =2c_\rho^{-s}\Gamma(s)L_K(s,\chi).
\]
The constant two does not change the functional equation or
zero divisor, and $c_\rho$ is exactly the Hankel parameter.

On ideals $\chi^\sigma(\mathfrak a)=\chi(\sigma\mathfrak a)$.
Conjugation preserves norms and conjugates conductors, so
reindexing the theta sum gives $E_{\chi^\sigma}^K=E_\chi^K$.
Their raw selections and full Gram projections are identical.
Thus one representative of the orbit $\{\chi,\chi^\sigma\}$
represents the Artin factor; summing both sectors counts it twice.

The normalized idele norms satisfy
$|N_{K/\Q}x|_\Q=|x|_K$. The inherited variable is therefore the
absolute norm $r$, with $t=\log r$ and Haar measure $dr/r$.
There is no rescaling of $t$ by two: the single norm channel
has leading term $2\log\Lambda$, not $4\log\Lambda$.

The zeros in \eqref{eq:artin-induced-L} coincide with full
multiplicity, and the defined projection is precisely the
Hecke projection of Theorem~\ref{nh:theorem}. That theorem
proves the trace and measure assertions and the polynomial
mass bound. Bessel's inequality also gives the pointwise
positivity in \eqref{eq:artin-induced-measures}.
There are no principal polar coordinates because $\chi\ne1$.
The unitary functional equation and coefficient conjugation
reflect a zero $\omega$ to $1-\overline\omega$; this does not
require $\chi^\sigma=\overline\chi$.
The half-plane thus confines the divisor to width $\delta$.
Theorem~\ref{thm:band-bound} and Corollary~\ref{cor:hecke-band}
give \eqref{eq:artin-induced-band}.
\end{proof}

This is a degree-two Galois representation, not a nonabelian
character of the abelian group $C_\Q$. Its cutoff is inherited
from the $K$-Hecke norm channel; no $\rho$-isotypic component
of a $C_\Q$ representation is constructed.
The cubic Kummer example above is an instance: its $S_3$
permutation representation is $1\oplus\rho$, so its irreducible
Artin factor has degree two despite the cubic field.
This direct norm-variable identification is a subcase of
Corollary~\ref{hnew:odd-artin}, which also covers the odd
two-dimensional representations not induced from this field.

Even two-dimensional representations require a different
archimedean transform. For an arbitrary irreducible degree-three representation,
an index-two character induction cannot apply. Even a monomial
induction from a cubic field requires a cutoff theorem for that
different field and its archimedean factors.
In general a zeros-only trace also requires exclusion of poles
in the relevant strip; Brauer induction has possibly negative
integer coefficients and yields a quotient of Hecke functions.
Even given holomorphy and an automorphic theta realization,
the Gram control, normalized zero-jet tails, endpoint
closures, graph approximation, local limit, polynomial frequency
mass, and trace transfer still have to be proved for that cutoff.
If the windowed construction of Section~\ref{glnw:section} is used,
one must additionally prove the normalization comparison
$\Tr((\widetilde Q-S_LQS_L)(V(h)\otimes I_d))\to0$.
Exact finite-section positivity alone supplies none of these
limits for those remaining representation classes.


\subsection{Moments and inertia}
Suzuki's operator records each distinct ordinate as one spectral
coordinate \citep[Theorems~1.1--1.2]{suzuki2026}. Its ordinary trace
must therefore be distinguished from the multiplicity-weighted
logarithmic derivative of the completed zeta function. The next
proposition states the resulting comparison as an abstract operator
fact, with all its hypotheses explicit.

\begin{proposition}[Moments and the integer deficit]
\label{prop:moment-inertia}
Let $0<\gamma_1<\gamma_2<\cdots$ tend to infinity, let $m_j\geq1$
be integers, and assume $\sum_jm_j\gamma_j^{-2}<\infty$.
Let $C$ be a self-adjoint Hilbert--Schmidt operator with simple
nonzero eigenvalues $\pm1/\gamma_j$. Suppose a unitary involution $J$
satisfies $JCJ=-C$, and finite-rank orthogonal projections $P_N$
increase strongly to $I$ and commute with $J$. Set
\[
 C_N=P_NCP_N,\qquad
 S_k=2\sum_jm_j\gamma_j^{-2k},\qquad
 D_{N,k}=S_k-\Tr C_N^{2k},\quad k\geq1.
\]
Then $D_{N,k}\geq0$, it decreases with $N$, and
\begin{equation}\label{eq:moment-excess-limit}
 \lim_ND_{N,k}=2\sum_j(m_j-1)\gamma_j^{-2k}.
\end{equation}
Define
\[
 n(T)=\sum_{\gamma_j\leq T}m_j,
 \qquad \iota_N(T)=\#\{\lambda>1/T:
                  \lambda\text{ is an eigenvalue of }C_N\}.
\]
Eigenvalues in $\iota_N(T)$ are counted with their multiplicities.
At a height $T$ different from every $\gamma_j$,
$\iota_N(T)$ increases to $\#\{j:\gamma_j<T\}$. Moreover
\begin{equation}\label{eq:moment-inertia-integral}
 D_{N,k}=4k\int_0^\infty
                [n(t)-\iota_N(t)]t^{-2k-1}\,dt.
\end{equation}
If $\gamma_-$ is the largest ordinate below such a $T$, then
\begin{equation}\label{eq:moment-gap-deficit}
 n(T)-\iota_N(T)
    \leq\frac{D_{N,k}}
                  {2(\gamma_-^{-2k}-T^{-2k})}.
\end{equation}
With the inner product conjugate-linear in its first argument,
any basis $(v_j)$ of $P_N\mathcal H$ defines
$G_N=(\langle v_i,v_j\rangle)_{i,j}$ and
$H_N=(\langle v_i,Cv_j\rangle)_{i,j}$. Then
$\iota_N(T)=\operatorname{ind}_+(TH_N-G_N)$.
\end{proposition}

\begin{proof}
Write $\lambda_j=1/\gamma_j$ in decreasing order and list the
positive eigenvalues $\mu_{N,j}$ of $C_N$ in decreasing order,
appending zeros. Min--max gives $0\leq\mu_{N,j}\leq\lambda_j$
and monotonicity under nested compression. The parity involution
pairs the positive and negative eigenvalues, whence
\begin{equation}\label{eq:moment-defect-decomposition}
 D_{N,k}=2\sum_j(m_j-1)\lambda_j^{2k}
              +2\sum_j(\lambda_j^{2k}-\mu_{N,j}^{2k}).
\end{equation}
This proves nonnegativity and monotonicity. Approximation of $C$
by a finite-rank operator and strong convergence of $P_N$ show
$C_N\to C$ in Hilbert--Schmidt norm. In the telescoping identity
for $C_N^{2k}-C^{2k}$, place the difference and one other factor
in Hilbert--Schmidt norm and the remaining factors in operator
norm. The inequality $\|AB\|_1\leq\|A\|_2\|B\|_2$ proves
trace-norm convergence and \eqref{eq:moment-excess-limit}.

Only finitely many $\lambda_j$ exceed $1/T$. Since none equals
that threshold, operator-norm convergence puts all of them above
it for large $N$; min--max excludes extra ones. This proves the
count limit. Apply
$\lambda^{2k}=2k\int_{1/\lambda}^\infty t^{-2k-1}dt$
to the full and compressed eigenvalue lists. Tonelli's theorem
applies separately to their nonnegative convergent sums, proving
\eqref{eq:moment-inertia-integral}; the values at isolated
thresholds do not change the integral.

For \eqref{eq:moment-gap-deficit}, put
$r=1/T$, $n_0=\#\{j:\gamma_j<T\}$, and $i=\iota_N(T)$.
For $i<j\leq n_0$ one has $\lambda_j\geq1/\gamma_-$ and
$\mu_{N,j}\leq r$. Each such term in the second sum of
\eqref{eq:moment-defect-decomposition} contributes at least
$2\Delta$, where $\Delta=\gamma_-^{-2k}-T^{-2k}>0$.
Each multiplicity excess below $T$ contributes at least
$2\gamma_-^{-2k}\geq2\Delta$ in the first sum. Discarding the
other nonnegative terms gives
$D_{N,k}\geq2\Delta[n(T)-i]$. Finally, congruence by
$G_N^{-1/2}$ transforms $TH_N-G_N$ into the matrix of
$TC_N-I$ in an orthonormal basis. Sylvester's law of inertia
proves the matrix assertion.
\end{proof}

Put $\xi(s)=\tfrac12s(s-1)\pi^{-s/2}\Gamma(s/2)\zeta(s)$ and
$\Xi(z)=\xi(1/2+iz)$. Under RH, its paired Hadamard product gives
$S_k=-2k[z^{2k}]\log(\Xi(z)/\Xi(0))$ with these arithmetic
multiplicities. Formula \eqref{eq:moment-excess-limit} can be
positive even when every zero is on the line, whereas
$\mathfrak B-\mathfrak W$ then vanishes. The defects are therefore
different. Even after subtracting the multiplicity term, a
compressed moment is not generally a linear discarded trace.
Writing $C$ in blocks $\bigl(\begin{smallmatrix}A&B^*\\B&D\end{smallmatrix}\bigr)$
relative to $P$ and $I-P$ gives
\[
 \Tr C^2-\Tr(PCP)^2=2\|B\|_2^2+\|D\|_2^2,
 \qquad
 \Tr((I-P)C^2)=\|B\|_2^2+\|D\|_2^2.
\]
Nor can the gap in \eqref{eq:moment-gap-deficit} simply be omitted.
An operator with eigenvalues $\pm(1+\varepsilon),0,0$ has a
two-dimensional compression with eigenvalues $\pm1$, by choosing
the squared cosine of each compression angle to be
$(1+\varepsilon)^{-1}$. At threshold one the count deficit is one
but the even moment loss is $2((1+\varepsilon)^{2k}-1)\to0$.

\section{A subordinate comparison with divisor shells}
\label{sec:divisor-shells}
The following two Euler products explain what a smoothed shell
count detects, while keeping its all-integer quantifier separate
from a prime-by-prime occupancy assertion.

Let $\eta$ be the nontrivial character modulo three. For $a$ prime
to three, take two residue labels for each positive divisor
$d\mid a$, namely $d$ and $4ad$ modulo three. Let $W(a)$ count
ordered pairs of opposite labels. If $p=4a-3>3$ is prime, these
are the residues of the divisors of $pa$, and every opposite pair
$d,e\mid pa$ gives positive integer denominators
\[
 b=\frac{pa(d+e)}{3e},\qquad
 c=\frac{pa(d+e)}{3d},\qquad
 \frac4p=\frac1a+\frac1b+\frac1c.
\]
Conversely, writing $b=ku,c=kv$ with $(u,v)=1$ in a solution
gives $3kuv=pa(u+v)$; thus $u,v\mid pa$ and $3\mid u+v$.
For composite $4a-3$ the residue-tag definition still specifies
$W(a)$, without asserting that it is a prime shell count.

\begin{theorem}[The two rational counts]
\label{thm:two-shell-counts}
Put $S=\{2,3,5,7\}$. Let $\chi$ be trivial or the cubic character
modulo seven with $\chi(3)=e^{2\pi i/3}$. For $h=1,25$ and
$w\in C_c^\infty((0,\infty))$ define
\[
 H_h(X)=\sum_{\substack{m\equiv1\ (6)\\(m,210)=1}}
               \chi(m)W(hm)w(m/X).
\]
Write $L^{(S)}$ for deletion of Euler factors at $S$, and set
\[
 P(s)=L^{(S)}(s,\chi),\quad Q(s)=L^{(S)}(s,\chi\eta),
 \quad Z(s)=L^{(S)}(2s,\chi^2).
\]
The two Dirichlet series are
\begin{equation}\label{eq:shell-numerators}
 D_1(s)=\frac{(P(s)^2-Q(s)^2)^2}{Z(s)},\qquad
 D_{25}(s)=\frac{9(P(s)^4+Q(s)^4)-2P(s)^2Q(s)^2}{Z(s)}.
\end{equation}
Let $J_h(X)=\operatorname*{Res}_{s=1}
X^s\widehat w(s)D_h(s)$, with
$\widehat w(s)=\int_0^\infty w(t)t^{s-1}dt$.
The critical-line assertion for $L(s,\chi^2)$ is equivalent to
\begin{equation}\label{eq:shell-quarter-power}
 H_h(X)=J_h(X)+O_{w,\varepsilon}(X^{1/4+\varepsilon})
 \quad(h=1,25)
\end{equation}
for every such $w$ and every $\varepsilon>0$.
Hence the verified Dirichlet theorem, when used as the external
input stated in the introduction, gives both remainder bounds.
For trivial $\chi$, $J_h$ is $X$ times a polynomial of degree at
most three in $\log X$; for cubic $\chi$ it is zero.
\end{theorem}

\begin{proof}
Put $\tau(m)=\sum_{d\mid m}1$ and
$b(m)=\sum_{d\mid m}\eta(d)$. When $m\equiv1\pmod3$,
the two labels attached to a divisor have the same sign.
Counting opposite ordered signs gives
$W(m)=2(\tau(m)^2-b(m)^2)$.
Since $(m,5)=1$ and $1+\eta(5)+\eta(5)^2=1$,
\[
 W(25m)=18\tau(m)^2-2b(m)^2.
\]
Divisor complementation gives $b(m)=\eta(m)b(m)$, so its square
already vanishes on $\eta(m)=-1$. At a prime outside $S$, with
$T=\chi(p)p^{-s}$, summing the two elementary geometric series
and their derivatives gives
\[
 \sum_{k\geq0}(k+1)^2T^k=\frac{1-T^2}{(1-T)^4},
 \qquad
 \sum_{k\geq0}\left(\sum_{j=0}^k\eta(p)^j\right)^2T^k
     =\frac{1-T^2}{(1-T)^2(1-\eta(p)T)^2}.
\]
The residue projector $(1+\eta(m))/2$ turns the first global
series into $(P^4+Q^4)/(2Z)$, while the second is $P^2Q^2/Z$.
Substitution proves \eqref{eq:shell-numerators}.

Assume the critical-line assertion. Then $1/Z$ is holomorphic
for $\Re s>1/4$, with at most a zero where $Z$ has a pole.
It has polynomial growth on every fixed closed strip strictly
inside this region. Apply the following argument to $L(z,\chi^2)$
in the variable $z=2s$, whose zero-free region is $\Re z>1/2$.
Fix a height $t_0$ and a chain
of disks, of bounded radii and number, joining the point in
question to an absolutely convergent Euler-product line.
Polynomial strip growth makes
$C\log(3+|t_0|)-\log|L(z,\chi^2)|$ positive and harmonic on
these disks for a sufficiently large fixed $C$. Harnack's
inequality, anchored on the Euler-product line, bounds it by
$O(\log(3+|t_0|))$. For the principal character first replace
$L(z,\chi^2)$ by $(z-1)L(z,\chi^2)/(z+1)$ to remove its pole.
The fixed distance to the zero-free boundary controls the chain
uniformly in $t_0$. The finite deleted factors
are nonzero in $\Re s>0$ and have bounded inverses there away
from that boundary. The numerator also has polynomial strip
growth, with a possible pole only at one. Mellin inversion,
rapid vertical decay of $\widehat w$, and shifting the contour
to $\Re s=1/4+\varepsilon$ prove
\eqref{eq:shell-quarter-power} and its stated main term.

Conversely, take the Mellin transform in $X$ with measure
$X^{-s-1}dX$. The sum vanishes for sufficiently small $X$, and
the asserted remainder continues $\widehat w(s)D_h(s)$
holomorphically to $\Re s>1/4$, except for its possible pole at
one. At any specified point choose a nonnegative bump sufficiently
near one that $\widehat w$ is nonzero there. Thus both $D_h$ are
holomorphic throughout that half-plane away from one.

If $L(\rho,\chi^2)=0$ with $1/2<\Re\rho<1$, both numerators
in \eqref{eq:shell-numerators} vanish at $s=\rho/2$. The first
implies $P^2=Q^2$; the second then equals $16P^4$. Hence
$L(\rho/2,\chi)=0$, since deleted Euler factors have no zeros
there. The functional equation produces a zero of
$L(s,\overline\chi)=L(s,\chi^2)$ at $1-\rho/2$.
The map $T(s)=1-s/2$ preserves the open right half of the strip
and satisfies
$T^n(s)=2/3+(-1/2)^n(s-2/3)$.
Any starting point other than $2/3$ would give infinitely many
distinct zeros accumulating at an interior point, which is
impossible for this nonzero meromorphic function.

The fixed point is nonzero. For trivial $\chi$, the alternating
Dirichlet series gives $\zeta(u)<0$ for $0<u<1$.
For the chosen cubic character and real $u>0$, its imaginary
part is $\sqrt3/2$ times
\[
 \sum_{k\geq0}\bigl((7k+3)^{-u}+(7k+4)^{-u}
                  -(7k+2)^{-u}-(7k+5)^{-u}\bigr)<0.
\]
Each term is strictly negative by convexity of $x^{-u}$, and
periodic cancellation justifies the grouped Dirichlet series.
The conjugate character is likewise nonzero at $2/3$.
Functional-equation reflection and coefficient conjugation
exclude left-half zeros as well. This proves the equivalence.
\end{proof}

Imposing primality of $4hm-3$ replaces these Euler products by
shifted prime--divisor correlations. It needs an additional
arithmetic estimate. The marked Type~II inputs of
\citet[Lemma~10.2, p.~64; Proposition~10.3, p.~67;
Lemma~12.1, p.~75]{openai2026primitive} supply averaged
replacement estimates. Varying-cofactor arguments use bounded
mark weights and a separate bound for failures to obtain positive
lower proportions of primes in the primitive-root, primitive-point,
and cyclic-reduction problems
\citep[Theorems~1--2]{bhargava2026roots,bhargava2026points,bhargava2026cyclic}.
These quantifiers differ from occupation of every shell in a
growing interval for almost every prime, and from simultaneous
failure of an entire menu. A signed rough-sieve identity that
bounds a prime sum plus a composite sum does not bound either
sum separately. At shift three, for example, the four primes
$1009,2689,3529,32089$ have quartic signs respectively
$-, -, +, +$, while their shells are respectively occupied,
empty, empty, occupied; the sign alone does not determine
occupancy.

A sufficient uniform upgrade can be stated precisely. For each
remaining prime $p\in[X,2X]$, let $M(p)$ be the nonnegative
integer number of primitive marks in an allowed menu, each
mapped to a solution. It suffices to prove, uniformly outside
an explicit exceptional set $\mathcal E_X$,
\[
 M(p)=A(p)+E(p),\qquad A(p)\geq1+\eta_X,
 \qquad |E(p)|\leq\eta_X,
\]
and either $|\mathcal E_X|<1$ or a separate solution for every
prime in $\mathcal E_X$. Applying this on dyadic intervals and
checking the finite initial range covers all primes. A bound on
mark multiplicity is needed to convert averaged mark mass into
counts of distinct primes or triples; an explicit map from each
positive mark already suffices for existence. No mean or
positive-proportion theorem above supplies this uniform upgrade.
